\section{Curves over a field}
\label{pa-chap:Curves}

The scheme-theoretic curve facts used here follow the standard treatments in Hartshorne
\cite{hartshorne-ag} and the Stacks Project \cite{stacks-project}.

Throughout this chapter \(k\) is a field. As in Chapter~\ref{pa-chap:Challenge}, a \emph{smooth
curve} over \(k\) is a \(k\)-scheme which is smooth of relative dimension one, proper, and
geometrically irreducible over \(k\). This chapter provides the geometric substrate for the
genus: smooth schemes over a field are reduced, so a smooth curve is geometrically integral;
a proper geometrically integral scheme over a field has no global functions besides the
constants; the projective line carries two affine charts with polynomial section rings and
Laurent-polynomial overlap; the section rings of the standard smooth charts of a curve are
Dedekind domains, so the local rings of the curve are valuation rings and its
specialization order has height one; a rational function transcendental over \(k\) therefore
spreads out to a locally quasi-finite morphism \(C \to \PP^1\); and Zariski's main theorem
upgrades it to a finite one.

\section{Smooth algebras over a domain are reduced}

\begin{lemma}[Unramified flat algebras over a domain are reduced]
  \label{pa-lem:unramified_flat_reduced}
  \lean{Algebra.FormallyUnramified.isReduced_of_isDomain}
  \leanok
  \dcref{custom}

  Let \(B\) be an integral domain and let \(S\) be a commutative \(B\)-algebra which is
  formally unramified, essentially of finite type, and flat over \(B\). Then \(S\) is reduced.
\end{lemma}
\begin{proof}
  \leanok
  Let \(K = \Frac B\) be the fraction field. The \(K\)-algebra \(K \otimes_B S\) is again
  formally unramified and essentially of finite type, both properties being stable under base
  change, and a formally unramified, essentially of finite type algebra over a field is
  reduced. It therefore suffices to embed \(S\) into \(K \otimes_B S\). The canonical map
  \(s \mapsto 1 \otimes s\) is, up to the isomorphism \(S \cong B \otimes_B S\), obtained from
  the injection \(B \hookrightarrow K\) by tensoring with \(S\); since \(S\) is flat over
  \(B\), it is injective. A subring of a reduced ring is reduced.
\end{proof}

\begin{lemma}[\'Etale algebras over a domain are reduced]
  \label{pa-lem:etale_reduced_domain}
  \lean{Algebra.Etale.isReduced_of_isDomain}
  \leanok
  \dcref{custom}

  An \'etale algebra over an integral domain is reduced.
\end{lemma}
\begin{proof}
  \leanok
  \uses{pa-lem:unramified_flat_reduced}
  An \'etale algebra is formally unramified and of finite presentation, in particular
  essentially of finite type, and it is flat. Apply \ref{pa-lem:unramified_flat_reduced}.
\end{proof}

\begin{lemma}[Standard smooth homomorphisms out of a domain]
  \label{pa-lem:standard_smooth_reduced_ringhom}
  \lean{RingHom.IsStandardSmooth.isReduced_of_isDomain}
  \leanok
  \dcref{custom}

  Let \(B\) be an integral domain and let \(B \to S\) be a standard smooth ring homomorphism,
  that is, one admitting a presentation \(S = B[x_1,\dots,x_n]/(f_1,\dots,f_c)\) in which the
  Jacobian matrix \((\partial f_u/\partial x_v)_{u,v \le c}\) has invertible determinant in
  \(S\). Then \(S\) is reduced.
\end{lemma}
\begin{proof}
  \leanok
  \uses{pa-lem:etale_reduced_domain}
  By the structure theorem for standard smooth algebras, \(S\) is \'etale over a polynomial
  ring over the base: there are \(m\) and a factorization
  \(B \to B[x_1,\dots,x_m] \to S\) with the second map \'etale. A polynomial ring over an
  integral domain is an integral domain, so \ref{pa-lem:etale_reduced_domain} applies.
\end{proof}

\begin{lemma}[Standard smooth algebras over a domain]
  \label{pa-lem:standard_smooth_reduced}
  \lean{Algebra.IsStandardSmooth.isReduced_of_isDomain}
  \leanok
  \dcref{custom}

  A standard smooth algebra over an integral domain is reduced.
\end{lemma}
\begin{proof}
  \leanok
  \uses{pa-lem:standard_smooth_reduced_ringhom}
  The structure homomorphism is a standard smooth ring homomorphism out of a domain; apply
  \ref{pa-lem:standard_smooth_reduced_ringhom}.
\end{proof}

\begin{theorem}[Smooth schemes over a field are reduced]
  \label{pa-thm:smooth_reduced_field}
  \lean{AlgebraicGeometry.Smooth.isReduced_of_field}
  \leanok
  \dcref{custom}

  Let \(K\) be a field and let \(X\) be a scheme which is smooth over \(\Spec K\). Then \(X\)
  is reduced.
\end{theorem}
\begin{proof}
  \leanok
  \uses{pa-lem:standard_smooth_reduced_ringhom}
  A smooth morphism is, locally on the source and the target, standard smooth: every point
  \(x \in X\) has an affine open neighbourhood \(V\), mapping into an affine open \(U\) of the
  base, such that the induced ring homomorphism \(\Gamma(U) \to \Gamma(V)\) is standard
  smooth. Since \(\Spec K\) has a single point and \(U\) is a nonempty open subset containing
  the image of \(x\), we have \(U = \Spec K\), and \(\Gamma(U) \cong K\) is a field, in
  particular a domain. By \ref{pa-lem:standard_smooth_reduced_ringhom} the ring \(\Gamma(V)\) is
  reduced, hence the affine scheme \(V\) is reduced. The opens \(V\) cover \(X\), and a scheme
  covered by reduced open subschemes is reduced.
\end{proof}

\begin{theorem}[Smooth morphisms are geometrically reduced]
  \label{pa-thm:smooth_geometrically_reduced}
  \lean{AlgebraicGeometry.Smooth.geometricallyReduced}
  \leanok
  \dcref{custom}

  A smooth morphism of schemes \(f : X \to Y\) is geometrically reduced: for every field \(K\)
  and every morphism \(\Spec K \to Y\), the base change \(X \times_Y \Spec K\) is a reduced
  scheme.
\end{theorem}
\begin{proof}
  \leanok
  \uses{pa-thm:smooth_reduced_field}
  Smoothness is stable under base change, so the projection
  \(X \times_Y \Spec K \to \Spec K\) is smooth; apply \ref{pa-thm:smooth_reduced_field}.
\end{proof}

\begin{lemma}[Smooth of relative dimension one implies smooth]
  \label{pa-lem:smooth_of_relDim_one}
  \lean{AlgebraicGeometry.Smooth.of_smoothOfRelativeDimension_one}
  \leanok
  \dcref{custom}

  A morphism of schemes which is smooth of relative dimension one is smooth.
\end{lemma}
\begin{proof}
  \leanok
  A morphism smooth of some relative dimension is smooth.
\end{proof}

\begin{lemma}[Smooth of a relative dimension implies geometrically reduced]
  \label{pa-lem:relDim_geometrically_reduced}
  \lean{AlgebraicGeometry.SmoothOfRelativeDimension.geometricallyReduced}
  \leanok
  \dcref{custom}

  For every \(n\), a morphism of schemes which is smooth of relative dimension \(n\) is
  geometrically reduced.
\end{lemma}
\begin{proof}
  \leanok
  \uses{pa-thm:smooth_geometrically_reduced}
  Such a morphism is smooth; apply \ref{pa-thm:smooth_geometrically_reduced}.
\end{proof}

\begin{theorem}[Smooth and geometrically irreducible implies geometrically integral]
  \label{pa-thm:smooth_geometrically_integral}
  \lean{AlgebraicGeometry.Smooth.geometricallyIntegral}
  \leanok
  \dcref{custom}

  A smooth, geometrically irreducible morphism of schemes is geometrically integral: every
  base change to the spectrum of a field is an integral scheme.
\end{theorem}
\begin{proof}
  \leanok
  \uses{pa-thm:smooth_geometrically_reduced}
  Every base change to the spectrum of a field is reduced
  (\ref{pa-thm:smooth_geometrically_reduced}) and irreducible (geometric irreducibility), and a
  reduced scheme with irreducible --- in particular nonempty --- underlying space is integral.
\end{proof}

\begin{theorem}[Relative-dimension form]
  \label{pa-thm:relDim_geometrically_integral}
  \lean{AlgebraicGeometry.SmoothOfRelativeDimension.geometricallyIntegral}
  \leanok
  \dcref{custom}

  For every \(n\), a geometrically irreducible morphism of schemes which is smooth of relative
  dimension \(n\) is geometrically integral.
\end{theorem}
\begin{proof}
  \leanok
  \uses{pa-thm:smooth_geometrically_integral}
  Such a morphism is smooth; apply \ref{pa-thm:smooth_geometrically_integral}.
\end{proof}

\section{Schemes over a field: the topological substrate}

In this section \(X\) is a scheme over \(k\), that is, a scheme equipped with a structure
morphism \(X \to \Spec k\).

\begin{lemma}[Irreducibility]
  \label{pa-lem:irreducible_of_geomirred}
  \lean{AlgebraicGeometry.irreducibleSpace_left_of_geometricallyIrreducible}
  \leanok
  \dcref{custom}

  If \(X\) is geometrically irreducible over \(k\), then the underlying topological space of
  \(X\) is irreducible; in particular \(X\) is nonempty.
\end{lemma}
\begin{proof}
  \leanok
  Geometric irreducibility applied to the trivial extension \(k \to k\) says that
  \(X \times_{\Spec k} \Spec k \cong X\) is irreducible.
\end{proof}

\begin{lemma}[Quasi-compactness]
  \label{pa-lem:curve_compact}
  \lean{AlgebraicGeometry.compactSpace_left_of_quasiCompact}
  \leanok
  \dcref{custom}

  If \(X\) is quasi-compact over \(\Spec k\) (for instance proper), then the underlying
  topological space of \(X\) is quasi-compact.
\end{lemma}
\begin{proof}
  \leanok
  \(\Spec k\) is a one-point space, hence quasi-compact, and \(X\) is the preimage of it
  under a quasi-compact morphism.
\end{proof}

\begin{lemma}[Quasi-separatedness]
  \label{pa-lem:curve_quasiseparated}
  \lean{AlgebraicGeometry.quasiSeparatedSpace_left_of_quasiSeparated}
  \leanok
  \dcref{custom}

  If \(X\) is quasi-separated over \(\Spec k\) (for instance proper), then the underlying
  topological space of \(X\) is quasi-separated: the intersection of any two quasi-compact
  opens is quasi-compact.
\end{lemma}
\begin{proof}
  \leanok
  \(\Spec k\) is affine, hence a quasi-separated scheme, and a scheme quasi-separated over a
  quasi-separated scheme is quasi-separated.
\end{proof}

\begin{proposition}[Integrality]
  \label{pa-prop:curve_integral}
  \lean{AlgebraicGeometry.isIntegral_left_of_geometricallyReduced}
  \leanok
  \dcref{custom}

  If \(X\) is geometrically reduced and geometrically irreducible over \(k\), then \(X\) is an
  integral scheme.
\end{proposition}
\begin{proof}
  \leanok
  \uses{pa-lem:irreducible_of_geomirred}
  Applying geometric reducedness to the trivial extension \(k \to k\) shows that \(X\) is
  reduced, and \(X\) is irreducible by \ref{pa-lem:irreducible_of_geomirred}. A reduced scheme
  with irreducible underlying space is integral.
\end{proof}

\section{Global sections of a proper geometrically integral scheme}

\begin{lemma}[Pushout of a field extension along itself]
  \label{pa-lem:pushout_field_epi}
  \lean{CommRingCat.epi_of_isPushout_of_isField}
  \leanok
  \dcref{custom}

  Let \(\varphi : A \to B\) and \(\psi : A \to C\) be homomorphisms of commutative rings with
  \(B\) nontrivial, and let \(D\), with coprojections \(\iota_B : B \to D\) and
  \(\iota_C : C \to D\), be a pushout of \((\varphi, \psi)\) (for instance
  \(D = B \otimes_A C\)). Assume that \(D\) is a field and that there is a ring isomorphism
  \(e : C \cong B\) with \(e \circ \psi = \varphi\). Then \(\varphi\) is an epimorphism in the
  category of commutative rings.
\end{lemma}
\begin{proof}
  \leanok
  Let \(\mu : D \to B\) be the homomorphism with \(\mu \circ \iota_B = \id_B\) and
  \(\mu \circ \iota_C = e\), given by the universal property of the pushout (the two maps
  agree on \(A\) because \(e \circ \psi = \varphi\)). Since \(D\) is a field and \(B\) is
  nontrivial, \(\mu\) is injective; it is surjective because \(\mu \circ \iota_B = \id_B\).
  Hence \(\mu\) is an isomorphism, and from
  \(\mu \circ \iota_B = \id_B = \mu \circ (\iota_C \circ e^{-1})\) we conclude
  \(\iota_B = \iota_C \circ e^{-1}\).

  Now let \(u, v : B \to W\) satisfy \(u \circ \varphi = v \circ \varphi\). Then
  \(u \circ \varphi = (v \circ e) \circ \psi\), so there is
  \(w : D \to W\) with \(w \circ \iota_B = u\) and \(w \circ \iota_C = v \circ e\). Therefore
  \[
    u = w \circ \iota_B = w \circ \iota_C \circ e^{-1} = v \circ e \circ e^{-1} = v,
  \]
  which proves that \(\varphi\) is an epimorphism.
\end{proof}

\begin{lemma}[Global sections form a field]
  \label{pa-lem:sections_field}
  \lean{AlgebraicGeometry.isField_of_universallyClosed_of_isField}
  \leanok
  \dcref{custom}

  Let \(R\) be a commutative ring which is a field, let \(X\) be an integral scheme, and let
  \(f : X \to \Spec R\) be universally closed. Then \(\Gamma(X, \struct X)\) is a field.
\end{lemma}
\begin{proof}
  \leanok
  \(\Gamma(X, \struct X)\) is an integral domain because \(X\) is integral. Since \(f\) is
  universally closed, the induced ring homomorphism
  \(R \cong \Gamma(\Spec R, \struct{\Spec R}) \to \Gamma(X, \struct X)\) is an integral
  extension. A domain which is integral over a field is a field.
\end{proof}

\begin{theorem}[Constant functions on a proper geometrically integral scheme]
  \label{pa-thm:sections_k}
  \lean{AlgebraicGeometry.bijective_appTop_of_isProper_of_geometricallyIntegral}
  \leanok
  \dcref{custom}

  Let \(K\) be a field and let \(X\) be a scheme, proper and geometrically integral over
  \(K\). Then the ring homomorphism
  \(\Gamma(\Spec K, \struct{\Spec K}) \to \Gamma(X, \struct X)\) induced by the structure
  morphism is bijective; identifying \(\Gamma(\Spec K, \struct{\Spec K}) = K\), this says
  \(\Gamma(X, \struct X) = K\).
\end{theorem}
\begin{proof}
  \leanok
  \uses{pa-lem:pushout_field_epi, pa-lem:sections_field, pa-lem:curve_compact, pa-lem:curve_quasiseparated}
  Write \(A = \Gamma(X, \struct X)\) and let \(\varphi : K \to A\) be the map in question.
  Applying geometric integrality to the trivial extension \(K \to K\) shows that \(X\) is
  integral; in particular \(X\) is nonempty and \(A\) is a nontrivial domain, so \(\varphi\),
  a ring homomorphism out of a field into a nontrivial ring, is injective. By
  \ref{pa-lem:sections_field} (properness includes universal closedness), \(A\) is a field, and
  since \(f\) is proper, \(A\) is moreover a finite \(K\)-module. As a finite ring
  homomorphism is surjective as soon as it is an epimorphism of commutative rings, it remains
  to prove that \(\varphi\) is an epimorphism.

  Let \(q : \Spec A \to \Spec K\) be the morphism induced by \(\varphi\) and form the base
  change \(Y = X \times_{\Spec K} \Spec A\). Since \(A\) is a field, geometric integrality of
  \(f\) makes \(Y\) integral, and the projection \(Y \to \Spec A\) is universally closed as a
  base change of \(f\); hence \(D := \Gamma(Y, \struct Y)\) is a field by
  \ref{pa-lem:sections_field}. On the other hand, \(X\) is quasi-compact and quasi-separated as
  a topological space (\ref{pa-lem:curve_compact}, \ref{pa-lem:curve_quasiseparated}), and \(q\) is
  flat because every vector space over \(K\) is flat; flat base change for global sections
  therefore exhibits the square of section rings
  \[
    \begin{array}{ccc}
      K & \longrightarrow & A \\
      \downarrow & & \downarrow \\
      A & \longrightarrow & D
    \end{array}
  \]
  (left column \(\varphi\), top row \(\varphi\) transported through the canonical isomorphism
  \(\Gamma(\Spec A, \struct{\Spec A}) \cong A\)) as a pushout square of commutative rings,
  \(D \cong A \otimes_K A\). Its pushout \(D\) is a field and its two legs are isomorphic
  copies of \(\varphi\), so \(\varphi\) is an epimorphism by \ref{pa-lem:pushout_field_epi}.
\end{proof}

\begin{corollary}[Isomorphism form]
  \label{pa-cor:sections_iso}
  \lean{AlgebraicGeometry.isIso_appTop_of_isProper_of_geometricallyIntegral}
  \leanok
  \dcref{custom}

  In the situation of \ref{pa-thm:sections_k}, the induced map
  \(\Gamma(\Spec K, \struct{\Spec K}) \to \Gamma(X, \struct X)\) is an isomorphism of rings.
\end{corollary}
\begin{proof}
  \leanok
  \uses{pa-thm:sections_k}
  A bijective ring homomorphism is an isomorphism.
\end{proof}

\begin{corollary}[Global sections of the curve]
  \label{pa-cor:curve_sections}
  \lean{AlgebraicGeometry.isIso_hom_appTop_of_geometricallyReduced}
  \leanok
  \dcref{custom}

  Let \(C\) be a scheme over \(k\) which is proper, geometrically irreducible and
  geometrically reduced over \(k\). Then the structure morphism induces an isomorphism
  \(\Gamma(\Spec k, \struct{\Spec k}) \cong \Gamma(C, \struct C)\); that is,
  \(\Gamma(C, \struct C) = k\).
\end{corollary}
\begin{proof}
  \leanok
  \uses{pa-cor:sections_iso}
  Every base change of \(C\) to the spectrum of a field is reduced and irreducible, hence
  integral; so \(C\) is geometrically integral over \(k\) and \ref{pa-cor:sections_iso} applies.
\end{proof}

\begin{remark}[The smooth curve]
  \label{pa-rem:curve_bundle}
  \dcref{custom}
  \uses{pa-lem:smooth_of_relDim_one, pa-lem:relDim_geometrically_reduced,
    pa-thm:relDim_geometrically_integral, pa-lem:irreducible_of_geomirred, pa-lem:curve_compact,
    pa-lem:curve_quasiseparated, pa-prop:curve_integral, pa-cor:curve_sections}
  Let \(C\) be a smooth curve over \(k\), i.e.\ proper, smooth of relative dimension one and
  geometrically irreducible over \(k\), as in Chapter~\ref{pa-chap:Challenge}. Combining the
  above: \(C\) is smooth over \(k\) (\ref{pa-lem:smooth_of_relDim_one}), geometrically reduced
  (\ref{pa-lem:relDim_geometrically_reduced}) and geometrically integral
  (\ref{pa-thm:relDim_geometrically_integral}); its underlying space is irreducible
  (\ref{pa-lem:irreducible_of_geomirred}), quasi-compact (\ref{pa-lem:curve_compact}) and
  quasi-separated (\ref{pa-lem:curve_quasiseparated}); \(C\) is an integral scheme
  (\ref{pa-prop:curve_integral}); and \(\Gamma(C, \struct C) = k\) (\ref{pa-cor:curve_sections}).
\end{remark}

\section{The projective line and its charts}

Let \(A = k[X_0, X_1]\) be the polynomial ring in two variables, graded by total degree. For a
homogeneous element \(f \in A\) we write \(A_{(f)}\) for the degree-zero part of the
localization \(A_f\); its elements are fractions \(a/f^n\) with \(a\) homogeneous of degree
\(n \deg f\).

\begin{definition}[The projective line]
  \label{pa-def:P1}
  \lean{AlgebraicGeometry.P1}
  \leanok
  \dcref{custom}

  The projective line over \(k\) is \(\PP^1 = \Proj A\), the Proj of the graded ring
  \(A = k[X_0, X_1]\).
\end{definition}

\begin{definition}[Degree zero]
  \label{pa-def:P1_gradeZero}
  \lean{AlgebraicGeometry.P1.gradeZeroAlgEquiv}
  \leanok
  \dcref{custom}

  The inclusion of constants \(k \to A_0\) into the degree-zero part of \(A\) is an
  isomorphism of \(k\)-algebras.
\end{definition}
\begin{proof}
  \leanok
  It is injective because a nonzero constant polynomial is nonzero, and surjective because a
  polynomial which is homogeneous of degree zero is constant.
\end{proof}

\begin{definition}[Structure morphism]
  \label{pa-def:P1_structureMap}
  \lean{AlgebraicGeometry.P1.structureMap}
  \leanok
  \dcref{custom}

  \uses{pa-def:P1, pa-def:P1_gradeZero}
  The structure morphism \(\PP^1 \to \Spec k\) is the canonical morphism
  \(\Proj A \to \Spec A_0\) followed by the isomorphism \(\Spec A_0 \cong \Spec k\) induced by
  \ref{pa-def:P1_gradeZero}.
\end{definition}

\begin{theorem}[Properness of the projective line]
  \label{pa-thm:P1_proper}
  \lean{AlgebraicGeometry.P1.isProper_structureMap}
  \leanok
  \dcref{custom}

  \uses{pa-def:P1_structureMap}
  \(\PP^1\) is proper over \(k\).
\end{theorem}
\begin{proof}
  \leanok
  \uses{pa-def:P1_gradeZero}
  \(A\) is a finitely generated algebra over its degree-zero part \(A_0 = k\) (generated by
  \(X_0\) and \(X_1\)), and the Proj of a graded ring which is finitely generated over its
  degree-zero part is proper over the spectrum of the degree-zero part. Composing with the
  isomorphism \(\Spec A_0 \cong \Spec k\) preserves properness.
\end{proof}

\begin{definition}[The projective line as a scheme over \(k\)]
  \label{pa-def:P1_asOver}
  \lean{AlgebraicGeometry.P1.asOver}
  \leanok
  \dcref{custom}

  \uses{pa-def:P1_structureMap}
  \(\PP^1\), regarded as a scheme over \(k\) via the structure morphism of
  \ref{pa-def:P1_structureMap}.
\end{definition}

\begin{definition}[The standard charts]
  \label{pa-def:P1_chartOpen}
  \lean{AlgebraicGeometry.P1.chartOpen}
  \leanok
  \dcref{custom}

  \uses{pa-def:P1}
  For \(i \in \{0,1\}\), the chart \(U_i = D_+(X_i) \subseteq \PP^1\) is the basic open locus
  where the homogeneous coordinate \(X_i\) does not vanish.
\end{definition}

\begin{lemma}[The charts are affine]
  \label{pa-lem:P1_chartOpen_affine}
  \lean{AlgebraicGeometry.P1.isAffineOpen_chartOpen}
  \leanok
  \dcref{custom}

  \uses{pa-def:P1_chartOpen}
  Each chart \(U_i\) is an affine open of \(\PP^1\).
\end{lemma}
\begin{proof}
  \leanok
  \(X_i\) is homogeneous of degree \(1 > 0\), and the basic open of \(\Proj\) at a homogeneous
  element of positive degree is affine.
\end{proof}

\begin{lemma}[The charts cover]
  \label{pa-lem:P1_chartOpen_sup}
  \lean{AlgebraicGeometry.P1.chartOpen_sup}
  \leanok
  \dcref{custom}

  \uses{pa-def:P1_chartOpen}
  \(U_0 \cup U_1 = \PP^1\).
\end{lemma}
\begin{proof}
  \leanok
  A point of \(\Proj A\) is a homogeneous prime ideal \(\mathfrak p\) of \(A\) not containing
  the irrelevant ideal \(A_+\). Every monomial of positive degree is divisible by \(X_0\) or
  by \(X_1\), so if both \(X_0 \in \mathfrak p\) and \(X_1 \in \mathfrak p\) then
  \(A_+ \subseteq \mathfrak p\), a contradiction. Hence every point avoids some \(X_i\), i.e.\
  lies in \(U_0 \cup U_1\).
\end{proof}

\begin{lemma}[The chart overlap]
  \label{pa-lem:P1_chartOpen_inf}
  \lean{AlgebraicGeometry.P1.chartOpen_inf}
  \leanok
  \dcref{custom}

  \uses{pa-def:P1_chartOpen}
  \(U_0 \cap U_1 = D_+(X_0 X_1)\).
\end{lemma}
\begin{proof}
  \leanok
  A homogeneous prime avoids the product \(X_0 X_1\) if and only if it avoids both factors.
\end{proof}

\begin{definition}[Affine models of the charts]
  \label{pa-def:P1_chartIota}
  \lean{AlgebraicGeometry.P1.chartι}
  \leanok
  \dcref{custom}

  \uses{pa-def:P1, pa-def:P1_chartOpen}
  For \(i \in \{0,1\}\), the canonical morphism \(\Spec A_{(X_i)} \to \PP^1\) associated with
  the basic open at the degree-one element \(X_i\); it is an open immersion.
\end{definition}

\begin{lemma}[Image of the affine model]
  \label{pa-lem:P1_chartIota_range}
  \lean{AlgebraicGeometry.P1.opensRange_chartι}
  \leanok
  \dcref{custom}

  \uses{pa-def:P1_chartIota, pa-def:P1_chartOpen}
  The image of \(\Spec A_{(X_i)} \to \PP^1\) is the chart \(U_i\); thus
  \(U_i \cong \Spec A_{(X_i)}\).
\end{lemma}
\begin{proof}
  \leanok
  This is the standard description of the basic opens of \(\Proj\): the homogeneous primes
  avoiding \(X_i\) correspond to the primes of the degree-zero localization \(A_{(X_i)}\).
\end{proof}

\begin{lemma}[The affine models lie over \(k\)]
  \label{pa-lem:P1_chartIota_structureMap}
  \lean{AlgebraicGeometry.P1.chartι_structureMap}
  \leanok
  \dcref{custom}

  \uses{pa-def:P1_chartIota, pa-def:P1_structureMap}
  The composite \(\Spec A_{(X_i)} \to \PP^1 \to \Spec k\) is the morphism induced by the
  \(k\)-algebra structure of \(A_{(X_i)}\); that is, the affine models are morphisms of
  schemes over \(k\).
\end{lemma}
\begin{proof}
  \leanok
  \uses{pa-def:P1_gradeZero}
  The canonical morphism \(\Proj A \to \Spec A_0\), restricted to the affine chart, is induced
  by the ring homomorphism \(A_0 \to A_{(X_i)}\); composing with the isomorphism
  \(k \cong A_0\) of \ref{pa-def:P1_gradeZero} gives the claim.
\end{proof}

\begin{definition}[Chart coordinates]
  \label{pa-def:P1_chartCoord}
  \lean{AlgebraicGeometry.P1.chartCoord}
  \leanok
  \dcref{custom}

  \uses{pa-def:P1}
  For \(i, j \in \{0,1\}\), the chart coordinate is the degree-zero fraction
  \(t_{ij} = X_j / X_i \in A_{(X_i)}\).
\end{definition}

\begin{definition}[Dehomogenization]
  \label{pa-def:P1_dehomogenize}
  \lean{AlgebraicGeometry.P1.dehomogenize}
  \leanok
  \dcref{custom}

  For \(i \in \{0,1\}\), the dehomogenization at \(X_i\) is the \(k\)-algebra homomorphism
  \(\delta_i : k[X_0, X_1] \to k[t]\) with \(\delta_i(X_i) = 1\) and \(\delta_i(X_j) = t\) for
  \(j \neq i\).
\end{definition}

\begin{definition}[From polynomials to the chart ring]
  \label{pa-def:P1_polyToAway}
  \lean{AlgebraicGeometry.P1.polyToAway}
  \leanok
  \dcref{custom}

  \uses{pa-def:P1_chartCoord}
  For \(i, j \in \{0,1\}\), the \(k\)-algebra homomorphism \(k[t] \to A_{(X_i)}\) given by
  evaluation at the chart coordinate, \(t \mapsto t_{ij}\).
\end{definition}

\begin{definition}[From the chart ring to polynomials]
  \label{pa-def:P1_awayToPoly}
  \lean{AlgebraicGeometry.P1.awayToPoly}
  \leanok
  \dcref{custom}

  \uses{pa-def:P1_dehomogenize}
  For \(i \in \{0,1\}\), the \(k\)-algebra homomorphism \(A_{(X_i)} \to k[t]\) sending a
  degree-zero fraction \(a / X_i^n\) (with \(a\) homogeneous of degree \(n\)) to
  \(\delta_i(a)\).
\end{definition}
\begin{proof}
  \leanok
  Well-definedness: \(\delta_i(X_i) = 1\) is a unit of \(k[t]\), so \(\delta_i\) extends to
  the localization \(A_{X_i} \to k[t]\) by the universal property of localization; on the
  degree-zero subring its value on \(a / X_i^n\) is
  \(\delta_i(a)\,\delta_i(X_i)^{-n} = \delta_i(a)\).
\end{proof}

\begin{lemma}[One composite is the identity]
  \label{pa-lem:P1_awayToPoly_polyToAway}
  \lean{AlgebraicGeometry.P1.awayToPoly_comp_polyToAway}
  \leanok
  \dcref{custom}

  \uses{pa-def:P1_polyToAway, pa-def:P1_awayToPoly}
  For \(i \neq j\), the composite \(k[t] \to A_{(X_i)} \to k[t]\) of
  \ref{pa-def:P1_polyToAway} followed by \ref{pa-def:P1_awayToPoly} is the identity.
\end{lemma}
\begin{proof}
  \leanok
  Both sides are \(k\)-algebra homomorphisms on \(k[t]\), so it suffices to check the
  generator: \(t \mapsto t_{ij} = X_j / X_i \mapsto \delta_i(X_j) = t\).
\end{proof}

\begin{lemma}[Monomial fractions are powers of the coordinate]
  \label{pa-lem:P1_away_monomial}
  \lean{AlgebraicGeometry.P1.away_mk_prod_eq}
  \leanok
  \dcref{custom}

  \uses{pa-def:P1_chartCoord}
  Let \(i \neq j\) and let \(e_0, e_1 \ge 0\) with \(e_0 + e_1 = a\). Then in \(A_{(X_i)}\)
  \[
    \frac{X_0^{e_0} X_1^{e_1}}{X_i^{a}} \;=\; t_{ij}^{\,e_j}.
  \]
\end{lemma}
\begin{proof}
  \leanok
  Since \(X_0^{e_0} X_1^{e_1} = X_i^{e_i} X_j^{e_j}\) and \(a = e_i + e_j\), the left-hand
  side is \(X_i^{e_i} X_j^{e_j} / X_i^{e_i + e_j} = X_j^{e_j} / X_i^{e_j}\), which is
  \(t_{ij}^{\,e_j}\).
\end{proof}

\begin{lemma}[Surjectivity onto the chart ring]
  \label{pa-lem:P1_polyToAway_surjective}
  \lean{AlgebraicGeometry.P1.polyToAway_surjective}
  \leanok
  \dcref{custom}

  \uses{pa-def:P1_polyToAway}
  For \(i \neq j\), the homomorphism \(k[t] \to A_{(X_i)}\) of \ref{pa-def:P1_polyToAway} is
  surjective.
\end{lemma}
\begin{proof}
  \leanok
  \uses{pa-lem:P1_away_monomial, pa-def:P1_gradeZero}
  Since \(A\) is generated over \(A_0\) by the degree-one elements \(X_0, X_1\), the
  degree-zero localization \(A_{(X_i)}\) is spanned, as a module over \(A_0 = k\)
  (\ref{pa-def:P1_gradeZero}), by the fractions \(m / X_i^a\) where \(m\) ranges over the
  monomials \(X_0^{e_0} X_1^{e_1}\) of degree \(a\). By \ref{pa-lem:P1_away_monomial} each such
  fraction equals \(t_{ij}^{\,e_j}\), which is the image of \(t^{e_j}\). The image of a ring
  homomorphism is closed under addition and under multiplication by scalars from \(k\), so it
  is all of \(A_{(X_i)}\).
\end{proof}

\begin{lemma}[The other composite is the identity]
  \label{pa-lem:P1_polyToAway_awayToPoly}
  \lean{AlgebraicGeometry.P1.polyToAway_comp_awayToPoly}
  \leanok
  \dcref{custom}

  \uses{pa-def:P1_polyToAway, pa-def:P1_awayToPoly}
  For \(i \neq j\), the composite \(A_{(X_i)} \to k[t] \to A_{(X_i)}\) of
  \ref{pa-def:P1_awayToPoly} followed by \ref{pa-def:P1_polyToAway} is the identity.
\end{lemma}
\begin{proof}
  \leanok
  \uses{pa-lem:P1_polyToAway_surjective, pa-lem:P1_awayToPoly_polyToAway}
  By \ref{pa-lem:P1_polyToAway_surjective} every element of \(A_{(X_i)}\) is of the form
  \(p(t_{ij})\) with \(p \in k[t]\), and on such elements the claim follows from
  \ref{pa-lem:P1_awayToPoly_polyToAway}.
\end{proof}

\begin{definition}[The chart identification]
  \label{pa-def:P1_awayAlgEquiv}
  \lean{AlgebraicGeometry.P1.awayAlgEquiv}
  \leanok
  \dcref{custom}

  \uses{pa-def:P1_polyToAway, pa-def:P1_awayToPoly}
  For \(i \neq j\), the isomorphism of \(k\)-algebras
  \(\alpha_i : A_{(X_i)} \cong k[t]\) given by the mutually inverse homomorphisms of
  \ref{pa-def:P1_awayToPoly} and \ref{pa-def:P1_polyToAway}; the chart coordinate \(t_{ij}\)
  corresponds to \(t\).
\end{definition}
\begin{proof}
  \leanok
  \uses{pa-lem:P1_awayToPoly_polyToAway, pa-lem:P1_polyToAway_awayToPoly}
  The two composites are the identity by \ref{pa-lem:P1_awayToPoly_polyToAway} and
  \ref{pa-lem:P1_polyToAway_awayToPoly}.
\end{proof}

\section{The chart overlap and Laurent polynomials}

Write \(k[T, T^{-1}]\) for the ring of Laurent polynomials over \(k\).

\begin{lemma}[Laurent span]
  \label{pa-lem:laurent_span}
  \lean{LaurentPolynomial.exists_toLaurent_add_aeval}
  \leanok
  \dcref{custom}

  Every Laurent polynomial \(f \in R[T, T^{-1}]\) over a commutative (semi)ring \(R\) can be
  written as \(f = p(T) + q(T^{-1})\) with \(p, q \in R[t]\) ordinary polynomials.
\end{lemma}
\begin{proof}
  \leanok
  Split \(f\) into its monomials and collect those with nonnegative exponent into \(p\) and
  those with negative exponent into \(q\).
\end{proof}

\begin{definition}[Restriction from the left chart]
  \label{pa-def:P1_awayToOverlapLeft}
  \lean{AlgebraicGeometry.P1.awayToOverlapLeft}
  \leanok
  \dcref{custom}

  \uses{pa-def:P1}
  The canonical ring homomorphism \(\lambda : A_{(X_0)} \to A_{(X_0 X_1)}\), sending
  \(a / X_0^{\,n}\) to \(a X_1^{\,n} / (X_0 X_1)^{n}\).
\end{definition}

\begin{definition}[Restriction from the right chart]
  \label{pa-def:P1_awayToOverlapRight}
  \lean{AlgebraicGeometry.P1.awayToOverlapRight}
  \leanok
  \dcref{custom}

  \uses{pa-def:P1}
  The canonical ring homomorphism \(\rho : A_{(X_1)} \to A_{(X_0 X_1)}\), sending
  \(a / X_1^{\,n}\) to \(a X_0^{\,n} / (X_0 X_1)^{n}\).
\end{definition}

\begin{lemma}[The two coordinates are mutually inverse on the overlap]
  \label{pa-lem:P1_coord_mul}
  \lean{AlgebraicGeometry.P1.awayToOverlap_mul_eq_one}
  \leanok
  \dcref{custom}

  \uses{pa-def:P1_awayToOverlapLeft, pa-def:P1_awayToOverlapRight, pa-def:P1_chartCoord}
  In \(A_{(X_0 X_1)}\) one has \(\lambda(t_{01}) \cdot \rho(t_{10}) = 1\).
\end{lemma}
\begin{proof}
  \leanok
  Computing with the formulas of \ref{pa-def:P1_awayToOverlapLeft} and
  \ref{pa-def:P1_awayToOverlapRight}:
  \(\lambda(t_{01}) = X_1^2 / (X_0 X_1)\) and \(\rho(t_{10}) = X_0^2 / (X_0 X_1)\), so
  \[
    \lambda(t_{01}) \cdot \rho(t_{10})
      = \frac{X_0^2 X_1^2}{(X_0 X_1)^2} = 1 .
  \]
\end{proof}

\begin{lemma}[The overlap ring is a localization of the chart ring]
  \label{pa-lem:P1_overlap_isLocalization}
  \lean{AlgebraicGeometry.P1.instAwayAwayNatMvPolynomialFinOfNatSubmoduleHomogeneousSubmoduleXChartCoordHMul}
  \leanok
  \dcref{custom}

  \uses{pa-def:P1_awayToOverlapLeft, pa-def:P1_chartCoord}
  The homomorphism \(\lambda\) of \ref{pa-def:P1_awayToOverlapLeft} identifies
  \(A_{(X_0 X_1)}\) with the localization of \(A_{(X_0)}\) at the powers of the chart
  coordinate \(t_{01}\).
\end{lemma}
\begin{proof}
  \leanok
  This is a general property of degree-zero localizations: for homogeneous elements \(f, g\)
  of degree one, \(A_{(fg)}\) is the localization of \(A_{(f)}\) away from the degree-zero
  fraction \(g/f\). Here \(f = X_0\), \(g = X_1\) and \(g/f = t_{01}\).
\end{proof}

\begin{definition}[The overlap identification]
  \label{pa-def:P1_overlapAlgEquiv}
  \lean{AlgebraicGeometry.P1.overlapAlgEquiv}
  \leanok
  \dcref{custom}

  \uses{pa-def:P1}
  The isomorphism of \(k\)-algebras \(\theta : A_{(X_0 X_1)} \cong k[T, T^{-1}]\), with
  \(T\) the image of the coordinate \(t_{01}\).
\end{definition}
\begin{proof}
  \leanok
  \uses{pa-lem:P1_overlap_isLocalization, pa-def:P1_awayAlgEquiv}
  By \ref{pa-lem:P1_overlap_isLocalization}, \(A_{(X_0X_1)}\) is the localization of
  \(A_{(X_0)}\) at the powers of \(t_{01}\). The isomorphism
  \(\alpha_0 : A_{(X_0)} \cong k[t]\) of \ref{pa-def:P1_awayAlgEquiv} carries \(t_{01}\) to
  \(t\), hence induces an isomorphism of the two localizations
  \(A_{(X_0X_1)} \cong k[t]_{t} = k[T, T^{-1}]\) (localizations at corresponding
  multiplicative sets are canonically isomorphic). The map is \(k\)-linear because
  \(\alpha_0\) is.
\end{proof}

\begin{lemma}[Left restriction is \(t \mapsto T\)]
  \label{pa-lem:P1_overlap_res_left}
  \lean{AlgebraicGeometry.P1.overlapAlgEquiv_awayToOverlapLeft}
  \leanok
  \dcref{custom}

  \uses{pa-def:P1_overlapAlgEquiv, pa-def:P1_awayAlgEquiv, pa-def:P1_awayToOverlapLeft}
  For every \(z \in A_{(X_0)}\) one has \(\theta(\lambda(z)) = p(T)\), where
  \(p = \alpha_0(z) \in k[t]\). In other words, under the identifications
  \(\alpha_0\) and \(\theta\), the restriction \(\lambda\) is the canonical inclusion
  \(k[t] \to k[T, T^{-1}]\), \(t \mapsto T\).
\end{lemma}
\begin{proof}
  \leanok
  \uses{pa-lem:P1_overlap_isLocalization}
  By construction (\ref{pa-def:P1_overlapAlgEquiv}), \(\theta\) is the comparison of
  localizations induced by \(\alpha_0\); under such a comparison the localization map
  \(\lambda\) of \ref{pa-lem:P1_overlap_isLocalization} corresponds to the localization map
  \(k[t] \to k[T, T^{-1}]\).
\end{proof}

\begin{lemma}[Right restriction is \(t \mapsto T^{-1}\)]
  \label{pa-lem:P1_overlap_res_right}
  \lean{AlgebraicGeometry.P1.overlapAlgEquiv_awayToOverlapRight}
  \leanok
  \dcref{custom}

  \uses{pa-def:P1_overlapAlgEquiv, pa-def:P1_awayAlgEquiv, pa-def:P1_awayToOverlapRight}
  For every \(z \in A_{(X_1)}\) one has \(\theta(\rho(z)) = q(T^{-1})\), where
  \(q = \alpha_1(z) \in k[t]\). In other words, under the identifications \(\alpha_1\) and
  \(\theta\), the restriction \(\rho\) is the homomorphism \(k[t] \to k[T, T^{-1}]\),
  \(t \mapsto T^{-1}\).
\end{lemma}
\begin{proof}
  \leanok
  \uses{pa-lem:P1_coord_mul, pa-lem:P1_overlap_res_left, pa-lem:P1_polyToAway_surjective}
  By \ref{pa-lem:P1_overlap_res_left} the image of \(\lambda(t_{01})\) under \(\theta\) is
  \(T\), and by \ref{pa-lem:P1_coord_mul} \(\lambda(t_{01})\,\rho(t_{10}) = 1\); applying
  \(\theta\) gives \(\theta(\rho(t_{10})) = T^{-1}\). By
  \ref{pa-lem:P1_polyToAway_surjective} (for the indices \((1,0)\)) every \(z \in A_{(X_1)}\)
  is a polynomial in \(t_{10}\) with coefficients in \(k\); both sides of the claim are
  \(k\)-algebra homomorphisms in \(z\), and they agree on \(t_{10}\).
\end{proof}

\begin{lemma}[Laurent span on the overlap ring]
  \label{pa-lem:P1_overlap_span}
  \lean{AlgebraicGeometry.P1.exists_awayToOverlap_add}
  \leanok
  \dcref{custom}

  \uses{pa-def:P1_awayToOverlapLeft, pa-def:P1_awayToOverlapRight}
  Every \(w \in A_{(X_0 X_1)}\) can be written as \(w = \lambda(z_0) + \rho(z_1)\) with
  \(z_0 \in A_{(X_0)}\) and \(z_1 \in A_{(X_1)}\).
\end{lemma}
\begin{proof}
  \leanok
  \uses{pa-lem:laurent_span, pa-def:P1_overlapAlgEquiv, pa-lem:P1_overlap_res_left,
    pa-lem:P1_overlap_res_right, pa-def:P1_awayAlgEquiv}
  Write \(\theta(w) = p(T) + q(T^{-1})\) by \ref{pa-lem:laurent_span} and set
  \(z_0 = \alpha_0^{-1}(p)\), \(z_1 = \alpha_1^{-1}(q)\). By
  \ref{pa-lem:P1_overlap_res_left} and \ref{pa-lem:P1_overlap_res_right},
  \(\theta(\lambda(z_0) + \rho(z_1)) = p(T) + q(T^{-1}) = \theta(w)\), and \(\theta\) is
  injective.
\end{proof}

\section{Section rings and the Laurent chart pair}

\begin{lemma}[The overlap is affine]
  \label{pa-lem:P1_overlap_affine}
  \lean{AlgebraicGeometry.P1.isAffineOpen_overlap}
  \leanok
  \dcref{custom}

  \uses{pa-lem:P1_chartOpen_inf}
  The chart overlap \(U_0 \cap U_1 = D_+(X_0 X_1)\) is an affine open of \(\PP^1\).
\end{lemma}
\begin{proof}
  \leanok
  \(X_0 X_1\) is homogeneous of degree \(2 > 0\), and the basic open of \(\Proj\) at a
  homogeneous element of positive degree is affine.
\end{proof}

\begin{definition}[Sections on the left chart]
  \label{pa-def:P1_chartSectionsEquiv0}
  \lean{AlgebraicGeometry.P1.chartSectionsEquiv₀}
  \leanok
  \dcref{custom}

  \uses{pa-def:P1_chartOpen, pa-def:P1_chartIota, pa-def:P1_awayAlgEquiv}
  The ring isomorphism \(\gamma_0 : \Gamma(\PP^1, U_0) \cong k[t]\): the canonical
  identification \(\Gamma(\PP^1, D_+(X_0)) \cong A_{(X_0)}\) (sections of \(\Proj\) over a
  basic open; equivalently, \(U_0 = \Spec A_{(X_0)}\) via \ref{pa-def:P1_chartIota}) followed by
  \(\alpha_0\) of \ref{pa-def:P1_awayAlgEquiv}.
\end{definition}

\begin{definition}[Sections on the right chart]
  \label{pa-def:P1_chartSectionsEquiv1}
  \lean{AlgebraicGeometry.P1.chartSectionsEquiv₁}
  \leanok
  \dcref{custom}

  \uses{pa-def:P1_chartOpen, pa-def:P1_chartIota, pa-def:P1_awayAlgEquiv}
  The ring isomorphism \(\gamma_1 : \Gamma(\PP^1, U_1) \cong k[t]\), defined as in
  \ref{pa-def:P1_chartSectionsEquiv0} with the identification
  \(\Gamma(\PP^1, D_+(X_1)) \cong A_{(X_1)}\) followed by \(\alpha_1\).
\end{definition}

\begin{definition}[Sections on the overlap]
  \label{pa-def:P1_overlapSectionsEquiv}
  \lean{AlgebraicGeometry.P1.overlapSectionsEquiv}
  \leanok
  \dcref{custom}

  \uses{pa-lem:P1_chartOpen_inf, pa-def:P1_overlapAlgEquiv}
  The ring isomorphism \(\gamma_{01} : \Gamma(\PP^1, U_0 \cap U_1) \cong k[T, T^{-1}]\): the
  canonical identification \(\Gamma(\PP^1, D_+(X_0X_1)) \cong A_{(X_0X_1)}\) followed by
  \(\theta\) of \ref{pa-def:P1_overlapAlgEquiv}.
\end{definition}

\begin{lemma}[Restriction of sections from the left chart]
  \label{pa-lem:P1_res_left}
  \lean{AlgebraicGeometry.P1.overlapSectionsEquiv_res_left}
  \leanok
  \dcref{custom}

  \uses{pa-def:P1_chartSectionsEquiv0, pa-def:P1_overlapSectionsEquiv}
  For every section \(a \in \Gamma(\PP^1, U_0)\),
  \[
    \gamma_{01}\bigl(a|_{U_0 \cap U_1}\bigr) = p(T), \qquad p = \gamma_0(a);
  \]
  under the identifications, restriction from the left chart to the overlap is
  \(t \mapsto T\).
\end{lemma}
\begin{proof}
  \leanok
  \uses{pa-lem:P1_overlap_res_left}
  The identifications \(\Gamma(\PP^1, D_+(f)) \cong A_{(f)}\) are natural with respect to
  inclusions of basic opens: restriction of sections from \(D_+(X_0)\) to \(D_+(X_0X_1)\)
  corresponds to the ring homomorphism \(\lambda\). Conclude by
  \ref{pa-lem:P1_overlap_res_left}.
\end{proof}

\begin{lemma}[Restriction of sections from the right chart]
  \label{pa-lem:P1_res_right}
  \lean{AlgebraicGeometry.P1.overlapSectionsEquiv_res_right}
  \leanok
  \dcref{custom}

  \uses{pa-def:P1_chartSectionsEquiv1, pa-def:P1_overlapSectionsEquiv}
  For every section \(b \in \Gamma(\PP^1, U_1)\),
  \[
    \gamma_{01}\bigl(b|_{U_0 \cap U_1}\bigr) = q(T^{-1}), \qquad q = \gamma_1(b);
  \]
  under the identifications, restriction from the right chart to the overlap is
  \(t \mapsto T^{-1}\).
\end{lemma}
\begin{proof}
  \leanok
  \uses{pa-lem:P1_overlap_res_right}
  As in \ref{pa-lem:P1_res_left}, restriction of sections corresponds to \(\rho\); conclude by
  \ref{pa-lem:P1_overlap_res_right}.
\end{proof}

\begin{definition}[Laurent chart pair]
  \label{pa-def:LaurentChartPair}
  \lean{AlgebraicGeometry.LaurentChartPair}
  \leanok
  \dcref{custom}

  \uses{pa-def:P1}
  A \emph{Laurent chart pair} on \(\PP^1\) consists of open subsets
  \(V_0, V_1, V_{01} \subseteq \PP^1\) with \(V_{01} \subseteq V_0\),
  \(V_{01} \subseteq V_1\) and \(V_0 \cap V_1 = V_{01}\), all three affine, covering the
  projective line (\(V_0 \cup V_1 = \PP^1\)), together with ring isomorphisms
  \[
    \Gamma_0 : \Gamma(\PP^1, V_0) \cong k[t], \quad
    \Gamma_1 : \Gamma(\PP^1, V_1) \cong k[t], \quad
    \Gamma_{01} : \Gamma(\PP^1, V_{01}) \cong k[T, T^{-1}],
  \]
  such that, under these identifications, restriction \(V_{01} \subseteq V_0\) is
  \(t \mapsto T\) and restriction \(V_{01} \subseteq V_1\) is \(t \mapsto T^{-1}\).
\end{definition}

\begin{theorem}[Laurent span for a chart pair]
  \label{pa-thm:laurentChartPair_span}
  \lean{AlgebraicGeometry.LaurentChartPair.exists_res_add_res}
  \leanok
  \dcref{custom}

  \uses{pa-def:LaurentChartPair}
  Let \((V_0, V_1, V_{01}, \Gamma_0, \Gamma_1, \Gamma_{01})\) be a Laurent chart pair on
  \(\PP^1\). Every section \(z \in \Gamma(\PP^1, V_{01})\) is the sum of a restriction from
  \(V_0\) and a restriction from \(V_1\): there are \(a \in \Gamma(\PP^1, V_0)\) and
  \(b \in \Gamma(\PP^1, V_1)\) with \(z = a|_{V_{01}} + b|_{V_{01}}\).
\end{theorem}
\begin{proof}
  \leanok
  \uses{pa-lem:laurent_span}
  Write \(\Gamma_{01}(z) = p(T) + q(T^{-1})\) by \ref{pa-lem:laurent_span} and set
  \(a = \Gamma_0^{-1}(p)\), \(b = \Gamma_1^{-1}(q)\). By the two restriction conditions,
  \(\Gamma_{01}(a|_{V_{01}} + b|_{V_{01}}) = p(T) + q(T^{-1}) = \Gamma_{01}(z)\), and
  \(\Gamma_{01}\) is injective.
\end{proof}

\begin{definition}[The canonical Laurent chart pair]
  \label{pa-def:P1_laurentChartPair}
  \lean{AlgebraicGeometry.P1.laurentChartPair}
  \leanok
  \dcref{custom}

  \uses{pa-def:LaurentChartPair, pa-def:P1_chartOpen, pa-def:P1_chartSectionsEquiv0,
    pa-def:P1_chartSectionsEquiv1, pa-def:P1_overlapSectionsEquiv}
  The Laurent chart pair on \(\PP^1\) given by the standard charts: \(V_0 = U_0\),
  \(V_1 = U_1\), \(V_{01} = D_+(X_0X_1)\), with the identifications \(\gamma_0, \gamma_1,
  \gamma_{01}\) of \ref{pa-def:P1_chartSectionsEquiv0}--\ref{pa-def:P1_overlapSectionsEquiv}.
\end{definition}
\begin{proof}
  \leanok
  \uses{pa-lem:P1_chartOpen_affine, pa-lem:P1_chartOpen_sup, pa-lem:P1_chartOpen_inf,
    pa-lem:P1_overlap_affine, pa-lem:P1_res_left, pa-lem:P1_res_right}
  The overlap conditions and affineness are \ref{pa-lem:P1_chartOpen_inf},
  \ref{pa-lem:P1_chartOpen_affine} and \ref{pa-lem:P1_overlap_affine}; the covering condition is
  \ref{pa-lem:P1_chartOpen_sup}; the two restriction conditions are \ref{pa-lem:P1_res_left} and
  \ref{pa-lem:P1_res_right}.
\end{proof}

\section{Unramified finite-type algebras over Dedekind domains}

The results of this section are pure commutative algebra, stated in full generality. The
first theorem strengthens the classical statement that a domain which is unramified and
\emph{module-finite} over a Dedekind domain is Dedekind: module-finiteness is relaxed to
finite type. The point of the generalization is where incomparability of primes comes from
--- for module-finite algebras it follows from integrality (lying over/incomparability),
while here it follows from quasi-finiteness, which every unramified algebra essentially of
finite type enjoys.

\begin{theorem}[Unramified finite-type domains over a Dedekind domain, DVR form]
  \label{pa-thm:unramified_finiteType_dedekind_dvr}
  \lean{IsDedekindDomainDvr.of_formallyUnramified_of_finiteType}
  \leanok
  \dcref{custom}

  Let \(A\) be a Dedekind domain and \(B\) an integral domain which is an \(A\)-algebra of
  finite type and formally unramified over \(A\). Then \(B\) is Noetherian, and the
  localization of \(B\) at every nonzero prime ideal is a discrete valuation ring.
\end{theorem}
\begin{proof}
  \leanok
  \(B\) is Noetherian as a finite-type algebra over the Noetherian ring \(A\) (Hilbert basis
  theorem). Let \(\mathfrak q\) be a nonzero prime of \(B\) and set
  \(\mathfrak p = \mathfrak q \cap A\). The localization \(B_{\mathfrak q}\) is a Noetherian
  local domain and not a field (its maximal ideal is nonzero because \(\mathfrak q\) is), so
  it suffices to show that its maximal ideal \(\mathfrak q B_{\mathfrak q}\) is principal.

  \emph{Step 1: \(\mathfrak p \neq 0\).} An algebra which is formally unramified and of
  finite type is quasi-finite, and over a quasi-finite algebra two comparable primes of
  \(B\) lying over the same prime of \(A\) coincide. If \(\mathfrak p\) were zero, the
  comparable primes \(0 \subseteq \mathfrak q\) of \(B\) would both contract to \(0\) in
  \(A\), forcing \(\mathfrak q = 0\), a contradiction. Since \(A\) is Dedekind, the nonzero
  prime \(\mathfrak p\) is maximal, and \(A_{\mathfrak p}\) is a discrete valuation ring; in
  particular \(\mathfrak p A_{\mathfrak p}\) is principal.

  \emph{Step 2: \(\mathfrak q\) is minimal over \(\mathfrak p B\).} The fiber ring
  \(B/\mathfrak p B\) is an algebra over the residue field \(A/\mathfrak p\) which is again
  formally unramified (base change); every module over a field is free, and an unramified
  algebra which is free as a module is module-finite. Hence \(B/\mathfrak p B\) is a finite
  module over the field \(A/\mathfrak p\), in particular an Artinian ring, so all of its
  primes are minimal; applying this to the image of \(\mathfrak q\) shows that
  \(\mathfrak q\) is a minimal prime over \(\mathfrak p B\).

  \emph{Step 3: \(\mathfrak q B_{\mathfrak q} = \mathfrak p B_{\mathfrak q}\).} Since
  \(\mathfrak p\) is maximal and \(B\) is unramified over \(A\), the extension
  \(\mathfrak p B\) is a radical ideal of \(B\) (the fiber ring \(B/\mathfrak p B\) is
  reduced, being unramified and essentially of finite type over a field). Localization
  commutes with radicals, so \(\mathfrak p B_{\mathfrak q}\) is a radical ideal of
  \(B_{\mathfrak q}\); and because \(\mathfrak q\) is minimal over \(\mathfrak p B\)
  (Step~2), the radical of \(\mathfrak p B_{\mathfrak q}\) is exactly the maximal ideal
  \(\mathfrak q B_{\mathfrak q}\). A radical ideal equals its own radical, so
  \(\mathfrak p B_{\mathfrak q} = \mathfrak q B_{\mathfrak q}\).

  \emph{Step 4: conclusion.} By Step~3 and Step~1, the maximal ideal
  \(\mathfrak q B_{\mathfrak q}\) is the extension of the principal ideal
  \(\mathfrak p A_{\mathfrak p} \subseteq A_{\mathfrak p}\) along
  \(A_{\mathfrak p} \to B_{\mathfrak q}\), hence principal. A Noetherian local domain which
  is not a field and whose maximal ideal is principal is a discrete valuation ring.
\end{proof}

\begin{corollary}[Unramified finite-type domains over a Dedekind domain]
  \label{pa-cor:unramified_finiteType_dedekind}
  \lean{IsDedekindDomain.of_formallyUnramified_of_finiteType}
  \leanok
  \dcref{custom}

  \uses{pa-thm:unramified_finiteType_dedekind_dvr}
  Let \(A\) be a Dedekind domain and \(B\) an integral domain which is an \(A\)-algebra of
  finite type and formally unramified over \(A\). Then \(B\) is a Dedekind domain.
\end{corollary}
\begin{proof}
  \leanok
  A Noetherian domain whose localizations at all nonzero primes are discrete valuation
  rings is a Dedekind domain; apply \ref{pa-thm:unramified_finiteType_dedekind_dvr}.
\end{proof}

\begin{lemma}[The polynomial ring in one variable is a principal ideal ring]
  \label{pa-lem:mvpoly_fin_one_pid}
  \lean{MvPolynomial.isPrincipalIdealRing_fin_one}
  \leanok
  \dcref{custom}

  Let \(k\) be a field. The polynomial ring over \(k\) on a one-element set of variables is
  a principal ideal ring.
\end{lemma}
\begin{proof}
  \leanok
  It is isomorphic as a ring to the ordinary one-variable polynomial ring \(k[t]\), which
  is a principal ideal domain (Euclidean division), and being a principal ideal ring
  transports along a surjective ring homomorphism.
\end{proof}

\begin{lemma}[The polynomial ring in one variable is Dedekind]
  \label{pa-lem:mvpoly_fin_one_dedekind}
  \lean{MvPolynomial.isDedekindDomain_fin_one}
  \leanok
  \dcref{custom}

  \uses{pa-lem:mvpoly_fin_one_pid}
  Let \(k\) be a field. The polynomial ring over \(k\) on a one-element set of variables is
  a Dedekind domain.
\end{lemma}
\begin{proof}
  \leanok
  A principal ideal domain is a Dedekind domain.
\end{proof}

\begin{theorem}[Structure theorem for standard smooth algebras]
  \label{pa-thm:standard_smooth_etale_structure}
  \lean{Algebra.IsStandardSmoothOfRelativeDimension.exists_etale_mvPolynomial}
  \mathlibok
  \source{stacks-project}

  \dcref{custom}
  Let \(R\) be a commutative ring and \(S\) an \(R\)-algebra which is standard smooth of
  relative dimension \(n\): it admits a presentation
  \(S = R[x_1,\dots,x_m]/(f_1,\dots,f_c)\) with \(m - c = n\) in which the Jacobian matrix
  \((\partial f_u/\partial x_v)_{u,v \le c}\) has invertible determinant in \(S\). Then the
  structure map factors as
  \[
    R \longrightarrow R[y_1, \dots, y_n] \longrightarrow S
  \]
  with the second homomorphism \'etale. (For such a presentation the module of K\"ahler
  differentials \(\Omega_{S/R}\) is free on the images of
  \(dx_{c+1}, \dots, dx_m\); consequently \(S\) is unramified, hence \'etale, over the
  polynomial subring on those last \(n\) coordinates.)
\end{theorem}

\begin{theorem}[Standard smooth domains of relative dimension one are Dedekind]
  \label{pa-thm:standard_smooth_dim_one_dedekind}
  \lean{Algebra.IsStandardSmoothOfRelativeDimension.isDedekindDomain}
  \leanok
  \dcref{custom}

  \uses{pa-cor:unramified_finiteType_dedekind, pa-lem:mvpoly_fin_one_dedekind,
    pa-thm:standard_smooth_etale_structure}
  Let \(k\) be a field and \(B\) an integral domain which is a standard smooth
  \(k\)-algebra of relative dimension \(1\). Then \(B\) is a Dedekind domain.
\end{theorem}
\begin{proof}
  \leanok
  By the structure theorem \ref{pa-thm:standard_smooth_etale_structure}, \(B\) is \'etale over
  a polynomial ring over \(k\) in one variable, which is a Dedekind domain
  (\ref{pa-lem:mvpoly_fin_one_dedekind}). An \'etale algebra is formally unramified and of
  finite type, so \ref{pa-cor:unramified_finiteType_dedekind} applies.
\end{proof}

\begin{theorem}[Standard smooth algebras of relative dimension one contain a transcendental]
  \label{pa-thm:standard_smooth_transcendental}
  \lean{Algebra.IsStandardSmoothOfRelativeDimension.exists_transcendental}
  \leanok
  \dcref{custom}

  \uses{pa-thm:standard_smooth_etale_structure}
  Let \(k\) be a field and \(B\) a nontrivial standard smooth \(k\)-algebra of relative
  dimension \(1\). Then \(B\) contains an element transcendental over \(k\).
\end{theorem}
\begin{proof}
  \leanok
  By the structure theorem \ref{pa-thm:standard_smooth_etale_structure}, there is an \'etale
  homomorphism \(g : k[y] \to B\) from a one-variable polynomial ring; we claim that the
  \'etale coordinate \(t = g(y)\) is transcendental over \(k\). Note that no dimension
  theory enters: the argument is pure torsion-freeness.

  First, \(g\) is injective. An \'etale algebra is flat, and a flat module over a domain is
  torsion-free; so \(B\) is a torsion-free \(k[y]\)-module. If \(g(P) = g(Q)\) for
  \(P \neq Q\), then the nonzero element \(P - Q\) of the domain \(k[y]\) would satisfy
  \((P - Q) \cdot 1_B = g(P - Q) = 0\) with \(1_B \neq 0\) (\(B\) is nontrivial),
  contradicting torsion-freeness.

  Now suppose \(P(t) = 0\) for some nonzero \(P \in k[T]\). Evaluating \(P\) at \(t = g(y)\)
  is the image under \(g\) of the evaluation of \(P\) at \(y\), so \(g(P(y)) = 0\); by
  injectivity \(P(y) = 0\) in \(k[y]\), contradicting the transcendence of the variable
  \(y\) over \(k\). Hence \(t\) is transcendental over \(k\).
\end{proof}

\section{Stalks, specializations and finiteness on curve-like schemes}

A scheme covered by affine opens whose section rings are Dedekind domains behaves like a
smooth curve: its stalks are valuation rings, and its specialization order has height one.
This section proves the pointwise statements from the single hypothesis ``the point lies in
an affine open with Dedekind section ring'', derives the finiteness of proper closed
subsets, and discharges the hypotheses for a scheme smooth of relative dimension one over a
field.

\begin{theorem}[Stalks over Dedekind section rings are valuation rings]
  \label{pa-thm:affine_dedekind_stalk_valuation}
  \lean{AlgebraicGeometry.IsAffineOpen.valuationRing_stalk}
  \leanok
  \dcref{custom}

  Let \(X\) be an integral scheme, \(V \subseteq X\) an affine open whose ring of sections
  \(\Gamma(X, V)\) is a Dedekind domain, and \(x \in V\). Then the stalk
  \(\struct{X,x}\) is a valuation ring.
\end{theorem}
\begin{proof}
  \leanok
  The stalk at \(x\) is the localization of \(\Gamma(X, V)\) at the prime ideal
  \(\mathfrak q\) corresponding to \(x\) under \(V \cong \Spec \Gamma(X, V)\). If
  \(\mathfrak q = 0\), its complement is the set of nonzero elements of the domain
  \(\Gamma(X, V)\), so the stalk is the fraction field --- a field, and every field is a
  valuation ring. If \(\mathfrak q \neq 0\), the localization of a Dedekind domain at a
  nonzero prime is a discrete valuation ring, and every discrete valuation ring is a
  valuation ring.
\end{proof}

\begin{theorem}[Height-one specialization over Dedekind section rings]
  \label{pa-thm:affine_dedekind_specializes}
  \lean{AlgebraicGeometry.IsAffineOpen.specializes_eq_genericPoint_or_eq}
  \leanok
  \dcref{custom}

  Let \(X\) be a scheme whose underlying space is irreducible, \(V \subseteq X\) an affine
  open with \(\Gamma(X, V)\) a Dedekind domain, and \(x \in V\). Then every generalization
  of \(x\) is the generic point of \(X\) or \(x\) itself.
\end{theorem}
\begin{proof}
  \leanok
  Let \(y\) be a generalization of \(x\). Open sets are stable under generalization, so
  \(y \in V\). Under the open immersion \(\Spec \Gamma(X, V) \cong V \subseteq X\) the
  points \(x, y\) correspond to primes \(\mathfrak q_x, \mathfrak q_y\) of
  \(\Gamma(X, V)\) with \(\mathfrak q_y \subseteq \mathfrak q_x\) (specialization order of
  the spectrum). If \(\mathfrak q_y = 0\), then \(y\) is the image of the generic point of
  \(\Spec \Gamma(X, V)\), and the generic point of a nonempty open subscheme of an
  irreducible scheme is the generic point of the whole scheme. If \(\mathfrak q_y \neq 0\),
  then \(\mathfrak q_y\) is maximal (nonzero primes of a Dedekind domain are maximal), so
  \(\mathfrak q_y = \mathfrak q_x\) and \(y = x\).
\end{proof}

\begin{lemma}[Non-generic points are closed]
  \label{pa-lem:height_one_closed_points}
  \lean{AlgebraicGeometry.Scheme.isClosed_singleton_of_forall_specializes}
  \leanok
  \dcref{custom}

  Let \(X\) be a scheme whose underlying space is irreducible and whose specialization
  order has height one: every generalization of any point is the generic point or the
  point itself. Then every point other than the generic point is closed.
\end{lemma}
\begin{proof}
  \leanok
  Let \(x \neq \xi\) (\(\xi\) the generic point) and let \(y\) be a point of the closure of
  \(\{x\}\), i.e.\ \(x\) is a generalization of \(y\). By hypothesis \(x = \xi\) or
  \(x = y\); the first is excluded, so \(y = x\). Hence the closure of \(\{x\}\) is
  \(\{x\}\).
\end{proof}

\begin{theorem}[Finiteness of proper closed subsets]
  \label{pa-thm:height_one_finite_closed}
  \lean{AlgebraicGeometry.Scheme.finite_of_isClosed_of_notMem_genericPoint}
  \leanok
  \dcref{custom}

  \uses{pa-lem:height_one_closed_points}
  Let \(X\) be a scheme whose underlying space is irreducible and Noetherian and whose
  specialization order has height one. Then every closed subset \(Z \subseteq X\) not
  containing the generic point is finite.
\end{theorem}
\begin{proof}
  \leanok
  A closed subset of a Noetherian space is a finite union of irreducible closed subsets.
  Schemes are sober, so each irreducible closed subset \(T\) in such a decomposition of
  \(Z\) is the closure of a unique point \(z \in T \subseteq Z\). Since the generic point
  does not lie in \(Z\), \(z\) is not the generic point, so \(\{z\}\) is closed by
  \ref{pa-lem:height_one_closed_points} and \(T = \overline{\{z\}} = \{z\}\). Hence \(Z\) is a
  finite union of singletons.
\end{proof}

\begin{theorem}[Dedekind charts on a curve]
  \label{pa-thm:relDim1_dedekind_charts}
  \lean{AlgebraicGeometry.SmoothOfRelativeDimension.exists_isDedekindDomain_section}
  \leanok
  \dcref{custom}

  \uses{pa-thm:standard_smooth_dim_one_dedekind}
  Let \(k\) be a field and \(X\) an integral scheme, smooth of relative dimension one over
  \(\Spec k\). Then every point of \(X\) lies in an affine open \(V\) whose ring of
  sections \(\Gamma(X, V)\) is a Dedekind domain.
\end{theorem}
\begin{proof}
  \leanok
  By definition of smoothness of relative dimension one, a point \(x\) has an affine open
  neighbourhood \(V\), mapping into an affine open \(U\) of \(\Spec k\), such that
  \(\Gamma(U) \to \Gamma(X, V)\) is standard smooth of relative dimension \(1\). The base
  \(\Spec k\) has a single point, and \(U\) is nonempty (it receives the image of \(x\)),
  so \(U = \Spec k\) and \(\Gamma(U) \cong k\). The ring \(\Gamma(X, V)\) is an integral
  domain, since \(V\) is a nonempty open of the integral scheme \(X\). Apply
  \ref{pa-thm:standard_smooth_dim_one_dedekind}.
\end{proof}

\begin{theorem}[The stalks of a curve are valuation rings]
  \label{pa-thm:relDim1_stalks_valuation}
  \lean{AlgebraicGeometry.SmoothOfRelativeDimension.valuationRing_stalk}
  \leanok
  \dcref{custom}

  \uses{pa-thm:relDim1_dedekind_charts, pa-thm:affine_dedekind_stalk_valuation}
  Let \(k\) be a field and \(X\) an integral scheme, smooth of relative dimension one over
  \(\Spec k\). Then every stalk of the structure sheaf of \(X\) is a valuation ring.
\end{theorem}
\begin{proof}
  \leanok
  Combine \ref{pa-thm:relDim1_dedekind_charts} with
  \ref{pa-thm:affine_dedekind_stalk_valuation}.
\end{proof}

\begin{theorem}[Height-one specialization on a curve]
  \label{pa-thm:relDim1_specializes}
  \lean{AlgebraicGeometry.SmoothOfRelativeDimension.specializes_eq_genericPoint_or_eq}
  \leanok
  \dcref{custom}

  \uses{pa-thm:relDim1_dedekind_charts, pa-thm:affine_dedekind_specializes}
  Let \(k\) be a field and \(X\) an integral scheme, smooth of relative dimension one over
  \(\Spec k\). Then every generalization of a point of \(X\) is the generic point or the
  point itself.
\end{theorem}
\begin{proof}
  \leanok
  Combine \ref{pa-thm:relDim1_dedekind_charts} with \ref{pa-thm:affine_dedekind_specializes};
  the underlying space of an integral scheme is irreducible.
\end{proof}

\begin{theorem}[Existence of a transcendental rational function]
  \label{pa-thm:relDim1_transcendental}
  \lean{AlgebraicGeometry.SmoothOfRelativeDimension.exists_transcendental_functionField}
  \leanok
  \dcref{custom}

  \uses{pa-thm:standard_smooth_transcendental}
  Let \(k\) be a field and \(X\) an integral scheme, smooth of relative dimension one over
  \(\Spec k\). Regard the function field \(K(X)\) --- the stalk of the structure sheaf at
  the generic point --- as a ring under \(k\) via the structure morphism (global sections
  followed by the germ at the generic point). Then \(K(X)\) contains an element \(f_0\)
  transcendental over \(k\): no nonzero polynomial with coefficients in \(k\) vanishes at
  \(f_0\).
\end{theorem}
\begin{proof}
  \leanok
  Apply the definition of smoothness of relative dimension one at the generic point
  \(\xi\): there is an affine open \(V \ni \xi\) such that \(k \to \Gamma(X, V)\) is
  standard smooth of relative dimension \(1\) (the base affine open is all of \(\Spec k\),
  as in the proof of \ref{pa-thm:relDim1_dedekind_charts}). The ring \(\Gamma(X, V)\) is
  nontrivial (\(V\) is nonempty), so by \ref{pa-thm:standard_smooth_transcendental} it
  contains an element \(t\) transcendental over \(k\); let \(f_0\) be the germ of \(t\) at
  \(\xi\).

  On an integral scheme, the germ map \(\Gamma(X, V) \to K(X)\) at the generic point is
  injective. Let \(P \in k[T]\) be nonzero. Evaluating \(P\) at \(f_0\) through the
  structure map \(k \to K(X)\) is the germ of the evaluation of \(P\) at \(t\) through
  \(k \to \Gamma(X, V)\) (the structure map factors through \(\Gamma(X, V)\), and germs
  commute with polynomial evaluation). The latter is nonzero by transcendence of \(t\)
  (with coefficients moved into \(\Gamma(X, V)\) along the injective map \(k \to
  \Gamma(X, V)\)), so by injectivity of the germ map \(P(f_0) \neq 0\).
\end{proof}

\section{The curve-morphism dichotomy}
\label{pa-sec:curve_morphism_dichotomy}

A morphism from a complete nonsingular curve to a curve is either constant or a finite
surjection. This section proves the bundle form of that dichotomy for the proper smooth
geometrically irreducible curve bundles of the challenge, together with its two general
ingredients: a morphism between proper bundles is itself proper, and a morphism whose
set-range is a single closed point factors through the spectrum of the residue field of
that point.

\begin{theorem}[Morphisms of proper bundles are proper]
  \label{pa-thm:overHom_isProper}
  \lean{AlgebraicGeometry.isProper_left, AlgebraicGeometry.locallyOfFiniteType_left,
    AlgebraicGeometry.quasiCompact_left}
  \leanok
  \dcref{custom}

  Let \(S\) be a scheme and \(D, E\) schemes over \(S\), with \(D \to S\) proper and
  \(E \to S\) separated. Then every morphism \(h : D \to E\) over \(S\) is proper; in
  particular \(h\) is locally of finite type and quasi-compact.
\end{theorem}
\begin{proof}
  \leanok
  The composite of \(h\) with the structure morphism \(E \to S\) is the structure
  morphism \(D \to S\), which is proper. Properness cancels along separated morphisms:
  if \(g \circ f\) is proper and \(g\) is separated, then \(f\) is proper --- \(f\)
  factors as its graph \((\mathrm{id}, f)\), a closed immersion because it is a base
  change of the diagonal of the separated \(g\), followed by a projection which is a
  base change of the proper \(g \circ f\); closed immersions are proper, and properness
  is stable under base change and composition. Apply this with \(f = h\) and \(g\) the
  structure morphism of \(E\). Properness comprises finite type, so \(h\) is locally of
  finite type; and \(h\) is universally closed, hence quasi-compact (the preimage of a
  quasi-compact open under a universally closed morphism is quasi-compact).
\end{proof}

\begin{lemma}[The kernel comparison at a one-point range]
  \label{pa-lem:ker_fromSpecResidueField_le}
  \lean{AlgebraicGeometry.Scheme.Hom.ker_fromSpecResidueField_le}
  \leanok
  \dcref{custom}

  Let \(f : X \to Y\) be a quasi-compact morphism of schemes with \(X\) reduced whose
  set-range is a single point \(y\): \(f(X) = \{y\}\). Then the kernel ideal sheaf of
  the canonical morphism \(\Spec \kappa(y) \to Y\) is contained in the kernel ideal
  sheaf of \(f\): for every open \(U \subseteq Y\), every section \(s \in \Gamma(Y, U)\)
  pulling back to \(0\) on \(\Spec \kappa(y)\) pulls back to \(0\) in
  \(\Gamma(X, f^{-1}U)\).
\end{lemma}
\begin{proof}
  \leanok
  Let \(s \in \Gamma(Y, U)\) pull back to \(0\) on \(\Spec \kappa(y)\). First,
  \(y \notin D(s)\), the basic open of \(s\): the point \(y\) lies in the range of
  \(\Spec \kappa(y) \to Y\), and if its preimage point lay in the preimage of \(D(s)\),
  that preimage --- the basic open of the pulled-back section, i.e.\ of \(0\) --- would
  be nonempty, which is absurd. Since the range of \(f\) is \(\{y\}\) and \(y \notin
  D(s)\), the preimage \(f^{-1}D(s)\) is empty. But \(f^{-1}D(s)\) is the basic open of
  the pulled-back section \(f^{\sharp}s \in \Gamma(X, f^{-1}U)\), so \(f^{\sharp}s\)
  has empty basic open: its image in every stalk of \(X\) lies in the maximal ideal, so
  on each affine chart it lies in every prime ideal, i.e.\ it is nilpotent. On the
  reduced scheme \(X\) it is therefore \(0\).
\end{proof}

\begin{theorem}[Factorization through the residue field of a closed point]
  \label{pa-thm:hom_factors_fromSpecResidueField}
  \lean{AlgebraicGeometry.Scheme.Hom.exists_comp_fromSpecResidueField_eq}
  \leanok
  \dcref{custom}

  \uses{pa-lem:ker_fromSpecResidueField_le}
  Let \(f : X \to Y\) be a quasi-compact morphism of schemes with \(X\) reduced whose
  set-range is a single \emph{closed} point \(y\): \(\{y\}\) is closed and
  \(f(X) = \{y\}\). Then \(f\) factors through the canonical morphism from the spectrum
  of the residue field: there is a morphism \(q : X \to \Spec \kappa(y)\) with
  \[
    f \;=\; \bigl(\Spec \kappa(y) \to Y\bigr) \circ q .
  \]
\end{theorem}
\begin{proof}
  \leanok
  Since \(\{y\}\) is closed, the canonical morphism \(\Spec \kappa(y) \to Y\) is a
  closed immersion (it presents the reduced induced structure on the closed set
  \(\{y\}\): a singleton is irreducible, its reduced subscheme is the spectrum of a
  field, namely the residue field of its point). A morphism factors through a closed
  immersion --- necessarily uniquely --- exactly when the kernel ideal sheaf of the
  immersion is contained in its own kernel ideal sheaf; this containment is
  \ref{pa-lem:ker_fromSpecResidueField_le}.
\end{proof}

\begin{theorem}[Proper closed subsets of the curve bundle are finite]
  \label{pa-thm:curve_finite_closed}
  \lean{AlgebraicGeometry.finite_of_isClosed_of_notMem_genericPoint_curve}
  \leanok
  \dcref{custom}

  Let \(K\) be a field and \(X\) an integral scheme, smooth of relative dimension one
  over \(\Spec K\), whose underlying space is quasi-compact. Then every closed subset
  \(Z \subseteq X\) not containing the generic point is finite.
\end{theorem}
\begin{proof}
  \leanok
  \uses{pa-thm:height_one_finite_closed, pa-thm:relDim1_specializes}
  The underlying space of \(X\) is irreducible, and its specialization order has height
  one: every generalization of a point is the generic point or the point itself
  (\ref{pa-thm:relDim1_specializes}). The space is moreover Noetherian: \(X\) is locally of
  finite type over the field \(K\) (smoothness of relative dimension one includes local
  finite presentation), so each affine chart is the spectrum of a finitely generated
  \(K\)-algebra --- a Noetherian ring by the Hilbert basis theorem --- and a
  quasi-compact space covered by finitely many Noetherian opens is Noetherian. Apply
  the finiteness of proper closed subsets (\ref{pa-thm:height_one_finite_closed}).
\end{proof}

\begin{theorem}[The curve-morphism dichotomy]
  \label{pa-thm:curveHom_isFinite_or_constant}
  \lean{AlgebraicGeometry.curveHom_isFinite_or_constant}
  \leanok
  \source{hartshorne-algebraic-geometry:page-0154}

  \dcref{custom}
  Let \(K\) be a field and let \(D, E\) be objects of \(\Over(\Spec K)\), both proper,
  smooth of relative dimension one and geometrically irreducible over \(K\). Then every
  morphism \(h : D \to E\) of \(\Over(\Spec K)\) satisfies at least one of the
  following: either (1) \(h\) is surjective and finite, or (2) there is a point \(x\)
  of \(E\) with \(\{x\}\) closed and \(h(D) = \{x\}\). The two cases may overlap: when
  the space of \(E\) is a single point, both hold.
\end{theorem}
\begin{proof}
  \leanok
  \uses{pa-thm:overHom_isProper, pa-thm:curve_finite_closed, pa-thm:closedpoint_isClosed,
    pa-lem:genericPoint_eq_of_surjective, pa-thm:smooth_geometrically_reduced,
    pa-prop:curve_integral, pa-lem:irreducible_of_geomirred}
  Both bundles are integral schemes: smoothness gives geometric reducedness
  (\ref{pa-thm:smooth_geometrically_reduced}), which together with geometric
  irreducibility gives integrality (\ref{pa-prop:curve_integral}); in particular the
  underlying spaces are irreducible (\ref{pa-lem:irreducible_of_geomirred}) and sober,
  with generic points \(\eta_D\), \(\eta_E\).

  \emph{The image is the closure of a point.} By \ref{pa-thm:overHom_isProper} the
  underlying morphism of \(h\) is proper, hence a closed map, so the set-range
  \(h(D)\) is closed in \(E\); and it is irreducible, being the continuous image of
  the irreducible space \(D\). A nonempty irreducible closed subset of a sober space is
  the closure of a unique point: \(h(D) = \overline{\{\zeta\}}\).

  \emph{Case \(\zeta \neq \eta_E\).} The singleton \(\{\zeta\}\) is closed
  (\ref{pa-thm:closedpoint_isClosed}), so \(h(D) = \overline{\{\zeta\}} = \{\zeta\}\), and
  case (2) holds with \(x = \zeta\).

  \emph{Case \(\zeta = \eta_E\).} Then \(h(D) = \overline{\{\eta_E\}} = E\), so \(h\)
  is surjective; and a surjection matches the generic points, \(h(\eta_D) = \eta_E\)
  (\ref{pa-lem:genericPoint_eq_of_surjective}).

  Suppose first that some point \(z \neq \eta_D\) also maps to \(\eta_E\). Then
  \(\{z\}\) is closed in \(D\) (\ref{pa-thm:closedpoint_isClosed}), and \(h\), a closed
  map, carries it to the closed set \(\{\eta_E\}\). A closed subset containing the
  generic point of an irreducible space is the whole space, so \(E\) has the single
  point \(\eta_E\); since the range is all of \(E = \{\eta_E\}\), case (2) holds with
  \(x = \eta_E\) --- this is the overlap of the two cases.

  Otherwise the fiber of \(\eta_E\) is exactly \(\{\eta_D\}\), and every fiber of
  \(h\) is finite: the fiber over a point \(y \neq \eta_E\) is closed --- \(\{y\}\) is
  closed by \ref{pa-thm:closedpoint_isClosed} --- and avoids \(\eta_D\), whose image is
  \(\eta_E\); so it is finite by \ref{pa-thm:curve_finite_closed}. Since \(h\) is locally
  of finite type and quasi-compact (\ref{pa-thm:overHom_isProper}), finiteness of all
  fibers makes it locally quasi-finite; and a proper, locally quasi-finite morphism is
  finite, by Zariski's main theorem. So case (1) holds.
\end{proof}

\section{The topology of the projective line and morphisms into it}

Two groups of results about \(\PP^1\) feed the construction of the finite morphism from
the curve: the topology of \(\PP^1\) (irreducibility, height-one specialization, closed
points, separatedness), and the construction of morphisms \(\Spec A \to \PP^1\) from a
single element of a ring \(A\) under \(k\), with the gluing identity
\([1 : u] = [u^{-1} : 1]\) for units.

\begin{lemma}[The chart models are domains]
  \label{pa-lem:P1_away_domains}
  \lean{AlgebraicGeometry.P1.instIsDomainAwayNatMvPolynomialFinOfNatSubmoduleHomogeneousSubmoduleX,
    AlgebraicGeometry.P1.instIsDomainAwayNatMvPolynomialFinOfNatSubmoduleHomogeneousSubmoduleX_1,
    AlgebraicGeometry.P1.instNontrivialAwayNatMvPolynomialFinOfNatSubmoduleHomogeneousSubmoduleHMulX}
  \leanok
  \dcref{custom}

  \uses{pa-def:P1_awayAlgEquiv, pa-def:P1_overlapAlgEquiv}
  The chart rings \(A_{(X_0)}\) and \(A_{(X_1)}\) are integral domains, and the overlap
  ring \(A_{(X_0X_1)}\) is nontrivial.
\end{lemma}
\begin{proof}
  \leanok
  Transport along the isomorphisms \(\alpha_i : A_{(X_i)} \cong k[t]\) of
  \ref{pa-def:P1_awayAlgEquiv} and \(\theta : A_{(X_0X_1)} \cong k[T,T^{-1}]\) of
  \ref{pa-def:P1_overlapAlgEquiv}: polynomial rings over a field are domains, and the Laurent
  polynomial ring is nontrivial.
\end{proof}

\begin{lemma}[The chart section rings are principal ideal rings]
  \label{pa-lem:P1_chartSections_pid}
  \lean{AlgebraicGeometry.P1.isPrincipalIdealRing_chartSections}
  \leanok
  \dcref{custom}

  \uses{pa-def:P1_chartSectionsEquiv0, pa-def:P1_chartSectionsEquiv1}
  For \(i \in \{0,1\}\), the ring \(\Gamma(\PP^1, U_i)\) is a principal ideal ring.
\end{lemma}
\begin{proof}
  \leanok
  Being a principal ideal ring transports along the surjective homomorphism
  \(\gamma_i^{-1} : k[t] \to \Gamma(\PP^1, U_i)\)
  (\ref{pa-def:P1_chartSectionsEquiv0}--\ref{pa-def:P1_chartSectionsEquiv1}), and \(k[t]\) is a
  principal ideal domain.
\end{proof}

\begin{lemma}[The chart section rings are domains]
  \label{pa-lem:P1_chartSections_domain}
  \lean{AlgebraicGeometry.P1.isDomain_chartSections}
  \leanok
  \dcref{custom}

  \uses{pa-def:P1_chartSectionsEquiv0, pa-def:P1_chartSectionsEquiv1}
  For \(i \in \{0,1\}\), the ring \(\Gamma(\PP^1, U_i)\) is an integral domain.
\end{lemma}
\begin{proof}
  \leanok
  Transport along \(\gamma_i : \Gamma(\PP^1, U_i) \cong k[t]\).
\end{proof}

\begin{theorem}[The chart section rings are Dedekind]
  \label{pa-thm:P1_chartSections_dedekind}
  \lean{AlgebraicGeometry.P1.isDedekindDomain_chartSections}
  \leanok
  \dcref{custom}

  \uses{pa-lem:P1_chartSections_pid, pa-lem:P1_chartSections_domain}
  For \(i \in \{0,1\}\), the ring \(\Gamma(\PP^1, U_i)\) is a Dedekind domain.
\end{theorem}
\begin{proof}
  \leanok
  A principal ideal domain is a Dedekind domain; combine
  \ref{pa-lem:P1_chartSections_pid} and \ref{pa-lem:P1_chartSections_domain}.
\end{proof}

\begin{lemma}[Every nonempty open meets the left chart]
  \label{pa-lem:P1_dense_chart}
  \lean{AlgebraicGeometry.P1.exists_mem_chartOpen_zero_of_isOpen}
  \leanok
  \dcref{custom}

  \uses{pa-lem:P1_chartOpen_sup, pa-lem:P1_chartIota_range, pa-lem:P1_away_domains}
  Every nonempty open subset \(W \subseteq \PP^1\) contains a point of the chart \(U_0\).
\end{lemma}
\begin{proof}
  \leanok
  Let \(w \in W\). The charts cover (\ref{pa-lem:P1_chartOpen_sup}), so \(w \in U_0\) or
  \(w \in U_1\); in the first case we are done. In the second, work in the affine model
  \(U_1 \cong \Spec A_{(X_1)}\) (\ref{pa-lem:P1_chartIota_range}), an irreducible space since
  \(A_{(X_1)}\) is a domain (\ref{pa-lem:P1_away_domains}). The preimage of \(W\) there is a
  nonempty open (it contains the preimage of \(w\)). The preimage of \(U_0\) there is open
  and nonempty: the overlap \(D_+(X_0X_1)\) is contained in both charts and is nonempty,
  because its coordinate ring \(A_{(X_0X_1)}\) is nontrivial (\ref{pa-lem:P1_away_domains}).
  Two nonempty opens of an irreducible space intersect; the image in \(\PP^1\) of a common
  point lies in \(W \cap U_0\).
\end{proof}

\begin{theorem}[Irreducibility of the projective line]
  \label{pa-thm:P1_irreducible}
  \lean{AlgebraicGeometry.P1.instIrreducibleSpaceCarrierCarrierCommRingCat}
  \leanok
  \dcref{custom}

  \uses{pa-lem:P1_dense_chart, pa-lem:P1_chartIota_range, pa-lem:P1_away_domains}
  The underlying topological space of \(\PP^1\) is irreducible (in particular nonempty).
\end{theorem}
\begin{proof}
  \leanok
  \(\PP^1\) is nonempty: the chart \(U_0 \cong \Spec A_{(X_0)}\) is the spectrum of a
  nontrivial ring. Let \(u, v\) be nonempty opens of \(\PP^1\). By
  \ref{pa-lem:P1_dense_chart} both meet \(U_0\), so their preimages in the affine model
  \(\Spec A_{(X_0)}\) (\ref{pa-lem:P1_chartIota_range}) are nonempty opens; the spectrum of
  the domain \(A_{(X_0)}\) (\ref{pa-lem:P1_away_domains}) is irreducible, so the preimages
  intersect, and the image of a common point lies in \(u \cap v\).
\end{proof}

\begin{lemma}[The generic point through the chart]
  \label{pa-lem:P1_generic_point_chart}
  \lean{AlgebraicGeometry.P1.chartι_base_genericPoint}
  \leanok
  \dcref{custom}

  \uses{pa-def:P1_chartIota, pa-thm:P1_irreducible}
  The affine model \(\Spec A_{(X_0)} \to \PP^1\) sends the generic point of
  \(\Spec A_{(X_0)}\) to the generic point of \(\PP^1\).
\end{lemma}
\begin{proof}
  \leanok
  An open immersion with nonempty source into an irreducible space has dense image, and
  the image of the generic point of the source is a generic point of the image, hence of
  the closure of the image, which is the whole space.
\end{proof}

\begin{lemma}[A second point of the projective line]
  \label{pa-lem:P1_exists_ne_generic}
  \lean{AlgebraicGeometry.P1.exists_ne_genericPoint}
  \leanok
  \dcref{custom}

  \uses{pa-def:P1_awayAlgEquiv, pa-lem:P1_generic_point_chart}
  \(\PP^1\) has a point distinct from its generic point.
\end{lemma}
\begin{proof}
  \leanok
  \(A_{(X_0)} \cong k[t]\) (\ref{pa-def:P1_awayAlgEquiv}) is not a field, so it has a nonzero
  prime ideal \(\mathfrak p\). The corresponding point of the chart differs from the
  generic point: the generic point corresponds to the zero ideal
  (\ref{pa-lem:P1_generic_point_chart}), and the chart immersion is injective.
\end{proof}

\begin{theorem}[Height-one specialization on the projective line]
  \label{pa-thm:P1_specializes}
  \lean{AlgebraicGeometry.P1.specializes_eq_genericPoint_or_eq}
  \leanok
  \dcref{custom}

  \uses{pa-lem:P1_chartOpen_sup, pa-lem:P1_chartOpen_affine, pa-thm:P1_chartSections_dedekind,
    pa-thm:affine_dedekind_specializes}
  Every generalization of a point of \(\PP^1\) is the generic point or the point itself.
\end{theorem}
\begin{proof}
  \leanok
  Every point lies in one of the two charts (\ref{pa-lem:P1_chartOpen_sup}), which are affine
  (\ref{pa-lem:P1_chartOpen_affine}) with Dedekind section rings
  (\ref{pa-thm:P1_chartSections_dedekind}); apply \ref{pa-thm:affine_dedekind_specializes}.
\end{proof}

\begin{lemma}[Non-generic points of the projective line are closed]
  \label{pa-lem:P1_closed_points}
  \lean{AlgebraicGeometry.P1.isClosed_singleton_of_ne_genericPoint}
  \leanok
  \dcref{custom}

  \uses{pa-lem:height_one_closed_points, pa-thm:P1_specializes}
  Every point of \(\PP^1\) other than the generic point is closed.
\end{lemma}
\begin{proof}
  \leanok
  Immediate from \ref{pa-lem:height_one_closed_points} and \ref{pa-thm:P1_specializes}.
\end{proof}

\begin{theorem}[Separatedness of the projective line]
  \label{pa-thm:P1_separated}
  \lean{AlgebraicGeometry.P1.instIsSeparated}
  \leanok
  \dcref{custom}

  \uses{pa-thm:P1_proper}
  \(\PP^1\) is a separated scheme.
\end{theorem}
\begin{proof}
  \leanok
  The structure morphism \(\PP^1 \to \Spec k\) is proper (\ref{pa-thm:P1_proper}), hence
  separated, and \(\Spec k\), being affine, is a separated scheme; a scheme separated over
  a separated scheme is separated.
\end{proof}

\begin{definition}[Evaluation of the chart ring]
  \label{pa-def:P1_chartEval}
  \lean{AlgebraicGeometry.P1.chartEval}
  \leanok
  \dcref{custom}

  \uses{pa-def:P1_awayToPoly}
  Let \(A\) be a commutative ring under \(k\), i.e.\ equipped with a ring homomorphism
  \(\sigma : k \to A\), let \(i \in \{0,1\}\) and \(a \in A\). The evaluation of the chart
  ring at \(a\) is the ring homomorphism
  \[
    A_{(X_i)} \longrightarrow A, \qquad
    z \longmapsto \bigl(\text{the polynomial } \alpha_i(z) \in k[t]\bigr)
      \text{ evaluated at } t = a \text{ with coefficients mapped by } \sigma
  \]
  (composite of \ref{pa-def:P1_awayToPoly} with polynomial evaluation). It sends the chart
  coordinate \(t_{ij}\) (\(j \neq i\)) to \(a\) and a constant \(c \in k\) to \(\sigma(c)\).
\end{definition}

\begin{definition}[Morphisms to the projective line through a chart]
  \label{pa-def:P1_fromSpecChart}
  \lean{AlgebraicGeometry.P1.fromSpecChart}
  \leanok
  \dcref{custom}

  \uses{pa-def:P1_chartEval, pa-def:P1_chartIota}
  With \(\sigma : k \to A\), \(i \in \{0,1\}\) and \(a \in A\) as in
  \ref{pa-def:P1_chartEval}: the morphism \(\Spec A \to \PP^1\) obtained by applying
  \(\Spec\) to the evaluation \(A_{(X_i)} \to A\) at \(a\) and composing with the affine
  model \(\Spec A_{(X_i)} \to \PP^1\) of \ref{pa-def:P1_chartIota}. For \(i = 0\) this is the
  ``point'' \([1 : a]\); for \(i = 1\) it is \([a : 1]\).
\end{definition}

\begin{lemma}[The chart morphisms lie over the base]
  \label{pa-lem:P1_fromSpecChart_over}
  \lean{AlgebraicGeometry.P1.fromSpecChart_structureMap}
  \leanok
  \dcref{custom}

  \uses{pa-def:P1_fromSpecChart, pa-def:P1_structureMap, pa-lem:P1_chartIota_structureMap}
  The composite of the morphism of \ref{pa-def:P1_fromSpecChart} with the structure morphism
  \(\PP^1 \to \Spec k\) is \(\Spec \sigma : \Spec A \to \Spec k\).
\end{lemma}
\begin{proof}
  \leanok
  \uses{pa-def:P1_chartEval}
  The affine model lies over \(k\) (\ref{pa-lem:P1_chartIota_structureMap}), and the
  evaluation homomorphism of \ref{pa-def:P1_chartEval} restricts to \(\sigma\) on the constants
  \(k \to A_{(X_i)}\); apply \(\Spec\) to this equality of ring homomorphisms.
\end{proof}

\begin{lemma}[Naturality of the chart morphisms]
  \label{pa-lem:P1_fromSpecChart_natural}
  \lean{AlgebraicGeometry.P1.SpecMap_fromSpecChart}
  \leanok
  \dcref{custom}

  \uses{pa-def:P1_fromSpecChart}
  Let \(g : A \to A'\) be a homomorphism of rings under \(k\) (so \(A'\) is under \(k\)
  via \(g \circ \sigma\)), \(i \in \{0,1\}\) and \(a \in A\). Then the composite
  \(\Spec A' \xrightarrow{\Spec g} \Spec A \to \PP^1\) of \ref{pa-def:P1_fromSpecChart} at
  \(a\) is the morphism of \ref{pa-def:P1_fromSpecChart} at \(g(a) \in A'\).
\end{lemma}
\begin{proof}
  \leanok
  Polynomial evaluation commutes with \(g\): evaluating at \(a\) and applying \(g\) is
  evaluating at \(g(a)\) with coefficients mapped by \(g \circ \sigma\). Apply \(\Spec\).
\end{proof}

\begin{lemma}[The chart restrictions factor the overlap model]
  \label{pa-lem:P1_SpecMap_overlap_chartIota}
  \lean{AlgebraicGeometry.P1.SpecMap_awayToOverlapLeft_chartι,
    AlgebraicGeometry.P1.SpecMap_awayToOverlapRight_chartι}
  \leanok
  \dcref{custom}

  \uses{pa-def:P1_awayToOverlapLeft, pa-def:P1_awayToOverlapRight, pa-def:P1_chartIota}
  Applying \(\Spec\) to the restriction homomorphism \(\lambda : A_{(X_0)} \to
  A_{(X_0X_1)}\) (resp.\ \(\rho : A_{(X_1)} \to A_{(X_0X_1)}\)) of
  \ref{pa-def:P1_awayToOverlapLeft} (resp.\ \ref{pa-def:P1_awayToOverlapRight}) and composing
  with the affine model of the chart \(U_0\) (resp.\ \(U_1\)) yields the affine model
  \(\Spec A_{(X_0X_1)} \to \PP^1\) of the overlap \(D_+(X_0X_1)\).
\end{lemma}
\begin{proof}
  \leanok
  This is the compatibility of the basic-open affine models of \(\Proj\) with inclusions
  of basic opens: the inclusion \(D_+(fg) \subseteq D_+(f)\) of basic opens corresponds,
  on affine models, to \(\Spec\) of the canonical map \(A_{(f)} \to A_{(fg)}\).
\end{proof}

\begin{theorem}[The gluing identity]
  \label{pa-thm:P1_gluing_identity}
  \lean{AlgebraicGeometry.P1.fromSpecChart_units}
  \leanok
  \dcref{custom}

  \uses{pa-def:P1_fromSpecChart, pa-lem:P1_SpecMap_overlap_chartIota,
    pa-lem:P1_overlap_isLocalization, pa-lem:P1_coord_mul, pa-lem:P1_polyToAway_surjective}
  Let \(\sigma : k \to A\) and let \(u \in A\) be a unit. Then the morphism
  \([1 : u] : \Spec A \to \PP^1\) through the chart \(U_0\) equals the morphism
  \([u^{-1} : 1]\) through the chart \(U_1\) (both as in \ref{pa-def:P1_fromSpecChart}).
\end{theorem}
\begin{proof}
  \leanok
  \uses{pa-def:P1_chartEval}
  Both morphisms factor through the chart overlap, where they are induced by one and the
  same ring homomorphism.

  The evaluation \(e_0 : A_{(X_0)} \to A\) at \(u\) (\ref{pa-def:P1_chartEval}) sends the
  coordinate \(t_{01}\) to the unit \(u\). Since \(A_{(X_0X_1)}\) is the localization of
  \(A_{(X_0)}\) at the powers of \(t_{01}\) (\ref{pa-lem:P1_overlap_isLocalization}), the
  universal property of localization extends \(e_0\) to a ring homomorphism
  \(\psi : A_{(X_0X_1)} \to A\) with \(\psi \circ \lambda = e_0\).

  From \(\lambda(t_{01}) \cdot \rho(t_{10}) = 1\) in \(A_{(X_0X_1)}\)
  (\ref{pa-lem:P1_coord_mul}) we get \(u \cdot \psi(\rho(t_{10})) = 1\), i.e.\
  \(\psi(\rho(t_{10})) = u^{-1}\). Moreover \(\psi \circ \rho\) agrees with the structure
  map on constants: the two restrictions \(\lambda\) and \(\rho\) restrict to the same map
  on the constants \(k \to A_{(X_0X_1)}\), and \(\psi \circ \lambda = e_0\) restricts to
  \(\sigma : k \to A\) on constants. Since \(A_{(X_1)}\) is generated as a \(k\)-algebra by
  \(t_{10}\) (\ref{pa-lem:P1_polyToAway_surjective}), it follows that \(\psi \circ \rho\) is
  exactly the evaluation \(e_1 : A_{(X_1)} \to A\) at \(u^{-1}\).

  Now apply \(\Spec\) and compose with the chart models. Writing \(\iota_0, \iota_1,
  \iota_{01}\) for the affine models of \(U_0\), \(U_1\) and \(D_+(X_0X_1)\), the
  factorization \(e_0 = \psi \circ \lambda\) and \ref{pa-lem:P1_SpecMap_overlap_chartIota}
  give
  \[
    [1 : u]
      = \iota_0 \circ \Spec e_0
      = \iota_0 \circ \Spec \lambda \circ \Spec \psi
      = \iota_{01} \circ \Spec \psi ,
  \]
  and the same computation for \(e_1 = \psi \circ \rho\) gives
  \([u^{-1} : 1] = \iota_1 \circ \Spec e_1 = \iota_{01} \circ \Spec \psi\). The two
  morphisms coincide.
\end{proof}

\begin{theorem}[Transcendental points land on the generic point]
  \label{pa-thm:P1_fromSpecChart_generic}
  \lean{AlgebraicGeometry.P1.fromSpecChart_base_genericPoint}
  \leanok
  \dcref{custom}

  \uses{pa-def:P1_fromSpecChart, pa-lem:P1_polyToAway_surjective, pa-lem:P1_generic_point_chart}
  Let \(A\) be an integral domain under \(\sigma : k \to A\) and let \(a \in A\) be
  transcendental in the sense that no nonzero polynomial over \(k\), evaluated at \(a\)
  with coefficients mapped by \(\sigma\), vanishes. Then the morphism
  \([1 : a] : \Spec A \to \PP^1\) of \ref{pa-def:P1_fromSpecChart} sends the generic point of
  \(\Spec A\) to the generic point of \(\PP^1\).
\end{theorem}
\begin{proof}
  \leanok
  \uses{pa-def:P1_chartEval}
  The evaluation \(A_{(X_0)} \to A\) at \(a\) is injective: every \(z \in A_{(X_0)}\) is a
  polynomial in the coordinate \(t_{01}\) with coefficients in \(k\)
  (\ref{pa-lem:P1_polyToAway_surjective}), its image is that polynomial evaluated at \(a\),
  and by hypothesis this vanishes only for the zero polynomial. The spectrum of an
  injective homomorphism between domains sends the zero ideal to the zero ideal, i.e.\ the
  generic point of \(\Spec A\) to the generic point of \(\Spec A_{(X_0)}\); conclude by
  \ref{pa-lem:P1_generic_point_chart}.
\end{proof}

\section{Spreading out a rational function to the projective line}

On an integral scheme \(X\), proper and smooth of relative dimension one over \(k\), a
rational function \(f_0 \in K(X)\) transcendental over \(k\) spreads out to a
\(k\)-morphism \(\pi : X \to \PP^1\) with finite fibers. This section provides the
stalk-level bricks and the spreading-out argument. Classically, the corresponding morphism
is produced from the inclusion of function fields \(k(f_0) \subseteq K(X)\) together with
the finiteness of \(K(X)\) over \(k(f_0)\); the route taken here never uses that
finiteness --- quasi-finiteness of \(\pi\) is obtained purely topologically from the
height-one specialization order of both spaces.

\begin{lemma}[Stalk morphisms are compatible with specialization]
  \label{pa-lem:fromSpecStalkOfMem_stalkSpecializes}
  \lean{AlgebraicGeometry.Scheme.Opens.fromSpecStalkOfMem_stalkSpecializes}
  \leanok
  \dcref{custom}

  Let \(X\) be a scheme, \(U \subseteq X\) open, \(x \in U\), and let \(y\) be a
  generalization of \(x\) (so \(y \in U\) as well). Write
  \(\struct{X,x} \to \struct{X,y}\) for the canonical map of stalks (restriction of
  germs). Then the composite
  \[
    \Spec \struct{X,y} \longrightarrow \Spec \struct{X,x} \longrightarrow U
  \]
  of the induced morphism of spectra with the canonical morphism
  \(\Spec \struct{X,x} \to U\) is the canonical morphism \(\Spec \struct{X,y} \to U\).
\end{lemma}
\begin{proof}
  \leanok
  The stalk of \(X\) at a point of \(U\) is canonically the stalk of the open subscheme
  \(U\) at that point, compatibly with the canonical maps of stalks along
  specializations (naturality of germs along the open immersion \(U \to X\)); and on
  \(X\) itself the analogous triangle of canonical morphisms
  \(\Spec \struct{X,y} \to \Spec \struct{X,x} \to X\) commutes, both composites being the
  morphism determined by the local ring \(\struct{X,y}\) and its ideal-theoretic image.
  Transport this triangle through the open immersion.
\end{proof}

\begin{definition}[Partial and rational maps]
  \label{pa-def:rational_map_machinery}
  \lean{AlgebraicGeometry.Scheme.RationalMap.ofFunctionField,
    AlgebraicGeometry.Scheme.RationalMap.toPartialMap,
    AlgebraicGeometry.Scheme.PartialMap.ofFromSpecStalk,
    AlgebraicGeometry.Scheme.RationalMap.mem_domain,
    AlgebraicGeometry.Scheme.RationalMap.eq_of_fromFunctionField_eq}
  \mathlibok
  \dcref{custom}

  Let \(X, Y\) be schemes over a base \(S\). A \emph{partial map} \(X \dashrightarrow Y\)
  is a morphism to \(Y\) from a dense open subscheme of \(X\) (its \emph{domain}); a
  \emph{rational map} is an equivalence class of partial maps, two being equivalent when
  they agree on a common dense open. For \(X\) irreducible with generic point \(\xi\) and
  function field \(K(X) = \struct{X,\xi}\), a partial map \(g\) restricts to a morphism
  \(\Spec K(X) \to Y\) (its \emph{function-field morphism}), which determines the rational
  map. For \(X\) integral and \(Y\) locally of finite type over \(S\):
  \begin{enumerate}
    \item every \(S\)-morphism \(\Spec K(X) \to Y\) arises as the function-field morphism
      of a unique rational \(S\)-map \(X \dashrightarrow Y\) (spreading out);
    \item if moreover an \(S\)-morphism \(\Spec \struct{X,x} \to Y\) restricting to the
      given one on \(\Spec K(X)\) exists at a point \(x\), then \(x\) lies in the domain
      of definition of that rational map (the union of the domains of its
      representatives);
    \item if \(Y\) is separated, a rational map has a largest representative, defined on
      its whole domain of definition.
  \end{enumerate}
\end{definition}

\begin{lemma}[Restriction of a partial map to the function field]
  \label{pa-lem:partialMap_fromFunctionField_eq}
  \lean{AlgebraicGeometry.Scheme.PartialMap.fromFunctionField_eq}
  \leanok
  \dcref{custom}

  \uses{pa-def:rational_map_machinery, pa-lem:fromSpecStalkOfMem_stalkSpecializes}
  Let \(X\) be a scheme with irreducible underlying space, \(g : X \dashrightarrow Y\) a
  partial map and \(x\) a point of its domain. Then the function-field morphism of \(g\)
  is the composite
  \[
    \Spec K(X) \longrightarrow \Spec \struct{X,x}
      \longrightarrow Y,
  \]
  where the first arrow is induced by the canonical map
  \(\struct{X,x} \to K(X) = \struct{X,\xi}\) and the second is the restriction of \(g\) to
  \(\Spec \struct{X,x}\).
\end{lemma}
\begin{proof}
  \leanok
  The function-field morphism is by definition the restriction of \(g\) to
  \(\Spec\) of the stalk at the generic point, which is a generalization of \(x\); apply
  \ref{pa-lem:fromSpecStalkOfMem_stalkSpecializes} inside the domain of \(g\).
\end{proof}

\begin{definition}[The structure map of a stalk]
  \label{pa-def:structureStalk}
  \lean{AlgebraicGeometry.Scheme.Hom.structureStalk}
  \leanok
  \dcref{custom}

  Let \(f : X \to \Spec k\) be a scheme over \(k\) and \(x \in X\). The structure map of
  the stalk at \(x\) is the composite ring homomorphism
  \(k \to \Gamma(X, X) \to \struct{X,x}\) of the map induced by \(f\) on global sections
  with the germ at \(x\); it makes every stalk of \(X\) a ring under \(k\).
\end{definition}

\begin{lemma}[Coherence of the structure maps of stalks]
  \label{pa-lem:structureStalk_coherence}
  \lean{AlgebraicGeometry.Scheme.Hom.Spec_map_structureStalk,
    AlgebraicGeometry.Scheme.Hom.structureStalk_stalkSpecializes}
  \leanok
  \dcref{custom}

  \uses{pa-def:structureStalk}
  Let \(f : X \to \Spec k\) and \(x \in X\).
  \begin{enumerate}
    \item Applying \(\Spec\) to the structure map of \ref{pa-def:structureStalk} yields the
      composite \(\Spec \struct{X,x} \to X \xrightarrow{f} \Spec k\) of the canonical
      morphism with \(f\).
    \item For a generalization \(y\) of \(x\), composing the structure map at \(x\) with
      the canonical map \(\struct{X,x} \to \struct{X,y}\) yields the structure map at
      \(y\).
  \end{enumerate}
\end{lemma}
\begin{proof}
  \leanok
  (1) is the compatibility of the canonical morphism \(\Spec \struct{X,x} \to X\) with
  global sections: its composite with \(f\) corresponds, under the anti-equivalence of
  affine schemes and rings, to the composite \(k \to \Gamma(X, X) \to \struct{X,x}\).
  (2) is the compatibility of germs with the canonical maps of stalks: the germ at \(x\)
  followed by restriction of germs to \(y\) is the germ at \(y\).
\end{proof}

\begin{theorem}[A smooth projective curve maps locally quasi-finitely and dominantly to the line]
  \label{pa-thm:exists_lqf_dominant_toP1}
  \lean{AlgebraicGeometry.exists_locallyQuasiFinite_isDominant_toP1}
  \leanok
  \dcref{custom}

  \uses{pa-def:P1_structureMap}
  Let \(X\) be an integral scheme, proper and smooth of relative dimension one over the
  field \(k\). Then there is a morphism \(\pi : X \to \PP^1\) over \(k\) which is locally
  quasi-finite and dominant.
\end{theorem}
\begin{proof}
  \leanok
  \uses{pa-thm:relDim1_transcendental, pa-thm:relDim1_stalks_valuation, pa-thm:relDim1_specializes,
    pa-def:structureStalk, pa-lem:structureStalk_coherence, pa-def:P1_fromSpecChart,
    pa-lem:P1_fromSpecChart_over, pa-lem:P1_fromSpecChart_natural, pa-thm:P1_gluing_identity,
    pa-thm:P1_fromSpecChart_generic, pa-def:rational_map_machinery,
    pa-lem:partialMap_fromFunctionField_eq, pa-lem:height_one_closed_points,
    pa-thm:height_one_finite_closed, pa-lem:P1_exists_ne_generic, pa-lem:P1_closed_points,
    pa-thm:P1_separated, pa-thm:P1_proper}
  Write \(\xi\) for the generic point of \(X\) and \(K(X) = \struct{X,\xi}\) for the
  function field, a field under \(k\) via the structure map of \ref{pa-def:structureStalk}.

  (a) \emph{The generic morphism.} By \ref{pa-thm:relDim1_transcendental} choose
  \(f_0 \in K(X)\) transcendental over \(k\); in particular \(f_0 \neq 0\) (the polynomial
  \(T\) does not vanish at \(f_0\)). Let
  \(\varphi_K = [1 : f_0] : \Spec K(X) \to \PP^1\) be the chart-\(0\) morphism of
  \ref{pa-def:P1_fromSpecChart}. It lies over \(k\), compatibly with the canonical morphism
  \(\Spec K(X) \to X\), by \ref{pa-lem:P1_fromSpecChart_over} and
  \ref{pa-lem:structureStalk_coherence}(1).

  (b) \emph{Local regularity.} Let \(x \in X\). The stalk \(\struct{X,x}\) is a valuation
  ring (\ref{pa-thm:relDim1_stalks_valuation}) with fraction field \(K(X)\), so \(f_0\) or
  \(f_0^{-1}\) lies in \(\struct{X,x}\). In the first case, with \(f_0 = a\) in
  \(\struct{X,x}\), set \(\varphi_x = [1 : a] : \Spec \struct{X,x} \to \PP^1\) through
  chart \(0\); in the second, with \(f_0^{-1} = b\), set \(\varphi_x = [b : 1]\) through
  chart \(1\). In both cases \(\varphi_x\) lies over \(k\)
  (\ref{pa-lem:P1_fromSpecChart_over}, \ref{pa-lem:structureStalk_coherence}(1)), and its
  restriction along \(\Spec K(X) \to \Spec \struct{X,x}\) is \(\varphi_K\): in the first
  case by naturality (\ref{pa-lem:P1_fromSpecChart_natural}) together with
  \ref{pa-lem:structureStalk_coherence}(2), since \(a \mapsto f_0\) under
  \(\struct{X,x} \to K(X)\); in the second case, naturality gives the chart-\(1\)
  morphism \([f_0^{-1} : 1]\) on \(\Spec K(X)\), which equals \([1 : f_0]\) by the gluing
  identity \ref{pa-thm:P1_gluing_identity} applied to the unit \(f_0\) of the field
  \(K(X)\).

  (c) \emph{Spreading out.} The projective line is locally of finite type and separated
  over \(k\) (\ref{pa-thm:P1_proper}). By \ref{pa-def:rational_map_machinery}(1), \(\varphi_K\)
  is the function-field morphism of a rational \(k\)-map
  \(F : X \dashrightarrow \PP^1\). Every point \(x\) lies in the domain of definition of
  \(F\): the morphism \(\varphi_x\) of (b) spreads out to a partial map defined at \(x\)
  (\ref{pa-def:rational_map_machinery}(2)) whose function-field morphism is the restriction
  of \(\varphi_x\) to \(\Spec K(X)\) (\ref{pa-lem:partialMap_fromFunctionField_eq}), i.e.\
  \(\varphi_K\) --- so this partial map represents \(F\). Since \(X\) is reduced and
  \(\PP^1\) separated (\ref{pa-thm:P1_separated}), \(F\) has a largest representative
  (\ref{pa-def:rational_map_machinery}(3)), defined on the whole of \(X\): an honest
  morphism \(\pi : X \to \PP^1\) whose restriction to \(\Spec K(X)\) is \(\varphi_K\).

  (d) \emph{\(\pi\) lies over \(k\).} The canonical morphism \(\Spec K(X) \to X\) hits the
  generic point, hence is dominant, and \(\PP^1\) is separated; two morphisms out of the
  reduced scheme \(X\) to a separated scheme agreeing after composition with a dominant
  morphism agree. Both \(\pi\) followed by the structure morphism of \(\PP^1\) and the
  structure morphism \(f\) of \(X\) restrict on \(\Spec K(X)\) to the structure morphism
  of \(\Spec K(X)\) (step (a)), so they agree.

  (e) \emph{\(\pi\) is proper and locally of finite type.} By (d), \(\pi\) followed by the
  separated structure morphism of \(\PP^1\) is the proper, locally of finite type
  morphism \(f\); properness and local finite-typeness descend along such factorizations
  (cancellation).

  (f) \emph{\(\pi\) maps \(\xi\) to the generic point of \(\PP^1\).} The point
  \(\pi(\xi)\) is the image of the unique point of \(\Spec K(X)\) under \(\varphi_K\),
  and \(\varphi_K = [1 : f_0]\) sends it to the generic point of \(\PP^1\) by
  \ref{pa-thm:P1_fromSpecChart_generic} (the hypothesis is the transcendence of \(f_0\), and
  the unique point of \(\Spec K(X)\) is its generic point).

  (g) \emph{All fibers of \(\pi\) are finite.} The underlying space of \(X\) is Noetherian:
  \(X\) is locally of finite type over a field, hence locally Noetherian, and
  quasi-compact (proper). Its specialization order has height one
  (\ref{pa-thm:relDim1_specializes}).

  \emph{The fiber of the generic point of \(\PP^1\) is contained in \(\{\xi\}\):} suppose
  \(x \neq \xi\) with \(\pi(x)\) generic. Then \(\{x\}\) is closed
  (\ref{pa-lem:height_one_closed_points}), and \(\pi\), being proper, is a closed map, so
  \(\{\pi(x)\}\) is closed in \(\PP^1\). But the closure of the generic point is all of
  \(\PP^1\), which has a second point (\ref{pa-lem:P1_exists_ne_generic}) --- a
  contradiction.

  \emph{Any other fiber is finite:} for \(y\) a non-generic point of \(\PP^1\), the
  singleton \(\{y\}\) is closed (\ref{pa-lem:P1_closed_points}), so the fiber
  \(\pi^{-1}(y)\) is closed in \(X\), and it avoids \(\xi\) because
  \(\pi(\xi)\) is the generic point of \(\PP^1\) (step (f)), which differs from \(y\).
  By \ref{pa-thm:height_one_finite_closed} the fiber is finite.

  (h) A morphism which is locally of finite type and has finite fibers is locally
  quasi-finite; by (e) and (g), \(\pi\) is locally quasi-finite.

  (i) \emph{\(\pi\) is dominant.} By (f) the image of \(\pi\) contains the generic point
  of \(\PP^1\), whose closure is all of \(\PP^1\); so the image of \(\pi\) is dense, i.e.\
  \(\pi\) is dominant.
\end{proof}

\begin{theorem}[A smooth projective curve maps locally quasi-finitely to the line]
  \label{pa-thm:exists_lqf_toP1}
  \lean{AlgebraicGeometry.exists_locallyQuasiFinite_toP1}
  \leanok
  \dcref{custom}

  \uses{pa-def:P1_structureMap}
  Let \(X\) be an integral scheme, proper and smooth of relative dimension one over the
  field \(k\). Then there is a morphism \(\pi : X \to \PP^1\) over \(k\) which is locally
  quasi-finite.
\end{theorem}
\begin{proof}
  \leanok
  \uses{pa-thm:exists_lqf_dominant_toP1}
  The morphism of \ref{pa-thm:exists_lqf_dominant_toP1}, forgetting that it is dominant.
\end{proof}

\section{A finite morphism to the projective line}

\begin{theorem}[Zariski's main theorem, cancellation form]
  \label{pa-thm:zmt_cancellation}
  \lean{AlgebraicGeometry.IsFinite.of_locallyQuasiFinite_of_comp}
  \leanok
  \dcref{custom}

  Let \(\pi : X \to Y\) and \(g : Y \to S\) be morphisms of schemes with \(\pi\) locally
  quasi-finite, \(g\) separated and \(g \circ \pi\) proper. Then \(\pi\) is finite.
\end{theorem}
\begin{proof}
  \leanok
  By the cancellation property of proper morphisms (using that \(g\) is separated), \(\pi\)
  is proper. A proper, locally quasi-finite morphism is finite (Zariski's main theorem).
\end{proof}

\begin{lemma}[Locally quasi-finite maps to the projective line are finite]
  \label{pa-lem:finite_toP1_of_lqf}
  \lean{AlgebraicGeometry.isFinite_toP1_of_locallyQuasiFinite}
  \leanok
  \dcref{custom}

  \uses{pa-def:P1_structureMap}
  Let \(X\) be a scheme, proper over \(k\), and let \(\pi : X \to \PP^1\) be a locally
  quasi-finite morphism over \(k\) (i.e.\ its composite with the structure morphism of
  \(\PP^1\) is the structure morphism of \(X\)). Then \(\pi\) is finite.
\end{lemma}
\begin{proof}
  \leanok
  \uses{pa-thm:zmt_cancellation, pa-thm:P1_proper}
  The structure morphism of \(\PP^1\) is proper (\ref{pa-thm:P1_proper}), in particular
  separated, and the composite of \(\pi\) with it is the structure morphism of \(X\), which
  is proper. Apply \ref{pa-thm:zmt_cancellation}.
\end{proof}

\begin{lemma}[Reduction to local quasi-finiteness]
  \label{pa-lem:exists_finite_of_lqf}
  \lean{AlgebraicGeometry.exists_isFinite_toP1_of_locallyQuasiFinite}
  \leanok
  \dcref{custom}

  \uses{pa-def:P1_structureMap}
  Let \(C\) be a scheme, proper over \(k\). If there exists a locally quasi-finite morphism
  \(C \to \PP^1\) over \(k\), then there exists a finite morphism \(C \to \PP^1\) over \(k\).
\end{lemma}
\begin{proof}
  \leanok
  \uses{pa-lem:finite_toP1_of_lqf}
  Immediate from \ref{pa-lem:finite_toP1_of_lqf}: the given locally quasi-finite morphism is
  itself finite.
\end{proof}

\begin{theorem}[Existence of a finite morphism to the projective line]
  \label{pa-thm:exists_finite_toP1}
  \lean{AlgebraicGeometry.exists_isFinite_toP1}
  \leanok
  \dcref{custom}

  \uses{pa-def:P1, pa-def:P1_structureMap}
  Let \(C\) be a smooth curve over \(k\) (proper, smooth of relative dimension one,
  geometrically irreducible). Then there exists a finite morphism \(\pi : C \to \PP^1\) over
  \(k\).
\end{theorem}
\begin{proof}
  \leanok
  \uses{pa-lem:relDim_geometrically_reduced, pa-prop:curve_integral, pa-thm:exists_lqf_toP1,
    pa-lem:exists_finite_of_lqf}
  \(C\) is geometrically reduced over \(k\) (\ref{pa-lem:relDim_geometrically_reduced}), and
  being also geometrically irreducible, it is an integral scheme
  (\ref{pa-prop:curve_integral}). By \ref{pa-thm:exists_lqf_toP1} there is a locally quasi-finite
  morphism \(C \to \PP^1\) over \(k\), and by \ref{pa-lem:exists_finite_of_lqf} it may be
  taken finite.
\end{proof}

\begin{theorem}[Existence of a finite dominant morphism to the projective line]
  \label{pa-thm:exists_finite_dominant_toP1}
  \lean{AlgebraicGeometry.exists_isFinite_isDominant_toP1}
  \leanok
  \dcref{custom}

  \uses{pa-def:P1, pa-def:P1_structureMap}
  Let \(C\) be a smooth curve over \(k\) (proper, smooth of relative dimension one,
  geometrically irreducible). Then there exists a finite morphism \(\pi : C \to \PP^1\)
  over \(k\) which is moreover \emph{dominant}. The dominance is part of the output of the
  construction --- the constructed \(\pi\) sends the generic point of \(C\) to the generic
  point of \(\PP^1\) --- and is not asserted as a formal consequence of finiteness alone.
\end{theorem}
\begin{proof}
  \leanok
  \uses{pa-lem:relDim_geometrically_reduced, pa-prop:curve_integral, pa-thm:exists_lqf_dominant_toP1,
    pa-lem:finite_toP1_of_lqf}
  \(C\) is geometrically reduced over \(k\) (\ref{pa-lem:relDim_geometrically_reduced}), and
  being also geometrically irreducible, it is an integral scheme
  (\ref{pa-prop:curve_integral}). By \ref{pa-thm:exists_lqf_dominant_toP1} there is a locally
  quasi-finite \emph{dominant} morphism \(\pi : C \to \PP^1\) over \(k\), and by
  \ref{pa-lem:finite_toP1_of_lqf} this same morphism \(\pi\) is finite. Since the finiteness
  upgrade concerns the very morphism already constructed, its dominance and its
  compatibility with the structure morphisms persist unchanged.
\end{proof}

\begin{lemma}[Preimages of the charts are affine]
  \label{pa-lem:preimage_chart_affine}
  \lean{AlgebraicGeometry.isAffineOpen_preimage_chartOpen}
  \leanok
  \dcref{custom}

  \uses{pa-def:P1_chartOpen}
  Let \(\pi : Y \to \PP^1\) be an affine morphism (for instance a finite one). Then
  \(\pi^{-1}(U_i)\) is an affine open of \(Y\) for \(i \in \{0,1\}\).
\end{lemma}
\begin{proof}
  \leanok
  \uses{pa-lem:P1_chartOpen_affine}
  The charts are affine (\ref{pa-lem:P1_chartOpen_affine}), and the preimage of an affine open
  under an affine morphism is affine.
\end{proof}

\begin{lemma}[Preimages of the charts cover]
  \label{pa-lem:preimage_chart_sup}
  \lean{AlgebraicGeometry.preimage_chartOpen_sup}
  \leanok
  \dcref{custom}

  \uses{pa-def:P1_chartOpen}
  For any morphism \(\pi : Y \to \PP^1\), the preimages of the two standard charts cover
  \(Y\): \(\pi^{-1}(U_0) \cup \pi^{-1}(U_1) = Y\).
\end{lemma}
\begin{proof}
  \leanok
  \uses{pa-lem:P1_chartOpen_sup}
  Preimages of a cover form a cover (\ref{pa-lem:P1_chartOpen_sup}).
\end{proof}

\begin{lemma}[Module-finiteness over the charts]
  \label{pa-lem:finite_app_chart}
  \lean{AlgebraicGeometry.finite_app_chartOpen}
  \leanok
  \dcref{custom}

  \uses{pa-def:P1_chartOpen}
  Let \(\pi : Y \to \PP^1\) be a finite morphism. For \(i \in \{0,1\}\), the induced ring
  homomorphism \(\Gamma(\PP^1, U_i) \to \Gamma(Y, \pi^{-1}(U_i))\) is module-finite. (By
  \ref{pa-def:P1_chartSectionsEquiv0}--\ref{pa-def:P1_chartSectionsEquiv1}, the source is the
  polynomial ring \(k[t]\).)
\end{lemma}
\begin{proof}
  \leanok
  \uses{pa-lem:P1_chartOpen_affine}
  A finite morphism induces module-finite maps on sections over affine opens; the charts are
  affine by \ref{pa-lem:P1_chartOpen_affine}.
\end{proof}

\begin{lemma}[Module-finiteness over the overlap]
  \label{pa-lem:finite_app_overlap}
  \lean{AlgebraicGeometry.finite_app_overlap}
  \leanok
  \dcref{custom}

  \uses{pa-lem:P1_chartOpen_inf}
  Let \(\pi : Y \to \PP^1\) be a finite morphism. The induced ring homomorphism
  \(\Gamma(\PP^1, U_0 \cap U_1) \to \Gamma(Y, \pi^{-1}(U_0 \cap U_1))\) is module-finite.
  (By \ref{pa-def:P1_overlapSectionsEquiv}, the source is the Laurent polynomial ring
  \(k[T,T^{-1}]\).)
\end{lemma}
\begin{proof}
  \leanok
  \uses{pa-lem:P1_overlap_affine}
  As in \ref{pa-lem:finite_app_chart}, with the overlap affine by \ref{pa-lem:P1_overlap_affine}.
\end{proof}

\section{An affine two-chart cover of the curve, and its stability under base change}

The two standard charts of \(\PP^1\) pull back, along a finite morphism \(C \to \PP^1\), to a
cover of the curve \(C\) by two affine opens with affine overlap. This section records that
cover as a structure in its own right, together with its stability under base change of the
test algebra.

\begin{definition}[Affine two-chart cover]
  \label{pa-def:AffineTwoCover}
  \lean{AlgebraicGeometry.Scheme.AffineTwoCover}
  \leanok
  \dcref{custom}

  Let \(Y\) be a scheme. An \emph{affine two-chart cover} of \(Y\) is a pair of affine open
  subschemes \(V_0, V_1 \subseteq Y\) with \(V_0 \cup V_1 = Y\), whose overlap \(V_0 \cap V_1\)
  is again affine.
\end{definition}

\begin{theorem}[Existence of an affine two-chart cover of the curve]
  \label{pa-thm:AffineTwoCover_nonempty_of_curve}
  \lean{AlgebraicGeometry.Scheme.AffineTwoCover.nonempty_of_curve}
  \leanok
  \dcref{custom}

  \uses{pa-def:AffineTwoCover, pa-thm:exists_finite_toP1, pa-lem:preimage_chart_affine,
    pa-lem:preimage_chart_sup, pa-lem:P1_chartOpen_inf, pa-lem:P1_overlap_affine}
  Let \(C\) be a smooth, proper, geometrically irreducible curve over a field \(k\). Then \(C\)
  admits an affine two-chart cover.
\end{theorem}
\begin{proof}
  \leanok
  By Theorem~\ref{pa-thm:exists_finite_toP1} there is a finite morphism \(\pi : C \to \PP^1\) over
  \(k\). Take \(V_0 := \pi^{-1}(U_0)\) and \(V_1 := \pi^{-1}(U_1)\), the preimages of the two
  standard charts of \(\PP^1\) (Definition~\ref{pa-def:P1_chartOpen}); each is affine, since \(\pi\)
  is finite, hence affine, and preimages of affine opens under affine morphisms are affine
  (Lemma~\ref{pa-lem:preimage_chart_affine}). The preimages cover \(C\)
  (Lemma~\ref{pa-lem:preimage_chart_sup}). For the overlap, preimage commutes with intersection, so
  \(V_0 \cap V_1 = \pi^{-1}(U_0 \cap U_1)\); since \(U_0 \cap U_1 = D_+(X_0X_1)\)
  (Lemma~\ref{pa-lem:P1_chartOpen_inf}) is affine (Lemma~\ref{pa-lem:P1_overlap_affine}), its preimage
  \(V_0 \cap V_1\) is affine as well, again because preimages of affine opens under the affine
  morphism \(\pi\) are affine. Hence \((V_0, V_1)\) is an affine two-chart cover of \(C\).
\end{proof}

\begin{theorem}[Base change of an affine two-chart cover]
  \label{pa-thm:AffineTwoCover_pullbackProd}
  \lean{AlgebraicGeometry.Scheme.AffineTwoCover.pullbackProd}
  \leanok
  \dcref{custom}

  \uses{pa-def:AffineTwoCover}
  Let \(k\) be a field, \(C\) an object of \(\Over(\Spec k)\), and \(D = (V_0,V_1)\) an affine
  two-chart cover of the underlying scheme of \(C\). Let \(R\) be a commutative \(k\)-algebra,
  and write \(C \otimes \Spec R\) for the fibre product of \(C\) with the object \(\Spec R\) of
  \(\Over(\Spec k)\) (the \(k\)-algebra structure of \(R\) making \(\Spec R\) an object of
  \(\Over(\Spec k)\)), with first projection \(\mathrm{pr}_1 : C \otimes \Spec R \to C\). Then
  \(\bigl(\mathrm{pr}_1^{-1}(V_0),\ \mathrm{pr}_1^{-1}(V_1)\bigr)\) is an affine two-chart cover
  of \(C \otimes \Spec R\).
\end{theorem}
\begin{proof}
  \leanok
  The projection \(\mathrm{pr}_1\) is an affine morphism: it is the base change, along the
  structure morphism of \(C\), of the structure morphism \(\Spec R \to \Spec k\) (a morphism
  between affine schemes, hence affine), and affineness of morphisms is stable under base
  change. Preimages of affine opens under an affine morphism are affine, so
  \(\mathrm{pr}_1^{-1}(V_0)\) and \(\mathrm{pr}_1^{-1}(V_1)\) are affine. Preimage commutes with
  union and with intersection, so these two preimages cover \(C \otimes \Spec R\) (since
  \(V_0 \cup V_1 = C\)) and their overlap is \(\mathrm{pr}_1^{-1}(V_0 \cap V_1)\), the preimage
  of the affine overlap of \(D\), hence itself affine. Thus the two preimages form an affine
  two-chart cover of \(C \otimes \Spec R\).
\end{proof}

\section{Product charts on a fibre product}

Let \(X\) and \(T\) be objects of \(\Over(\Spec k)\), with the two projections
\(\mathrm{fst} \colon X \otimes T \to X\) and \(\mathrm{snd} \colon X \otimes T \to T\)
exhibiting the underlying scheme of the fibre product \(X \otimes T\) as
\(X \times_{\Spec k} T\). For opens \(U\) of \(X\) and \(V\) of \(T\) this section builds the
open \(\mathfrak W(U,V) \subseteq X \otimes T\) cut out by the two projections and, for \(U, V\)
affine, identifies its section ring with the tensor product \(\Gamma(X,U) \otimes_k \Gamma(T,V)\).
All \(k\)-algebra structures on section rings are those obtained by pulling back along the
structure morphism to \(\Spec k\).

\begin{definition}[The product chart]
  \label{pa-def:productChart}
  \lean{AlgebraicGeometry.Over.productChart, AlgebraicGeometry.Over.isAffineOpen_productChart}
  \leanok
  \dcref{custom}

  For opens \(U\) of \(X\) and \(V\) of \(T\), the \emph{product chart} is the open
  \[
    \mathfrak W(U,V) \;:=\; \mathrm{fst}^{-1}(U) \cap \mathrm{snd}^{-1}(V) \;\subseteq\; X \otimes T .
  \]
  It is monotone in \(U\) and \(V\). If \(U\) and \(V\) are affine, then \(\mathfrak W(U,V)\) is
  affine.
\end{definition}
\begin{proof}
  \leanok
  Monotonicity follows from monotonicity of preimages and of intersections. For affineness, the
  two projections exhibit \(X \otimes T\) as the fibre product \(X \times_{\Spec k} T\).
  Restricting this cartesian square to the opens \(U \subseteq X\), \(V \subseteq T\) and the
  whole base \(\Spec k\) presents the restriction of \(X \otimes T\) to
  \(\mathrm{fst}^{-1}(U) \cap \mathrm{snd}^{-1}(V) = \mathfrak W(U,V)\) as the fibre product
  \(U \times_{\Spec k} V\) of the restricted schemes. When \(U\), \(V\) and the base \(\Spec k\)
  are affine, this is a fibre product of affine schemes over an affine scheme, hence affine; so
  \(\mathfrak W(U,V)\) is an affine open.
\end{proof}

\begin{definition}[The sections identification]
  \label{pa-def:productChartSections}
  \lean{AlgebraicGeometry.Over.isPushout_productChartSections,
    AlgebraicGeometry.Over.productChartSections,
    AlgebraicGeometry.Over.productChartSectionsAlgEquiv,
    AlgebraicGeometry.Over.productChartSections_tmul,
    AlgebraicGeometry.Over.productChartSections_naturality}
  \leanok
  \dcref{custom}

  \uses{pa-def:productChart}
  Let \(U\) and \(V\) be affine. The square of section rings
  \[
    \begin{array}{ccc}
      \Gamma(\Spec k, \struct{}) & \longrightarrow & \Gamma(X, U) \\
      \downarrow & & \downarrow \\
      \Gamma(T, V) & \longrightarrow & \Gamma\bigl(X \otimes T,\ \mathfrak W(U,V)\bigr)
    \end{array}
  \]
  --- the top and left edges the pullbacks along the two structure morphisms to \(\Spec k\), the
  bottom and right edges the two projection pullbacks \(\mathrm{fst}^\sharp\),
  \(\mathrm{snd}^\sharp\) --- is a pushout of commutative rings. Since
  \(\Gamma(\Spec k, \struct{}) = k\), it follows that there is a \(k\)-algebra isomorphism
  \[
    \Gamma(X, U) \otimes_k \Gamma(T, V) \;\xrightarrow{\ \sim\ }\;
      \Gamma\bigl(X \otimes T,\ \mathfrak W(U,V)\bigr), \qquad
      s \otimes t \;\longmapsto\; \mathrm{fst}^\sharp(s) \cdot \mathrm{snd}^\sharp(t),
  \]
  natural in restriction in both factors: for affine opens \(U' \subseteq U\) and
  \(V' \subseteq V\), restricting an identified section to \(\mathfrak W(U', V')\) equals
  identifying the section obtained by restricting the two tensor factors to \(U'\) and \(V'\).
\end{definition}
\begin{proof}
  \leanok
  The square is a pushout by the base-change theorem for the sections of a fibre product, in the
  affine case: for a cartesian square of schemes and compatible opens over which the base and the
  two source opens are affine, the induced square of section rings is a pushout; here the base
  \(\Spec k\) and the opens \(U\), \(V\) are affine, and no flatness hypothesis is needed. On the
  other hand the tensor product \(\Gamma(X,U) \otimes_k \Gamma(T,V)\) is the pushout of the span
  \(k \to \Gamma(X,U)\), \(k \to \Gamma(T,V)\) of \(k\)-algebra structure maps. Comparing the two
  pushouts of this common span yields the displayed \(k\)-algebra isomorphism; it sends the two
  coprojections \(s \otimes 1\) and \(1 \otimes t\) to \(\mathrm{fst}^\sharp(s)\) and
  \(\mathrm{snd}^\sharp(t)\), hence a pure tensor \(s \otimes t = (s \otimes 1)(1 \otimes t)\) to
  the product \(\mathrm{fst}^\sharp(s) \cdot \mathrm{snd}^\sharp(t)\). Naturality is checked on
  pure tensors: both sides send \(s \otimes t\) to the restriction of
  \(\mathrm{fst}^\sharp(s) \cdot \mathrm{snd}^\sharp(t)\) to \(\mathfrak W(U',V')\), because
  projection pullback commutes with restriction of sections.
\end{proof}

\begin{theorem}[The lift-app rule]
  \label{pa-thm:appLE_productChartSections}
  \lean{AlgebraicGeometry.Over.appLE_productChartSections_tmul,
    AlgebraicGeometry.Over.appLE_productChartSections_tmul_lift,
    AlgebraicGeometry.Over.appLE_productChartSections_sub}
  \leanok
  \dcref{custom}

  \uses{pa-def:productChartSections}
  Let \(U\), \(V\) be affine and let \(q \colon Z \to X \otimes T\) be a morphism of
  \(\Over(\Spec k)\) with components \(a := \mathrm{fst} \circ q \colon Z \to X\) and
  \(b := \mathrm{snd} \circ q \colon Z \to T\). Then on any open \(W\) of \(Z\) contained in
  \(q^{-1}\mathfrak W(U,V)\), in \(a^{-1}(U)\) and in \(b^{-1}(V)\), the pullback \(q^\sharp\)
  sends the identified pure tensor and its difference shape to
  \[
    q^\sharp\bigl(s \otimes t\bigr) = a^\sharp(s) \cdot b^\sharp(t), \qquad
    q^\sharp\bigl(s \otimes 1 - 1 \otimes t\bigr) = a^\sharp(s) - b^\sharp(t).
  \]
  In particular, for the morphism \(\langle a, b\rangle \colon Z \to X \otimes T\) into the
  product determined by two maps \(a \colon Z \to X\), \(b \colon Z \to T\), the pullback
  \(\langle a, b\rangle^\sharp\) sends \(s \otimes t\) to \(a^\sharp(s) \cdot b^\sharp(t)\).
\end{theorem}
\begin{proof}
  \leanok
  The identification sends \(s \otimes t\) to
  \(\mathrm{fst}^\sharp(s) \cdot \mathrm{snd}^\sharp(t)\)
  (Definition~\ref{pa-def:productChartSections}). Pulling back along \(q\) and using functoriality
  of section pullback, \(q^\sharp \circ \mathrm{fst}^\sharp = (\mathrm{fst} \circ q)^\sharp = a^\sharp\)
  and \(q^\sharp \circ \mathrm{snd}^\sharp = b^\sharp\) over the opens in question, so
  \(q^\sharp(s \otimes t) = a^\sharp(s) \cdot b^\sharp(t)\). The difference formula follows by
  additivity together with \(s \otimes 1 \mapsto a^\sharp(s)\), \(1 \otimes t \mapsto b^\sharp(t)\)
  and \(a^\sharp(1) = b^\sharp(1) = 1\). For \(\langle a,b\rangle\) one has
  \(\mathrm{fst} \circ \langle a,b\rangle = a\) and \(\mathrm{snd} \circ \langle a,b\rangle = b\),
  so the general rule applies with \(q = \langle a, b\rangle\).
\end{proof}

\begin{lemma}[Basic open of a pure tensor]
  \label{pa-lem:basicOpen_productChartSections_tmul}
  \lean{AlgebraicGeometry.Over.basicOpen_productChartSections_tmul}
  \leanok
  \dcref{custom}

  \uses{pa-def:productChart, pa-def:productChartSections}
  Let \(U\), \(V\) be affine, \(s \in \Gamma(X,U)\) and \(t \in \Gamma(T,V)\). The basic open of
  the identified pure tensor is the product chart of the two factor basic opens:
  \[
    D\bigl(\mathrm{fst}^\sharp(s) \cdot \mathrm{snd}^\sharp(t)\bigr)
      \;=\; \mathfrak W\bigl(D(s),\, D(t)\bigr).
  \]
\end{lemma}
\begin{proof}
  \leanok
  The basic open of a product is the intersection of the two basic opens, and the basic open of a
  pullback \(\mathrm{fst}^\sharp(s)\) is the preimage \(\mathrm{fst}^{-1}(D(s))\) of the basic open
  of \(s\) (likewise for \(\mathrm{snd}^\sharp(t)\)). Hence
  \(D(\mathrm{fst}^\sharp(s) \cdot \mathrm{snd}^\sharp(t)) = \mathrm{fst}^{-1}(D(s)) \cap \mathrm{snd}^{-1}(D(t))
  = \mathfrak W(D(s), D(t))\), using \(D(s) \subseteq U\) and \(D(t) \subseteq V\).
\end{proof}

\begin{theorem}[Sections over a basic open localize the tensor ring]
  \label{pa-thm:isLocalization_away_basicOpen_productChartSections}
  \lean{AlgebraicGeometry.Over.isLocalization_away_basicOpen_productChartSections,
    AlgebraicGeometry.Over.isAffineOpen_basicOpen_productChartSections}
  \leanok
  \dcref{custom}

  \uses{pa-def:productChart, pa-def:productChartSections}
  Let \(U\), \(V\) be affine and \(x \in \Gamma(X,U) \otimes_k \Gamma(T,V)\); write \(\xi\) for
  the section of \(\mathfrak W(U,V)\) identified with \(x\). Then the basic open
  \(D(\xi) \subseteq \mathfrak W(U,V)\) is affine, and its section ring, regarded as a
  \(\Gamma(X,U) \otimes_k \Gamma(T,V)\)-algebra through the identification followed by restriction
  to \(D(\xi)\), is the localization of \(\Gamma(X,U) \otimes_k \Gamma(T,V)\) away from \(x\).
\end{theorem}
\begin{proof}
  \leanok
  The chart \(\mathfrak W(U,V)\) is affine (Definition~\ref{pa-def:productChart}), so a basic open of
  it is affine, and the section ring over \(D(\xi)\) is the localization of
  \(\Gamma(X \otimes T, \mathfrak W(U,V))\) away from \(\xi\). Transporting along the ring
  isomorphism of Definition~\ref{pa-def:productChartSections}, which carries \(x\) to \(\xi\) and the
  powers of \(x\) to the powers of \(\xi\), identifies this with the localization of
  \(\Gamma(X,U) \otimes_k \Gamma(T,V)\) away from \(x\).
\end{proof}

\section{The diagonal of the curve square}

Fix an object \(C\) of \(\Over(\Spec k)\) and form the fibre-product square \(C \otimes C\) with
its two projections \(\mathrm{fst}, \mathrm{snd} \colon C \otimes C \to C\). The diagonal is the
graph of the identity; when \(C\) is separated over \(\Spec k\) it is a closed immersion, and its
image is the closed set along which the per-chart local equation of the next section cuts out the
diagonal.

\begin{definition}[The diagonal of the curve square]
  \label{pa-def:diagonal}
  \lean{AlgebraicGeometry.Over.diagonal,
    AlgebraicGeometry.Over.isClosedImmersion_diagonal_left,
    AlgebraicGeometry.Over.isClosed_range_diagonal,
    AlgebraicGeometry.Over.diagonalComplement,
    AlgebraicGeometry.Over.coe_support_ker_diagonal,
    AlgebraicGeometry.Over.diagonal_preimage_productChart}
  \leanok
  \source{kleiman-picard}

  \uses{pa-def:productChart}
  \dcref{custom}
  The \emph{diagonal} of \(C \otimes C\) is the morphism into the product determined by the
  identity in both components,
  \[
    \delta \;:=\; \langle \id_C, \id_C \rangle \colon C \longrightarrow C \otimes C, \qquad
    \mathrm{fst} \circ \delta = \id_C, \quad \mathrm{snd} \circ \delta = \id_C .
  \]
  For opens \(U, V\) of \(C\) its preimage of a product chart is the intersection of the factors,
  \(\delta^{-1}\mathfrak W(U,V) = U \cap V\). If moreover \(C\) is separated over \(\Spec k\),
  then \(\delta\) is a closed immersion; its image \(\Delta := \delta(C)\) is closed, with open
  complement \(\Delta^{\mathrm{c}} := (C \otimes C) \setminus \Delta\); and the support of the
  kernel ideal sheaf of \(\delta\) is exactly \(\Delta\).
\end{definition}
\begin{proof}
  \leanok
  The morphism \(\langle \id_C, \id_C\rangle\) is the unique one with both composites equal to
  \(\id_C\), by the universal property of the product. A point \(z\) lies in
  \(\delta^{-1}\mathfrak W(U,V)\) iff \(\mathrm{fst}(\delta z) \in U\) and
  \(\mathrm{snd}(\delta z) \in V\); since \(\mathrm{fst} \circ \delta = \mathrm{snd} \circ \delta = \id_C\),
  this means \(z \in U\) and \(z \in V\), i.e. \(z \in U \cap V\).

  Now assume \(C\) separated. The second projection \(\mathrm{snd} \colon C \otimes C \to C\) is
  the base change of the structure morphism \(C \to \Spec k\) along itself, hence separated
  because \(C\) is. The composite \(\mathrm{snd} \circ \delta = \id_C\) is a closed immersion; and
  if a composite \(g \circ f\) is a closed immersion and \(g\) is separated, then \(f\) is a
  closed immersion. Applying this with \(f = \delta\), \(g = \mathrm{snd}\) shows \(\delta\) is a
  closed immersion. A closed immersion is a closed embedding, so its image \(\Delta\) is closed
  and \(\Delta^{\mathrm{c}}\) is open. Finally, the support of the kernel ideal sheaf of a
  quasi-compact morphism is the closure of its image; since \(\Delta\) is closed, this support
  equals \(\Delta\).
\end{proof}

\section{The per-chart local equation of the diagonal}

Fix an affine open \(U\) of \(C\) and write \(B := \Gamma(C, U)\). Suppose \(B\) carries the
abstract \'etale-coordinate data of Definition~\ref{pa-def:diagonalSetup} and
Theorem~\ref{pa-thm:exists_idempotent_lift}: a \(k[X]\)-algebra structure making \(k \to k[X] \to B\)
a tower, an \'etale coordinate \(u \in B\) (the image of \(X\)), and an idempotent lift
\(\mathrm{elift} \in B \otimes_k B\) whose base change to \(B \otimes_{k[X]} B\) is idempotent and
generates the \(k[X]\)-diagonal ideal. On such a chart the diagonal of \(C \otimes C\) is, in one
affine piece, the vanishing locus of the single diagonal equation \(u \otimes 1 - 1 \otimes u\);
this section makes that precise and records the regularity of the equation.

\begin{definition}[The diagonal member and its equation]
  \label{pa-def:diagonalChart}
  \lean{AlgebraicGeometry.Over.diagonalChart,
    AlgebraicGeometry.Over.isAffineOpen_diagonalChart,
    AlgebraicGeometry.Over.diagonalChartEqn,
    AlgebraicGeometry.Over.diagonal_appLE_productChartSections,
    AlgebraicGeometry.Over.diagonal_preimage_diagonalChart}
  \leanok
  \dcref{custom}

  \uses{pa-def:productChart, pa-def:productChartSections, pa-def:diagonal, pa-def:diagonalSetup}
  Let \(\vartheta \in \Gamma(C \otimes C, \mathfrak W(U,U))\) be the section identified with
  \(1 - \mathrm{elift}\) (Definition~\ref{pa-def:productChartSections}). The \emph{diagonal member}
  attached to the chart \(U\) is the affine basic open
  \[
    \mathfrak D(U) \;:=\; D(\vartheta) \;\subseteq\; \mathfrak W(U,U).
  \]
  The \emph{diagonal equation} \(e_U \in \Gamma(C \otimes C, \mathfrak D(U))\) is the restriction
  to \(\mathfrak D(U)\) of the section of \(\mathfrak W(U,U)\) identified with the diagonal
  generator \(u \otimes 1 - 1 \otimes u\). Writing \(\xi_x\) for the section of
  \(\mathfrak W(U,U)\) identified with \(x \in B \otimes_k B\), the pullback along the diagonal is
  multiplication: over \(\delta^{-1}\mathfrak W(U,U) = U\),
  \[
    \delta^\sharp(\xi_x) \;=\; \mu(x) \;\in\; \Gamma(C, U), \qquad
    \mu \colon B \otimes_k B \to B, \quad x \otimes y \mapsto xy .
  \]
  Consequently, if \(\mu(\mathrm{elift}) = 0\), then \(\delta^{-1}\mathfrak D(U) = U\).
\end{definition}
\begin{proof}
  \leanok
  \uses{pa-thm:appLE_productChartSections, pa-thm:isLocalization_away_basicOpen_productChartSections}
  Affineness of \(\mathfrak D(U)\) is affineness of a basic open of the affine chart
  \(\mathfrak W(U,U)\) (Theorem~\ref{pa-thm:isLocalization_away_basicOpen_productChartSections}). For
  the pullback formula, apply the lift-app rule (Theorem~\ref{pa-thm:appLE_productChartSections}) to
  \(q = \delta\), whose components are \(\mathrm{fst} \circ \delta = \mathrm{snd} \circ \delta = \id_C\):
  on a pure tensor \(s \otimes t\) it gives
  \(\delta^\sharp(\xi_{s \otimes t}) = \id^\sharp(s) \cdot \id^\sharp(t) = st = \mu(s \otimes t)\),
  and both sides are additive in \(x\), so \(\delta^\sharp(\xi_x) = \mu(x)\) for every \(x\).
  Finally, \(\mathfrak D(U) = D(\vartheta)\) and \(\vartheta = \xi_{1 - \mathrm{elift}}\), so its
  \(\delta\)-preimage is the basic open of
  \(\delta^\sharp(\vartheta) = \mu(1 - \mathrm{elift}) = 1 - \mu(\mathrm{elift}) = 1\) when
  \(\mu(\mathrm{elift}) = 0\); the basic open of \(1\) is everything, whence
  \(\delta^{-1}\mathfrak D(U) = \delta^{-1}\mathfrak W(U,U) = U \cap U = U\) using
  \(\delta^{-1}\mathfrak W(U,U) = U\) (Definition~\ref{pa-def:diagonal}).
\end{proof}

\begin{theorem}[The per-chart kernel of the diagonal is principal]
  \label{pa-thm:diagonal_ker_ideal_diagonalChart}
  \lean{AlgebraicGeometry.Over.diagonal_ker_ideal_diagonalChart}
  \leanok
  \source{kleiman-picard}

  \uses{pa-def:diagonal, pa-def:diagonalChart}
  \dcref{custom}
  Assume \(C\) is separated over \(\Spec k\). On the diagonal member \(\mathfrak D(U)\) the kernel
  ideal sheaf of the diagonal is principal, generated by the diagonal equation:
  \[
    \bigl(\ker \delta\bigr)\bigl(\mathfrak D(U)\bigr) \;=\; (e_U).
  \]
\end{theorem}
\begin{proof}
  \leanok
  \uses{pa-thm:isLocalization_away_basicOpen_productChartSections, pa-thm:ker_map_localization_eq,
    pa-thm:exists_diagonalLocalization, pa-def:diagonalSetup}
  Because the base change of \(\mathrm{elift}\) generates the \(k[X]\)-diagonal ideal, the
  multiplication kills it, \(\mu(\mathrm{elift}) = 0\)
  (Theorem~\ref{pa-thm:exists_diagonalLocalization}); hence \(\delta^{-1}\mathfrak D(U) = U\)
  (Definition~\ref{pa-def:diagonalChart}), and restriction of sections along this equality of opens
  is a ring isomorphism, in particular injective. On the affine open \(\mathfrak D(U)\) the value
  of the kernel ideal sheaf is the kernel of the section map
  \(\delta^\sharp \colon \Gamma(C \otimes C, \mathfrak D(U)) \to \Gamma(C, U)\).

  By Theorem~\ref{pa-thm:isLocalization_away_basicOpen_productChartSections} the ring
  \(\Gamma(C \otimes C, \mathfrak D(U))\) is the localization of \(B \otimes_k B\) away from
  \(1 - \mathrm{elift}\), through the identification followed by restriction; write
  \(\lambda \colon B \otimes_k B \to \Gamma(C \otimes C, \mathfrak D(U))\) for its structure map,
  so that \(e_U = \lambda(u \otimes 1 - 1 \otimes u)\) and, by the pullback formula of
  Definition~\ref{pa-def:diagonalChart}, \(\delta^\sharp(\lambda(z)) = \mu(z)\big|_{U}\) for all
  \(z \in B \otimes_k B\).

  \emph{The equation lies in the kernel.} One has
  \(\delta^\sharp(e_U) = \mu(u \otimes 1 - 1 \otimes u)\big|_{U} = 0\), because
  \(\mu(u \otimes 1 - 1 \otimes u) = u - u = 0\) (Definition~\ref{pa-def:diagonalSetup}); so
  \((e_U) \subseteq (\ker \delta)(\mathfrak D(U))\).

  \emph{Every kernel element is a multiple of the equation.} Let \(s\) lie in the kernel. As
  \(\Gamma(C \otimes C, \mathfrak D(U))\) is a localization away from \(1 - \mathrm{elift}\), write
  \(s \cdot \lambda(m) = \lambda(r)\) with \(r \in B \otimes_k B\) and \(m\) a power of
  \(1 - \mathrm{elift}\), whose image \(\lambda(m)\) is a unit. Applying \(\delta^\sharp\),
  \(\delta^\sharp(\lambda(r)) = \delta^\sharp(s) \cdot \delta^\sharp(\lambda(m)) = 0\); but
  \(\delta^\sharp(\lambda(r)) = \mu(r)\big|_{U}\) and restriction along
  \(\delta^{-1}\mathfrak D(U) = U\) is injective, so \(\mu(r) = 0\), i.e. \(r\) lies in the
  \(k\)-diagonal ideal \(\ker \mu\). By the localized principality of the \(k\)-diagonal ideal
  (Theorem~\ref{pa-thm:ker_map_localization_eq}), the extension of \(\ker \mu\) along \(\lambda\) is
  generated by \(\lambda(u \otimes 1 - 1 \otimes u) = e_U\); hence \(\lambda(r) \in (e_U)\). Since
  \(\lambda(m)\) is a unit and \(s = \lambda(r) \cdot \lambda(m)^{-1}\), we conclude \(s \in (e_U)\).
\end{proof}

\begin{lemma}[Germs of a regular section on an affine open]
  \label{pa-lem:affineOpen_germ_mem_nonZeroDivisors}
  \lean{AlgebraicGeometry.IsAffineOpen.germ_mem_nonZeroDivisors}
  \leanok
  \dcref{custom}

  \uses{pa-lem:flat_mem_nonZeroDivisors}
  Let \(W\) be an affine open of a scheme \(X\) and \(s \in \Gamma(X, W)\) a nonzerodivisor. Then
  for every point \(y \in W\) the germ of \(s\) at \(y\) is a nonzerodivisor of the stalk
  \(\struct{X, y}\).
\end{lemma}
\begin{proof}
  \leanok
  On an affine open the stalk \(\struct{X, y}\) is the localization of \(\Gamma(X, W)\) at the
  prime corresponding to \(y\), and the germ map is the localization map. A localization is flat,
  so the germ map is a flat ring homomorphism, and a flat ring homomorphism carries the
  nonzerodivisor \(s\) to a nonzerodivisor (Lemma~\ref{pa-lem:flat_mem_nonZeroDivisors}).
\end{proof}

\begin{theorem}[Regularity of the diagonal equation]
  \label{pa-thm:diagonalChartEqn_regular}
  \lean{AlgebraicGeometry.Over.diagonalChartEqn_mem_nonZeroDivisors,
    AlgebraicGeometry.Over.germ_diagonalChartEqn_mem_nonZeroDivisors}
  \leanok
  \dcref{custom}

  \uses{pa-def:diagonalChart}
  Assume in addition that \(B\) is flat over \(k[X]\). Then the diagonal equation \(e_U\) is a
  nonzerodivisor of \(\Gamma(C \otimes C, \mathfrak D(U))\), and its germ at every point of
  \(\mathfrak D(U)\) is a nonzerodivisor of the corresponding stalk.
\end{theorem}
\begin{proof}
  \leanok
  \uses{pa-thm:diagonal_regularity_engine, pa-lem:flat_mem_nonZeroDivisors,
    pa-lem:affineOpen_germ_mem_nonZeroDivisors, pa-def:productChartSections,
    pa-thm:isLocalization_away_basicOpen_productChartSections}
  With \(B\) flat over \(k[X]\), the diagonal generator \(u \otimes 1 - 1 \otimes u\) is a
  nonzerodivisor of \(B \otimes_k B\) (Theorem~\ref{pa-thm:diagonal_regularity_engine}). The structure
  map \(\lambda \colon B \otimes_k B \to \Gamma(C \otimes C, \mathfrak D(U))\) factors as the ring
  isomorphism of Definition~\ref{pa-def:productChartSections} followed by the localization to the
  basic open \(\mathfrak D(U)\)
  (Theorem~\ref{pa-thm:isLocalization_away_basicOpen_productChartSections}); both are flat, so
  \(\lambda\) is flat. A flat ring homomorphism preserves nonzerodivisors
  (Lemma~\ref{pa-lem:flat_mem_nonZeroDivisors}), and \(e_U = \lambda(u \otimes 1 - 1 \otimes u)\),
  so \(e_U\) is a nonzerodivisor. Since \(\mathfrak D(U)\) is affine
  (Definition~\ref{pa-def:diagonalChart}), the germ regularity follows from the section regularity by
  Lemma~\ref{pa-lem:affineOpen_germ_mem_nonZeroDivisors}.
\end{proof}

\begin{theorem}[The diagonal member meets the diagonal complement in a basic open]
  \label{pa-thm:diagonalChart_inf_diagonalComplement}
  \lean{AlgebraicGeometry.Over.diagonalChart_inf_diagonalComplement,
    AlgebraicGeometry.Over.isUnit_res_diagonalChartEqn}
  \leanok
  \dcref{custom}

  \uses{pa-def:diagonal, pa-def:diagonalChart}
  Assume \(C\) is separated over \(\Spec k\). The overlap of the diagonal member with the diagonal
  complement is the basic open of the diagonal equation,
  \[
    \mathfrak D(U) \cap \Delta^{\mathrm{c}} \;=\; D(e_U),
  \]
  and the restriction of \(e_U\) to this overlap is a unit.
\end{theorem}
\begin{proof}
  \leanok
  \uses{pa-thm:diagonal_ker_ideal_diagonalChart}
  On the affine \(\mathfrak D(U)\) the kernel ideal sheaf of the diagonal is generated by \(e_U\)
  (Theorem~\ref{pa-thm:diagonal_ker_ideal_diagonalChart}), so the support of the kernel ideal sheaf,
  intersected with \(\mathfrak D(U)\), is the zero locus \(V(e_U) \cap \mathfrak D(U)\). That
  support is the diagonal \(\Delta\) (Definition~\ref{pa-def:diagonal}), whence
  \(\Delta \cap \mathfrak D(U) = V(e_U) \cap \mathfrak D(U)\). Taking complements inside
  \(\mathfrak D(U)\) gives
  \(\mathfrak D(U) \cap \Delta^{\mathrm{c}} = \mathfrak D(U) \setminus V(e_U) = D(e_U)\). A section
  is a unit on its own basic open, so the restriction of \(e_U\) to
  \(D(e_U) = \mathfrak D(U) \cap \Delta^{\mathrm{c}}\) is a unit.
\end{proof}

\section{The frozen étale-coordinate chart data}

The per-chart local equation of the previous section runs on any affine chart carrying the
abstract étale-coordinate data of Definition~\ref{pa-def:diagonalSetup}. To assemble a single
divisor on the whole square \(C \otimes C\) we must \emph{freeze} one such chart, with its
coordinate and idempotent lift, per point of the curve; this section makes the choice and
records the per-chart deliverables specialised to it.

\begin{definition}[The frozen chart data of a smooth curve]
  \label{pa-def:diagonalChartData}
  \lean{AlgebraicGeometry.Over.exists_diagonalChartAt, AlgebraicGeometry.Over.DiagonalChartData,
    AlgebraicGeometry.Over.diagonalChartData,
    AlgebraicGeometry.Over.DiagonalChartData.elift,
    AlgebraicGeometry.Over.DiagonalChartData.lmul'_elift,
    AlgebraicGeometry.Over.DiagonalChartData.isIdempotentElem_baseChange_elift,
    AlgebraicGeometry.Over.DiagonalChartData.ker_lmul'_eq_span_baseChange_elift}
  \leanok
  \dcref{custom}

  \uses{pa-def:diagonalSetup, pa-thm:exists_diagonalLocalization}
  Let \(C\) be smooth of relative dimension one over \(\Spec k\). Then every point \(p\) of \(C\)
  has an affine open \(U \ni p\) whose section ring \(B := \Gamma(C, U)\) carries an étale
  \(k[X]\)-algebra structure --- an étale coordinate \(u \in B\) --- compatible with the section
  \(k\)-algebra structure (making \(k \to k[X] \to B\) a tower).

  A \emph{chart datum} for \(C\) is a choice, one per point \(p\), of such a chart
  \(\mathrm{chart}(p) \ni p\) with its étale coordinate. On each chart the coordinate is flat,
  formally unramified and essentially of finite type, so the packaged diagonal localisation
  (Theorem~\ref{pa-thm:exists_diagonalLocalization}) supplies a frozen idempotent lift
  \(\mathrm{elift}(p) \in B \otimes_k B\) with \(\mu_k(\mathrm{elift}(p)) = 0\), whose base change
  is idempotent and generates the \(k[X]\)-diagonal ideal.
\end{definition}
\begin{proof}
  \leanok
  Smoothness of relative dimension one over the one-point base \(\Spec k\) yields, at each \(p\), a
  standard-smooth affine chart whose section ring is standard smooth of relative dimension one over
  \(\Gamma(\Spec k, \top)\). Transport the base along the ring isomorphism \(k \cong \Gamma(\Spec
  k, \top)\) --- standard-smoothness is stable under composition with a relative-dimension-zero
  isomorphism, dispatching the base-field-versus-\(\Gamma(\Spec k)\) subtlety at the level of ring
  homomorphisms --- to get \(\Gamma(C, U)\) standard smooth of relative dimension one over \(k\).
  The structure theorem for standard-smooth homomorphisms produces an étale factorisation
  \(k[X_1] \to \Gamma(C, U)\) through the one-variable polynomial ring; transporting the source
  along \(k[X_1] \cong k[X]\) gives the étale coordinate. Étaleness entails flatness, formal
  unramifiedness and essential finiteness, the hypotheses under which
  Theorem~\ref{pa-thm:exists_diagonalLocalization} produces the idempotent lift and its properties.
  Choosing one such chart per point gives the chart datum; every point of the curve has one, so a
  datum exists.
\end{proof}

\begin{theorem}[The frozen diagonal equation on a chart]
  \label{pa-thm:diagonalChartData_equation}
  \lean{AlgebraicGeometry.Over.DiagonalChartData.chartEqn,
    AlgebraicGeometry.Over.DiagonalChartData.ker_ideal_eq_span_chartEqn,
    AlgebraicGeometry.Over.DiagonalChartData.germ_chartEqn_mem_nonZeroDivisors,
    AlgebraicGeometry.Over.DiagonalChartData.chartEqn_inf_diagonalComplement,
    AlgebraicGeometry.Over.DiagonalChartData.span_res_chartEqn,
    AlgebraicGeometry.Over.DiagonalChartData.res_chartEqn_mem_nonZeroDivisors}
  \leanok
  \dcref{custom}

  \uses{pa-def:diagonalChartData, pa-def:diagonalChart, pa-thm:diagonal_ker_ideal_diagonalChart,
    pa-thm:diagonalChartEqn_regular, pa-thm:diagonalChart_inf_diagonalComplement}
  Fix a chart datum on a separated \(C\) and a point \(p\), and write \(\mathfrak D_p\) for the
  diagonal member \(\mathfrak D(\mathrm{chart}(p))\) and \(e_p\) for its diagonal equation
  (Definition~\ref{pa-def:diagonalChart}) built with the frozen coordinate. Then:
  \begin{itemize}
    \item the kernel ideal sheaf of the diagonal on \(\mathfrak D_p\) is principal, \((\ker
      \delta)(\mathfrak D_p) = (e_p)\);
    \item the germ of \(e_p\) at every point of \(\mathfrak D_p\) is a nonzerodivisor;
    \item the diagonal member meets the diagonal complement in the basic open of the equation,
      \(\mathfrak D_p \cap \Delta^{\mathrm c} = D(e_p)\);
    \item for every affine open \(W \le \mathfrak D_p\), the restriction \(e_p|_W\) generates the
      kernel ideal sheaf there, \(\mathrm{span}\{e_p|_W\} = (\ker \delta)(W)\), and \(e_p|_W\) is a
      nonzerodivisor of \(\Gamma(C \otimes C, W)\).
  \end{itemize}
\end{theorem}
\begin{proof}
  \leanok
  The first three are the per-chart deliverables (Theorems~\ref{pa-thm:diagonal_ker_ideal_diagonalChart},
  \ref{pa-thm:diagonalChartEqn_regular}, \ref{pa-thm:diagonalChart_inf_diagonalComplement}) read on the
  frozen chart, whose coordinate instance is baked into the definition of \(e_p\). For the fourth,
  the ideal sheaf of a closed subscheme is compatible with restriction to a smaller affine open, so
  the principal generator on \(\mathfrak D_p\) restricts to a generator on \(W\): \(\mathrm{span}\{e_p|_W\}
  = (\ker \delta)(W)\). The section regularity of \(e_p|_W\) is the propagation of germ regularity:
  a section whose every germ is a nonzerodivisor is itself a nonzerodivisor, because sections inject
  into the product of their germs; the germs of \(e_p|_W\) are those of \(e_p\), which are
  nonzerodivisors.
\end{proof}

\begin{lemma}[Regular generators of a principal ideal differ by a unit]
  \label{pa-lem:exists_isUnit_of_span_singleton_eq}
  \lean{AlgebraicGeometry.Over.exists_isUnit_of_span_singleton_eq}
  \leanok
  \dcref{custom}

  Let \(R\) be a commutative ring and \(a, b \in R\) with \(\mathrm{span}\{a\} = \mathrm{span}\{b\}\).
  If \(a\) is a nonzerodivisor, then \(a\) and \(b\) differ by a unit: there is a unit \(u \in
  R^{\times}\) with \(a = u\,b\).
\end{lemma}
\begin{proof}
  \leanok
  From \(a \in \mathrm{span}\{b\}\) and \(b \in \mathrm{span}\{a\}\) write \(a = c\,b\) and \(b =
  d\,a\). Substituting, \(a = c\,d\,a\), so \((1 - c\,d)\,a = 0\); as \(a\) is a nonzerodivisor this
  forces \(c\,d = 1\), so \(c\) is a unit with inverse \(d\). Taking \(u = c\) gives \(a = u\,b\).
\end{proof}

\begin{definition}[The diagonal pointed cover and its equation]
  \label{pa-def:diagonalCover}
  \lean{AlgebraicGeometry.Over.diagonalCover, AlgebraicGeometry.Over.diagonalEqn,
    AlgebraicGeometry.Over.diagonalCover_opens_of_mem,
    AlgebraicGeometry.Over.diagonalCover_opens_of_notMem,
    AlgebraicGeometry.Over.diagonalEqn_res_of_mem,
    AlgebraicGeometry.Over.diagonalEqn_of_notMem}
  \leanok
  \dcref{custom}

  \uses{pa-def:diagonal, pa-def:diagonalChart, pa-def:diagonalChartData, pa-thm:diagonalChartData_equation}
  Fix a chart datum on a separated \(C\). The \emph{diagonal pointed cover} of \(C \otimes C\)
  assigns to a point \(z\):
  \[
    \mathcal U(z) \;=\;
      \begin{cases}
        \mathfrak D_{\mathrm{fst}(z)} & z \in \Delta \ (\text{i.e. } z \text{ in the range of } \delta),\\
        \Delta^{\mathrm c} & z \notin \Delta,
      \end{cases}
  \]
  the diagonal member of the chart at the first projection of \(z\) on the diagonal, the complement
  off it. The \emph{diagonal equation} is \(e_z = e_{\mathrm{fst}(z)}|_{\mathcal U(z)}\) on the
  diagonal and \(e_z = 1\) off it.
\end{definition}
\begin{proof}
  \leanok
  Each assignment is an open containing \(z\): on the diagonal, \(z = \delta(\mathrm{fst}(z))\)
  lies in \(\mathfrak D_{\mathrm{fst}(z)}\) because \(\delta^{-1}\mathfrak D(U) = U\) contains
  \(\mathrm{fst}(z)\) (Definition~\ref{pa-def:diagonalChart}, using \(\mu_k(\mathrm{elift}) = 0\), and
  \(\mathrm{fst}(z) \in \mathrm{chart}(\mathrm{fst}(z))\)); off the diagonal, \(z \in
  \Delta^{\mathrm c}\) by definition of the complement. The equation is well typed by restriction
  along \(\mathcal U(z) = \mathfrak D_{\mathrm{fst}(z)}\), and its stated restriction behaviour is
  the definition unfolded.
\end{proof}

\begin{theorem}[The diagonal as an effective divisor in local equations]
  \label{pa-thm:diagonalLocalEquations}
  \lean{AlgebraicGeometry.Over.diagonalLocalEquationsOfData,
    AlgebraicGeometry.Over.diagonalPicClassOfData,
    AlgebraicGeometry.Over.diagonalLocalEquations, AlgebraicGeometry.Over.diagonalPicClass}
  \leanok
  \source{kleiman-picard}

  \uses{pa-def:diagonalCover, pa-thm:diagonalChartData_equation, pa-lem:exists_isUnit_of_span_singleton_eq,
    pa-def:localEquations, pa-def:localEquations_picClass, pa-def:diagonalChartData}
  \dcref{custom}
  Let \(C\) be smooth of relative dimension one and separated over \(\Spec k\). The diagonal cover
  and equation of Definition~\ref{pa-def:diagonalCover} (for the frozen chart datum) form a system of
  local equations (Definition~\ref{pa-def:localEquations}) on \(C \otimes C\), cutting out the
  diagonal \(\Delta\). Its Picard class is the class of the diagonal,
  \[
    \mathrm{diagonalPicClass}(C) \;=\; \struct{C \otimes C}(\Delta) \;\in\; \Pic(C \otimes C).
  \]
\end{theorem}
\begin{proof}
  \leanok
  \emph{Regularity.} On a diagonal member the equation is \(e_p\), germ-regular everywhere
  (Theorem~\ref{pa-thm:diagonalChartData_equation}); off the diagonal it is the constant \(1\), whose
  germs are units, hence nonzerodivisors.

  \emph{Overlap ratios are units.} Consider two cover members over points \(x, y\). If both are on
  the diagonal, their overlap \(W = \mathcal U(x) \cap \mathcal U(y)\) is affine and contained in
  both diagonal members; there both \(e_{\mathrm{fst}(x)}|_W\) and \(e_{\mathrm{fst}(y)}|_W\)
  generate the same kernel ideal sheaf \((\ker \delta)(W)\)
  (Theorem~\ref{pa-thm:diagonalChartData_equation}), and \(e_{\mathrm{fst}(x)}|_W\) is a nonzerodivisor,
  so by Lemma~\ref{pa-lem:exists_isUnit_of_span_singleton_eq} the two differ by a unit. If one is on
  the diagonal and the other off, the overlap lies in \(\mathfrak D \cap \Delta^{\mathrm c} =
  D(e_p)\) (Theorem~\ref{pa-thm:diagonalChartData_equation}), where \(e_p\) restricts to a unit, and
  the off-diagonal equation is \(1\); so the ratio is a unit (or its inverse). If both are off the
  diagonal, both equations are \(1\) and the ratio is \(1\). Thus the data is a local-equation
  system.

  Its Picard class (Definition~\ref{pa-def:localEquations_picClass}) is by construction the class of
  the effective divisor it cuts out, which is the diagonal \(\Delta\); that is, \(\struct{C \otimes
  C}(\Delta)\). The diagonal is an effective (relative Cartier) divisor precisely because \(C\) is
  smooth, matching Kleiman's criterion.
\end{proof}

\section{The graph of a curve point as an effective divisor}

Fix \(C\) smooth of relative dimension one and separated over \(\Spec k\). For a test object \(T\)
and a point \(t \colon T \to C\) (a morphism of \(\Over(\Spec k)\)), the \emph{graph}
\(\Gamma_t \subseteq C \otimes T\) is realised as the pullback of the diagonal divisor along the
graph lift \(\langle \mathrm{fst}, \mathrm{snd} \circ t\rangle \colon C \otimes T \to C \otimes C\).
The pullback of the local-equation system is available once the pulled equations stay regular; this
is the content of the section.

\begin{lemma}[The graph-lift whisker square]
  \label{pa-lem:whiskerLeft_comp_graphLift}
  \lean{AlgebraicGeometry.Over.whiskerLeft_comp_graphLift}
  \leanok
  \dcref{custom}

  For \(g \colon T' \to T\) and \(t \colon T \to C\), whiskering the graph lift of \(t\) by
  \(C \lhd g\) is the graph lift of the restricted point \(g \circ t\):
  \[
    (C \lhd g) \circ \langle \mathrm{fst}_{C,T},\, \mathrm{snd}_{C,T} \circ t\rangle
      \;=\; \langle \mathrm{fst}_{C,T'},\, \mathrm{snd}_{C,T'} \circ g \circ t\rangle .
  \]
\end{lemma}
\begin{proof}
  \leanok
  Both sides are maps into the product \(C \otimes C\); by the universal property it suffices to
  compare their two projections. Post-composing with \(\mathrm{fst}\): the left side gives \((C
  \lhd g) \circ \mathrm{fst}_{C,T} = \mathrm{fst}_{C,T'}\), the right side gives
  \(\mathrm{fst}_{C,T'}\). Post-composing with \(\mathrm{snd}\): the left side gives \((C \lhd g)
  \circ \mathrm{snd}_{C,T} \circ t = \mathrm{snd}_{C,T'} \circ g \circ t\) (as \(\mathrm{snd} \circ
  (C \lhd g) = g \circ \mathrm{snd}\)), matching the right side.
\end{proof}

\begin{theorem}[Regularity of the pulled-back diagonal equation]
  \label{pa-thm:graph_pullback_regular}
  \lean{AlgebraicGeometry.Over.graph_pullback_regular}
  \leanok
  \dcref{custom}

  \uses{pa-def:diagonalCover, pa-thm:diagonalChartData_equation, pa-def:diagonalChart, pa-def:productChart,
    pa-def:productChartSections, pa-thm:appLE_productChartSections, pa-thm:diagonal_regularity_engine,
    pa-lem:flat_mem_nonZeroDivisors, pa-lem:affineOpen_germ_mem_nonZeroDivisors,
    pa-def:localEquations_pullback, pa-def:diagonalChartData, pa-def:diagonalSetup}
  For a chart datum on a separated \(C\) and a point \(t \colon T \to C\), the pullback of the
  diagonal equation along the graph lift \(q = \langle \mathrm{fst}, \mathrm{snd} \circ t\rangle\)
  is germ-regular at every point of the pullback cover.
\end{theorem}
\begin{proof}
  \leanok
  Fix a point \(y\) of the pullback cover and a point \(z\) of its member. Off the diagonal (i.e.\
  \(q(y) \notin \Delta\)) the diagonal equation at \(q(y)\) is \(1\), so its \(q\)-pullback is
  \(q^{\sharp}(1) = 1\), whose germ is a unit, hence a nonzerodivisor.

  On the diagonal, set \(p_0 = \mathrm{fst}(q(y))\), let \(U = \mathrm{chart}(p_0)\) with étale
  coordinate \(u\), so the member is \(q^{-1}\mathfrak D(U)\). Choose an affine open \(V\) of \(T\)
  with \((\mathrm{snd})(z) \in V \subseteq t^{-1}U\), and pass to the affine product chart
  \(\mathfrak W(U, V)\) (Definition~\ref{pa-def:productChart}) through \(z\); on the overlap \(W\) of
  the member with \(\mathfrak W(U, V)\), the pulled equation restricts to the product-chart section
  of \(u \otimes 1 - 1 \otimes b\), where \(b := t^{\sharp}(u) \in \Gamma(T, V)\). This is the
  lift-app rule (Theorem~\ref{pa-thm:appLE_productChartSections}) for \(q\), whose components are
  \(\mathrm{fst} \circ q = \mathrm{fst}\) and \(\mathrm{snd} \circ q = \mathrm{snd} \circ t\): the
  diagonal generator \(u \otimes 1 - 1 \otimes u\) (Definition~\ref{pa-def:diagonalSetup}) pulls back
  to \(\mathrm{fst}^{\sharp}(u) - (\mathrm{snd} \circ t)^{\sharp}(u) = \mathrm{fst}^{\sharp}(u) -
  \mathrm{snd}^{\sharp}(b)\), i.e.\ the identified \(u \otimes 1 - 1 \otimes b\)
  (Definition~\ref{pa-def:productChartSections}).

  Now \(u \otimes 1 - 1 \otimes b\) is a nonzerodivisor of \(\Gamma(C, U) \otimes_k \Gamma(T, V)\)
  by the regularity engine (Theorem~\ref{pa-thm:diagonal_regularity_engine}), which holds uniformly in
  the second tensor factor \(\Gamma(T, V)\) and the element \(b\) --- no integrality of \(C \otimes
  T\) is used. The product-chart identification is a ring isomorphism, hence flat, so it carries
  this nonzerodivisor to a nonzerodivisor of \(\Gamma(C \otimes T, \mathfrak W(U, V))\)
  (Lemma~\ref{pa-lem:flat_mem_nonZeroDivisors}). Since \(\mathfrak W(U, V)\) is affine, the germ at
  \(z\) is a nonzerodivisor (Lemma~\ref{pa-lem:affineOpen_germ_mem_nonZeroDivisors}); matching this
  germ to the pullback germ across the overlap \(W\) (germs are compatible with restriction) gives
  the claim.
\end{proof}

\begin{definition}[The graph divisor and its Picard class]
  \label{pa-def:graphLocalEquations}
  \lean{AlgebraicGeometry.Over.graphLocalEquations, AlgebraicGeometry.Over.graphPicClass,
    AlgebraicGeometry.Over.graphPicClass_eq}
  \leanok
  \source{kleiman-picard}

  \uses{pa-thm:diagonalLocalEquations, pa-def:localEquations_pullback, pa-thm:graph_pullback_regular,
    pa-thm:localEquations_picClass_pullback}
  \dcref{custom}
  Let \(C\) be smooth of relative dimension one and separated over \(\Spec k\), \(T\) a test object
  and \(t \colon T \to C\) a point. The \emph{graph divisor} \(\Gamma_t \subseteq C \otimes T\) is
  cut out by the pullback of the diagonal local-equation system
  (Theorem~\ref{pa-thm:diagonalLocalEquations}) along the graph lift \(\langle \mathrm{fst},
  \mathrm{snd} \circ t\rangle\); the pullback is a local-equation system by
  Theorem~\ref{pa-thm:graph_pullback_regular}. Its Picard class \(\mathrm{graphPicClass}(C, t) =
  \struct{C \otimes T}(\Gamma_t)\) is the pullback of the diagonal class:
  \[
    \struct{C \otimes T}(\Gamma_t)
      \;=\; \langle \mathrm{fst}, \mathrm{snd} \circ t\rangle^{*}\,\struct{C \otimes C}(\Delta).
  \]
\end{definition}
\begin{proof}
  \leanok
  The diagonal is a local-equation system (Theorem~\ref{pa-thm:diagonalLocalEquations}); its pullback
  along the graph lift is again a local-equation system by
  Definition~\ref{pa-def:localEquations_pullback}, whose regularity hypothesis is discharged by
  Theorem~\ref{pa-thm:graph_pullback_regular}. The Picard class of a pullback of a local-equation
  system is the pullback of its Picard class (Theorem~\ref{pa-thm:localEquations_picClass_pullback}),
  giving the displayed identity. This is Kleiman's graph divisor \(\Gamma_t\), a relative effective
  divisor on \(X \times T\).
\end{proof}

\begin{theorem}[Base change of the graph class]
  \label{pa-thm:graphLocalEquations_base_change}
  \lean{AlgebraicGeometry.Over.graphLocalEquations_base_change}
  \leanok
  \dcref{custom}

  \uses{pa-def:graphLocalEquations, pa-lem:whiskerLeft_comp_graphLift}
  For \(g \colon T' \to T\) and \(t \colon T \to C\), the class of the graph of the restricted
  point \(g \circ t\) is the pullback of the class of the graph of \(t\) along \(C \lhd g\):
  \[
    \struct{C \otimes T'}(\Gamma_{g \circ t})
      \;=\; (C \lhd g)^{*}\,\struct{C \otimes T}(\Gamma_t).
  \]
\end{theorem}
\begin{proof}
  \leanok
  By Definition~\ref{pa-def:graphLocalEquations} each side is a pullback of the diagonal class
  \(\struct{C \otimes C}(\Delta)\): the left along the graph lift of \(g \circ t\), the right along
  the graph lift of \(t\) followed by \(C \lhd g\). Contravariant functoriality of the pullback of
  Picard classes turns the latter into the pullback along the composite \((C \lhd g) \circ
  \langle\mathrm{fst}, \mathrm{snd} \circ t\rangle\), which is the graph lift of \(g \circ t\) by
  the whisker square (Lemma~\ref{pa-lem:whiskerLeft_comp_graphLift}). Hence the two pullbacks agree.
\end{proof}

\section{The graph square and the graph fibre}
\label{pa-sec:graphFibre}

The graph divisor of the previous section is supported, over a field test, on a single
point: the graph is the base change of the diagonal, and the fibre of a one-point test is
one point. This section proves the pullback square and extracts the support statement.

\begin{lemma}[The graph square commutes]
  \label{pa-lem:sectionOfPoint_comp_graphLift}
  \lean{AlgebraicGeometry.Over.sectionOfPoint_comp_graphLift}
  \leanok
  \dcref{custom}

  \uses{pa-def:sectionOfPoint, pa-def:diagonal}
  For a point \(t \colon T \to C\), the section of the projection followed by the graph
  lift is the point followed by the diagonal:
  \[
    \langle \mathrm{fst}, \mathrm{snd} \circ t\rangle \circ g_t
      \;=\; \Delta \circ t
    \qquad\text{as morphisms } T \to C \otimes C .
  \]
\end{lemma}
\begin{proof}
  \leanok
  Compare the two projections. Onto the first factor: \(\mathrm{fst} \circ q \circ g_t =
  \mathrm{fst} \circ g_t = t\) and \(\mathrm{fst} \circ \Delta \circ t = t\). Onto the
  second: \(\mathrm{snd} \circ q \circ g_t = t \circ \mathrm{snd} \circ g_t = t\) and
  \(\mathrm{snd} \circ \Delta \circ t = t\), using the projection identities of the section
  (Definition~\ref{pa-def:sectionOfPoint}) and of the diagonal (Definition~\ref{pa-def:diagonal}).
\end{proof}

\begin{theorem}[The graph square is a pullback]
  \label{pa-thm:isPullback_graphLift}
  \lean{AlgebraicGeometry.Over.isPullback_graphLift}
  \leanok
  \source{kleiman-picard}

  \uses{pa-lem:sectionOfPoint_comp_graphLift, pa-lem:over_isPullback_left, pa-def:sectionOfPoint,
    pa-def:diagonal}
  \dcref{custom}
  For any test object \(T\) and point \(t \colon T \to C\), the commuting square of
  underlying schemes
  \[
    \begin{array}{ccc}
      T_{\mathrm{left}} & \xrightarrow{\ t\ } & C_{\mathrm{left}} \\
      \downarrow g_t & & \downarrow \Delta \\
      (C \otimes T)_{\mathrm{left}} & \xrightarrow{\ q\ } & (C \otimes C)_{\mathrm{left}}
    \end{array}
  \]
  (\(q = \langle\mathrm{fst}, \mathrm{snd} \circ t\rangle\) the graph lift) is a pullback:
  the graph of \(t\) is the base change of the diagonal along the graph lift — Kleiman's
  \(\Gamma_g = (1 \times g)^{-1}\Delta\). No separatedness or properness is used.
\end{theorem}
\begin{proof}
  \leanok
  Two paste-cancellations on the product bridge (Lemma~\ref{pa-lem:over_isPullback_left}).
  First, the middle square: the square with horizontals \(q\) and \(t\) and verticals the
  two second projections is a pullback, by right-cancellation of the product square of
  \(C \otimes C\) against the product square of \(C \otimes T\) (whose first projection
  factors as \(\mathrm{fst} \circ q\)). Second, the top square: the square with verticals
  \(g_t\) and \(\Delta\) pastes with the middle square to the square with verticals
  \(g_t \circ \mathrm{snd} = \id\) and \(\Delta \circ \mathrm{snd} = \id\), which is a
  pullback (vertical isomorphisms); bottom-cancellation against the middle square, using
  commutativity (Lemma~\ref{pa-lem:sectionOfPoint_comp_graphLift}), makes the top square a
  pullback.
\end{proof}

\begin{theorem}[The range of the section is the graph]
  \label{pa-thm:range_sectionOfPoint_left}
  \lean{AlgebraicGeometry.Over.range_sectionOfPoint_left}
  \leanok
  \dcref{custom}

  \uses{pa-thm:isPullback_graphLift}
  For any point \(t \colon T \to C\), the topological range of the section of the
  projection is the graph-lift preimage of the diagonal:
  \[
    \operatorname{range}(g_t) \;=\; q^{-1}\bigl(\operatorname{range}\Delta\bigr)
    \quad\text{in } (C \otimes T)_{\mathrm{left}} .
  \]
\end{theorem}
\begin{proof}
  \leanok
  A pullback of schemes surjects onto the topological fibre product
  (Theorem~\ref{pa-thm:isPullback_graphLift} plus the surjectivity of scheme pullbacks onto
  fibre products of spaces); the fibre product of \(\Delta\) and \(q\) projects onto
  exactly the points of \((C \otimes T)_{\mathrm{left}}\) whose \(q\)-image lies in the
  range of \(\Delta\), and the projection to \((C \otimes T)_{\mathrm{left}}\) is carried
  by the square onto the range of \(g_t\).
\end{proof}

\begin{definition}[The graph point of a field test]
  \label{pa-def:graphPoint}
  \lean{AlgebraicGeometry.Over.graphPoint,
    AlgebraicGeometry.Over.range_sectionOfPoint_left_eq_singleton,
    AlgebraicGeometry.Over.mem_range_diagonal_graphLift_iff,
    AlgebraicGeometry.Over.fst_graphPoint, AlgebraicGeometry.Over.graphLift_graphPoint}
  \leanok
  \dcref{custom}

  \uses{pa-thm:range_sectionOfPoint_left, pa-def:sectionOfPoint}
  Let \(K\) be a field extension of \(k\) and \(t \colon \Spec K \to C\) a field point. The
  \emph{graph point} \(x_t \in (C \otimes \Spec K)_{\mathrm{left}}\) is the image of the
  unique point of \(\Spec K\) under the section \(g_t\). It is the entire graph fibre:
  \[
    q(z) \in \operatorname{range}\Delta \iff z = x_t ,
  \]
  its first projection is the image point of \(t\), and \(q(x_t) = \Delta(t(\ast))\).
\end{definition}
\begin{proof}
  \leanok
  \(\Spec K\) has a single point, so
  \(\operatorname{range}(g_t) = \{x_t\}\); by
  Theorem~\ref{pa-thm:range_sectionOfPoint_left} the \(q\)-preimage of the diagonal is that
  singleton, which is the displayed equivalence. The two projection formulas follow from
  \(\mathrm{fst} \circ g_t = t\) and the commuting graph square
  (Definition~\ref{pa-def:sectionOfPoint},
  Lemma~\ref{pa-lem:sectionOfPoint_comp_graphLift}).
\end{proof}

\section{Rational points over a separably closed field}
\label{pa-sec:separablyClosedPoints}

Over a separably closed field \(K\) the \(K\)-rational points of a smooth curve are
\emph{dense}: every nonempty open contains one. Over an imperfect separably closed field
not every closed point is rational (on \(\mathbb{A}^1_K\) there are closed points with residue
field \(K(a^{1/p})\)), so the algebraically-closed argument through closed points fails;
what survives is density, through the standard-smooth chart, which is \'etale over the
affine line. There is no classical source anchored here; the statements are as formalised.

\begin{lemma}[A separably closed field is infinite]
  \label{pa-lem:sepClosed_infinite}
  \lean{AlgebraicGeometry.SeparablyClosed.instInfiniteOfIsSepClosed}
  \leanok
  \dcref{custom}

  A separably closed field is infinite.
\end{lemma}
\begin{proof}
  \leanok
  Suppose \(K\) had exactly \(n\) elements. The polynomial \(f = X^{n+1} - 1\) is separable
  over \(K\): in characteristic zero this is clear, and in characteristic \(p > 0\) the
  cardinality \(n\) is a power of \(p\), so \(n + 1 \equiv 1 \not\equiv 0 \pmod p\); hence
  the derivative \((n+1)X^{n}\) is a nonzero monomial, sharing no root with \(f\) (a common
  root would satisfy \(x = 0\) and \(x^{n+1} = 1\)). A separable polynomial over a separably
  closed field splits with distinct roots, so \(K\) contains \(n + 1\) distinct roots of
  \(f\) — more than \(n\) elements, a contradiction. (This mirrors the classical argument
  that an algebraically closed field is infinite, with the splitting input replaced by
  separable splitting.)
\end{proof}

\begin{definition}[Rational points of the affine line]
  \label{pa-def:ratPt}
  \lean{AlgebraicGeometry.SeparablyClosed.ratPt,
    AlgebraicGeometry.SeparablyClosed.mem_ratPt_basicOpen}
  \leanok
  \dcref{custom}

  For a field \(K\) and a value \(x \in K\), the \(K\)-rational point of
  \(\mathbb{A}^1 = \Spec K[X]\) at \(x\): the maximal ideal \(\ker(\mathrm{ev}_x) = (X - x)\). It
  lies in a basic open \(D(h)\) iff \(h(x) \neq 0\).
\end{definition}
\begin{proof}
  \leanok
  \(\mathrm{ev}_x\) is a surjection onto the field \(K\), so its kernel is maximal, hence
  prime. Membership in \(D(h)\) means \(h \notin \ker(\mathrm{ev}_x)\), i.e.\
  \(h(x) \neq 0\).
\end{proof}

\begin{lemma}[A rational base point under an \'etale algebra]
  \label{pa-lem:exists_ratPt_mem_range}
  \lean{AlgebraicGeometry.SeparablyClosed.exists_ratPt_mem_range}
  \leanok
  \dcref{custom}

  \uses{pa-def:ratPt, pa-lem:sepClosed_infinite}
  Let \(K\) be an infinite field and \(S \neq 0\) a \(K\)-algebra \'etale over
  \(K[X]\). Then the image of \(\Spec S \to \mathbb{A}^1\) contains a \(K\)-rational point.
\end{lemma}
\begin{proof}
  \leanok
  \'Etale maps are flat and of finite presentation, so the image of \(\Spec S\) is open
  (going-down plus finite presentation makes the induced map on spectra open); it is
  nonempty since \(S \neq 0\). A nonempty open of \(\Spec K[X]\) contains a nonempty basic
  open \(D(h)\), \(h \neq 0\). Over the infinite field \(K\) the nonzero polynomial \(h\)
  has a non-vanishing value \(x\) (a nonzero polynomial function on an infinite field is
  not identically zero), and \(\mathrm{ratPt}(x) \in D(h)\)
  (Definition~\ref{pa-def:ratPt}) lies in the image.
\end{proof}

\begin{lemma}[\'Etale algebras over a separably closed field split]
  \label{pa-lem:exists_algHom_to_base}
  \lean{AlgebraicGeometry.SeparablyClosed.exists_algHom_to_base}
  \leanok
  \dcref{custom}

  A nonzero \'etale algebra \(A\) over a separably closed field \(K\) admits a \(K\)-algebra
  homomorphism \(A \to K\).
\end{lemma}
\begin{proof}
  \leanok
  Over a separably closed field, a (formally) \'etale algebra splits as a finite product of
  copies of \(K\); \(A \neq 0\) means the product is nonempty, and any factor projection is
  a \(K\)-point.
\end{proof}

\begin{theorem}[The commutative-algebra core of the density theorem]
  \label{pa-thm:exists_algHom_of_etale_mvPoly}
  \lean{AlgebraicGeometry.SeparablyClosed.exists_algHom_of_etale_mvPoly}
  \leanok
  \dcref{custom}

  \uses{pa-lem:exists_ratPt_mem_range, pa-lem:exists_algHom_to_base, pa-lem:sepClosed_infinite,
    pa-def:ratPt}
  Let \(K\) be separably closed and \(S \neq 0\) a \(K\)-algebra \'etale over \(K[X]\).
  Then \(S\) has a \(K\)-point: a \(K\)-algebra homomorphism \(S \to K\).
\end{theorem}
\begin{proof}
  \leanok
  \(K\) is infinite (Lemma~\ref{pa-lem:sepClosed_infinite}); pick a \(K\)-rational point
  \(x\) of \(\mathbb{A}^1\) in the image of \(\Spec S\)
  (Lemma~\ref{pa-lem:exists_ratPt_mem_range}). The fibre
  \(S_x := K \otimes_{K[X]} S\) (base change along \(\mathrm{ev}_x\)) is \'etale over \(K\)
  (base change of \'etale is \'etale) and nonzero: a prime \(q\) of \(S\) over
  \(\mathrm{ratPt}(x)\) yields a \(K[X]\)-algebra map \(S_x \to S/q\) — lift
  \(K \to S/q\) through the surjection \(K[X] \twoheadrightarrow K\), whose kernel dies in
  \(S/q\) precisely because \(q\) lies over \(\ker(\mathrm{ev}_x)\), and pair it with the
  residue map — and \(S/q \neq 0\). By the splitting lemma
  (Lemma~\ref{pa-lem:exists_algHom_to_base}) the fibre has a \(K\)-point, which composes with
  \(S \to K \otimes_{K[X]} S\), \(s \mapsto 1 \otimes s\), to a \(K\)-point of \(S\).
\end{proof}

\begin{theorem}[A nonempty smooth curve over a separably closed field has a rational point]
  \label{pa-thm:exists_rationalPoint_sepClosed}
  \lean{AlgebraicGeometry.SeparablyClosed.exists_rationalPoint_of_smoothOfRelativeDimension_one}
  \leanok
  \dcref{custom}

  \uses{pa-thm:exists_algHom_of_etale_mvPoly, pa-thm:standard_smooth_etale_structure}
  Let \(K\) be separably closed and \(X\) a nonempty scheme smooth of relative dimension
  \(1\) over \(K\). Then \(X\) has a \(K\)-point over \(K\): a morphism
  \(p \colon \Spec K \to X\) with \(p \circ f = \id\).
\end{theorem}
\begin{proof}
  \leanok
  Pick a point \(x_0 \in X\) and a standard-smooth chart \(V \ni x_0\) of relative
  dimension \(1\). Its section ring is \'etale over a one-variable polynomial ring
  (Theorem~\ref{pa-thm:standard_smooth_etale_structure}), and transporting along the
  identification \(\Gamma(\Spec K, \top) \cong K\) makes \(\Gamma(X, V)\) an \'etale
  \(K[X]\)-algebra, nonzero since \(V \ni x_0\) is nonempty. The core
  (Theorem~\ref{pa-thm:exists_algHom_of_etale_mvPoly}) yields a \(K\)-algebra map
  \(\varphi \colon \Gamma(X,V) \to K\); \(\Spec\varphi\) composed with the canonical map of
  the affine open \(V\) is a scheme point \(\Spec K \to X\). It lies over \(K\): composing
  with the structure morphism corresponds, on rings, to
  \(\varphi \circ (\text{chart map } K \to \Gamma(X,V))\), which is the identity of \(K\)
  since \(\varphi\) is a \(K\)-algebra map.
\end{proof}

\begin{theorem}[Density of rational points over a separably closed field]
  \label{pa-thm:exists_rationalPoint_mem}
  \lean{AlgebraicGeometry.SeparablyClosed.exists_rationalPoint_mem,
    AlgebraicGeometry.Over.exists_rationalPoint_mem}
  \leanok
  \dcref{custom}

  \uses{pa-thm:exists_rationalPoint_sepClosed}
  Let \(K\) be separably closed, \(f \colon X \to \Spec K\) smooth of relative dimension
  \(1\), and \(W \subseteq X\) a nonempty open. Then there is a \(K\)-point
  \(p \colon \Spec K \to X\) over \(K\) whose topological image lies in \(W\). In
  particular, for a curve pack \(C\) over a separably closed field, every nonempty open of
  the underlying space of \(C\) contains a rational point.
\end{theorem}
\begin{proof}
  \leanok
  Smoothness of relative dimension \(1\) is local on the source, so the open subscheme
  \(W\) is itself smooth of relative dimension \(1\) over \(K\), and nonempty; by
  Theorem~\ref{pa-thm:exists_rationalPoint_sepClosed} it has a \(K\)-point over \(K\), whose
  composite with the open immersion \(W \to X\) is the required point of \(X\), with image
  in \(W\).
\end{proof}

\section{Coheight and the Krull dimension of stalks}
\label{pa-sec:coheight_stalk}

For a point \(z\) of a topological space, the \emph{coheight} of \(z\) is the supremum
of the lengths of strictly increasing chains starting at \(z\) in the specialisation
preorder (\(y \leq z\) iff \(z\) specialises to \(y\), i.e.\ \(y \in \overline{\{z\}}\));
on a scheme it measures the codimension of the closure of \(z\). This section identifies
it with the Krull dimension of the stalk.

\begin{lemma}[Coheight is local on opens]
  \label{pa-lem:coheight_open_embedding}
  \lean{Order.coheight_eq_of_isOpenEmbedding}
  \leanok
  \dcref{custom}

  Let \(X\) be a topological space, \(U \subseteq X\) an open subset and \(z \in U\).
  The coheight of \(z\) computed in \(X\) (for the specialisation preorder) equals the
  coheight of \(z\) computed in the subspace \(U\).
\end{lemma}
\begin{proof}
  \leanok
  (\(\geq\)) The inclusion \(U \to X\) is continuous, hence monotone for the
  specialisation preorders, and injective; a strictly increasing chain in \(U\) starting
  at \(z\) therefore maps to one in \(X\) of the same length (strictness is preserved:
  if the images were comparable both ways they would be equal).

  (\(\leq\)) Conversely, let \(z = p_0 < p_1 < \dots < p_n\) be a chain in \(X\)
  starting at \(z\). Each \(p_i\) is a generisation of \(z\): \(z \in
  \overline{\{p_i\}}\). Since \(U\) is open and \(z \in U\), every generisation of a
  point of \(U\) lies in \(U\) (an open set containing a point of the closure of
  \(\{p_i\}\) meets \(\{p_i\}\)), so all \(p_i \in U\). The specialisation preorder of
  the subspace \(U\) agrees with that of \(X\) on points of \(U\) when \(U\) is open, so
  the chain lives in \(U\) with the same strict inequalities and the same length.
\end{proof}

\begin{lemma}[Coheight in an affine scheme is the height of the prime]
  \label{pa-lem:coheight_spec_height}
  \lean{Order.coheight_spec_eq_height_primeSpectrum}
  \leanok
  \dcref{custom}

  Let \(R\) be a commutative ring and \(p\) a point of \(\Spec R\), with corresponding
  prime ideal \(\mathfrak p\). The coheight of \(p\) in the specialisation preorder of
  \(\Spec R\) equals the height of \(\mathfrak p\) (the supremum of lengths of chains of
  primes descending from \(\mathfrak p\)).
\end{lemma}
\begin{proof}
  \leanok
  In \(\Spec R\), a point \(q\) specialises to \(p\) exactly when
  \(\mathfrak p \supseteq \mathfrak q\); so the specialisation preorder on
  \(\Spec R\) is the opposite of the inclusion order on primes, and the order
  isomorphism \(\Spec R \cong (\text{primes of } R)^{\mathrm{op}}\) converts coheight at
  \(p\) into height at \(\mathfrak p\) (chains going up from \(p\) in the
  specialisation preorder are chains of primes going down from \(\mathfrak p\)).
\end{proof}

\begin{theorem}[Stalk dimension equals coheight]
  \label{pa-thm:stalk_dim_coheight}
  \lean{AlgebraicGeometry.Scheme.ringKrullDim_stalk_eq_coheight}
  \leanok
  \dcref{custom}

  \uses{pa-lem:coheight_open_embedding, pa-lem:coheight_spec_height}
  Let \(X\) be a scheme and \(z \in X\) a point. Then the Krull dimension of the stalk
  \(\struct{X,z}\) equals the coheight of \(z\) in the underlying topological space of
  \(X\).
\end{theorem}
\begin{proof}
  \leanok
  Choose an affine open \(U \ni z\), say \(U \cong \Spec A\) with
  \(A = \Gamma(X, U)\), and let \(\mathfrak p\) be the prime of \(A\) corresponding to
  \(z\). The stalk \(\struct{X,z}\) is the localisation \(A_{\mathfrak p}\), whose Krull
  dimension is the height of \(\mathfrak p\) (primes of the localisation are the primes
  of \(A\) contained in \(\mathfrak p\)). By \ref{pa-lem:coheight_open_embedding}, the
  coheight of \(z\) in \(X\) equals its coheight in \(U\); the isomorphism
  \(U \cong \Spec A\) is a homeomorphism, hence an isomorphism of specialisation
  preorders, so this equals the coheight of \(\mathfrak p\) in \(\Spec A\), which is the
  height of \(\mathfrak p\) by \ref{pa-lem:coheight_spec_height}. Chaining the equalities
  gives \(\dim \struct{X,z} = \mathrm{coheight}(z)\).
\end{proof}

\begin{lemma}[Codimension-one points have one-dimensional stalks]
  \label{pa-lem:stalk_dimLE_one}
  \lean{AlgebraicGeometry.Scheme.ringKrullDimLE_of_coheight_eq_one}
  \leanok
  \dcref{custom}

  \uses{pa-thm:stalk_dim_coheight}
  Let \(X\) be a scheme and \(z \in X\) with \(\mathrm{coheight}(z) = 1\). Then the
  stalk \(\struct{X,z}\) has Krull dimension \(\leq 1\).
\end{lemma}
\begin{proof}
  \leanok
  By \ref{pa-thm:stalk_dim_coheight}, \(\dim \struct{X,z} = \mathrm{coheight}(z) = 1\).
\end{proof}

\section{The dimension of a standard-smooth algebra at a maximal ideal}
\label{pa-sec:standard_smooth_dimension}

This section proves the dimension half of the smooth-implies-regular pipeline at closed
points: for a standard-smooth algebra of relative dimension \(n\) over a field, every
maximal ideal has height at least \(n\).

\begin{theorem}[Maximal ideals of polynomial rings have full height]
  \label{pa-thm:mvpoly_maximal_height}
  \lean{MvPolynomial.height_eq_natCard_of_isMaximal}
  \leanok
  \dcref{custom}

  Let \(k\) be a field, \(\iota\) a finite index type, and \(M\) a maximal ideal of the
  polynomial ring \(k[x_i : i \in \iota]\). Then the height of \(M\) is \(\#\iota\).
\end{theorem}
\begin{proof}
  \leanok
  Induct on the finite type \(\iota\) (along equivalences, the empty type, and adding
  one variable). Heights are invariant under ring isomorphisms, which handles transport
  along an equivalence \(\iota \cong \iota'\). If \(\iota\) is empty the ring is the
  field \(k\), of Krull dimension \(0\), so every maximal ideal has height \(0 =
  \#\iota\). For \(\iota = \{*\} \sqcup \alpha\), identify
  \(k[x_i : i \in \iota] \cong R[X]\) with \(R = k[x_i : i \in \alpha]\), and let
  \(M \subseteq R[X]\) be maximal. The ring \(R\), a finitely generated algebra over a
  field, is a Jacobson ring, so the contraction \(\mathfrak p = M \cap R\) is again
  \emph{maximal}. For a maximal ideal of \(R[X]\) lying over a maximal ideal
  \(\mathfrak p\) of a Noetherian ring one has the height formula
  \(\mathrm{ht}(M) = \mathrm{ht}(\mathfrak p) + 1\) (the fibre
  \(R[X] \otimes_R R/\mathfrak p = (R/\mathfrak p)[X]\) is a principal ideal domain over
  the field \(R/\mathfrak p\), and \(M\) generates a nonzero --- because maximal ---
  prime there). By the inductive hypothesis \(\mathrm{ht}(\mathfrak p) = \#\alpha\), so
  \(\mathrm{ht}(M) = \#\alpha + 1 = \#\iota\).
\end{proof}

\begin{theorem}[Height lower bound at maximal ideals of standard-smooth algebras]
  \label{pa-thm:standard_smooth_height_lower}
  \lean{Algebra.IsStandardSmoothOfRelativeDimension.natCast_le_height_of_isMaximal}
  \leanok
  \source{stacks-project}

  \uses{pa-thm:mvpoly_maximal_height}
  \dcref{custom}
  Let \(k\) be a field and \(S\) a standard-smooth \(k\)-algebra of relative dimension
  \(n\). Then every maximal ideal \(\mathfrak m \subseteq S\) has
  \(\mathrm{ht}(\mathfrak m) \geq n\).
\end{theorem}
\begin{proof}
  \leanok
  Choose a submersive presentation \(S = P/\ker\), where
  \(P = k[x_i : i \in \iota]\) is a polynomial ring on \(\#\iota\) generators, the
  kernel of \(P \twoheadrightarrow S\) is generated by \(\#\sigma\) relations, and
  \(n = \#\iota - \#\sigma\) (with \(\#\sigma \leq \#\iota\)). The pullback
  \(M \subseteq P\) of \(\mathfrak m\) is a maximal ideal (the quotient map is
  surjective), of height \(\#\iota\) by \ref{pa-thm:mvpoly_maximal_height}, and it contains
  the kernel. Krull's height theorem, in the relative form
  \(\mathrm{ht}(M) \leq \mathrm{ht}(M/\ker) + \mu(\ker)\) for a prime containing an
  ideal generated by \(\mu(\ker)\) elements, gives
  \[
    \#\iota = \mathrm{ht}(M) \leq \mathrm{ht}(\mathfrak m) + \#\sigma
  \]
  (the first isomorphism theorem identifies \(M/\ker\) with \(\mathfrak m\), and heights
  are preserved by this identification). Subtracting \(\#\sigma\) yields
  \(n = \#\iota - \#\sigma \leq \mathrm{ht}(\mathfrak m)\).
\end{proof}

\begin{lemma}[Localised form of the height lower bound]
  \label{pa-lem:standard_smooth_dim_lower_localized}
  \lean{Algebra.IsStandardSmoothOfRelativeDimension.le_ringKrullDim_of_isLocalization_atPrime}
  \leanok
  \dcref{custom}

  \uses{pa-thm:standard_smooth_height_lower}
  In the setting of \ref{pa-thm:standard_smooth_height_lower}, for any localisation
  \(S_{\mathfrak m}\) of \(S\) at a maximal ideal \(\mathfrak m\),
  \(\dim S_{\mathfrak m} \geq n\).
\end{lemma}
\begin{proof}
  \leanok
  The Krull dimension of the localisation at a prime is the height of that prime
  (primes of \(S_{\mathfrak m}\) are the primes of \(S\) contained in
  \(\mathfrak m\)); apply \ref{pa-thm:standard_smooth_height_lower}.
\end{proof}

\begin{lemma}[Regularity from the cotangent upper bound]
  \label{pa-lem:regular_of_cotangent_le_dim}
  \lean{IsRegularLocalRing.of_finrank_cotangentSpace_le_ringKrullDim}
  \leanok
  \dcref{custom}

  Let \((A, \mathfrak m)\) be a Noetherian local ring with
  \(\dim_\kappa \mathfrak m/\mathfrak m^2 \leq \dim A\). Then \(A\) is a regular local
  ring.
\end{lemma}
\begin{proof}
  \leanok
  By Nakayama, \(\dim_\kappa \mathfrak m/\mathfrak m^2\) is the minimal number
  \(\mu(\mathfrak m)\) of generators of \(\mathfrak m\), and Krull's height theorem
  gives the reverse inequality \(\dim A \leq \mu(\mathfrak m)\) in any Noetherian local
  ring. So the hypothesis forces equality, which is the definition of regularity.
\end{proof}

\section{The trdeg--height inequality and smooth-prime regularity}
\label{pa-sec:smooth_prime_regularity}

This section closes the smooth-implies-regular pipeline at \emph{arbitrary} primes of a
standard-smooth algebra over a perfect field, with no recourse to Serre's homological
characterisation of regularity (and hence no need for the theorem that localisations of
regular local rings are regular): the cotangent dimension at a prime is computed
directly from the conormal sequence, and matched against a transcendence-degree/height
inequality proved by induction on the number of variables.

\begin{lemma}[One-variable step of the trdeg--height inequality]
  \label{pa-lem:polynomial_step_height_trdeg}
  \lean{Polynomial.step_height_trdeg_of_isPrime}
  \leanok
  \source{stacks-project}

  \dcref{custom}
  Let \(k\) be a field, \(R\) a Noetherian \(k\)-algebra, \(P \subseteq R[X]\) a prime
  ideal and \(\mathfrak p = P \cap R\). Then one of the following holds:
  \begin{enumerate}
    \item \(\mathrm{ht}(\mathfrak p) + 1 \leq \mathrm{ht}(P)\) and
      \(\mathrm{trdeg}_k(R/\mathfrak p) \leq \mathrm{trdeg}_k(R[X]/P)\); or
    \item \(\mathrm{ht}(P) = \mathrm{ht}(\mathfrak p)\) and
      \(\mathrm{trdeg}_k(R/\mathfrak p) + 1 \leq \mathrm{trdeg}_k(R[X]/P)\).
  \end{enumerate}
\end{lemma}
\begin{proof}
  \leanok
  Write \(A = R/\mathfrak p\), a domain. The extension \(R \to R[X]\) is free, hence
  flat, hence has the going-down property, so the dimension formula for a prime lying
  over \(\mathfrak p\) applies (Stacks 00ON):
  \(\mathrm{ht}(P) = \mathrm{ht}(\mathfrak p) + \mathrm{ht}(\bar P)\), where
  \(\bar P\) is the image of \(P\) in the fibre ring
  \(R[X]/\mathfrak p R[X] \cong A[X]\) --- a prime \(Q \subseteq A[X]\) contracting to
  \((0)\) in \(A\) (because \(P \cap R = \mathfrak p\)).

  \emph{Case \(Q = (0)\).} Then \(\mathrm{ht}(P) = \mathrm{ht}(\mathfrak p)\), and
  \(R[X]/P \cong A[X]/Q = A[X]\) as \(k\)-algebras. Adjoining a variable to the domain
  \(A\) raises the transcendence degree by one
  (\(\mathrm{trdeg}_k A[X] \geq \mathrm{trdeg}_k A + \mathrm{trdeg}_A A[X]
  = \mathrm{trdeg}_k A + 1\)), giving case (2).

  \emph{Case \(Q \neq (0)\).} Then \((0) \subsetneq Q\) is a chain of primes in
  \(A[X]\), so \(\mathrm{ht}(Q) \geq 1\) and
  \(\mathrm{ht}(P) \geq \mathrm{ht}(\mathfrak p) + 1\). Moreover \(Q\) contracts to
  \((0)\) in \(A\), so the composite \(A \to A[X] \to A[X]/Q \cong R[X]/P\) is
  injective, whence
  \(\mathrm{trdeg}_k(A) \leq \mathrm{trdeg}_k(R[X]/P)\), giving case (1).
\end{proof}

\begin{theorem}[The trdeg--height inequality in polynomial rings]
  \label{pa-thm:trdeg_height_polynomial}
  \lean{MvPolynomial.exists_le_trdeg_and_natCard_le_height_add}
  \leanok
  \dcref{custom}

  \uses{pa-lem:polynomial_step_height_trdeg}
  Let \(k\) be a field, \(\iota\) a finite index type, and \(P\) a prime ideal of
  \(k[x_i : i \in \iota]\). Then there is a natural number \(d\) with
  \[
    d \leq \mathrm{trdeg}_k\bigl(k[x_i : i \in \iota]/P\bigr)
    \quad\text{and}\quad
    \#\iota \leq \mathrm{ht}(P) + d.
  \]
\end{theorem}
\begin{proof}
  \leanok
  Induct on the finite type \(\iota\) (along equivalences, the empty type, and adding
  one variable); heights and transcendence degrees are invariant under \(k\)-algebra
  isomorphisms, which handles equivalences. For \(\iota\) empty take \(d = 0\). For
  \(\iota = \{*\} \sqcup \alpha\), identify the ring with \(R[X]\),
  \(R = k[x_i : i \in \alpha]\), and let \(\mathfrak p = P \cap R\); by the inductive
  hypothesis pick \(d\) with \(d \leq \mathrm{trdeg}_k(R/\mathfrak p)\) and
  \(\#\alpha \leq \mathrm{ht}(\mathfrak p) + d\). Apply
  \ref{pa-lem:polynomial_step_height_trdeg}. In case (1) (height jump) keep the witness
  \(d\):
  \(d \leq \mathrm{trdeg}_k(R/\mathfrak p) \leq \mathrm{trdeg}_k(R[X]/P)\) and
  \(\#\iota = \#\alpha + 1 \leq \mathrm{ht}(\mathfrak p) + 1 + d \leq
  \mathrm{ht}(P) + d\). In case (2) (transcendence jump) take \(d + 1\):
  \(d + 1 \leq \mathrm{trdeg}_k(R/\mathfrak p) + 1 \leq \mathrm{trdeg}_k(R[X]/P)\) and
  \(\#\iota = \#\alpha + 1 \leq \mathrm{ht}(\mathfrak p) + d + 1 = \mathrm{ht}(P) +
  (d + 1)\).
\end{proof}

\begin{theorem}[The trdeg--height inequality for standard-smooth algebras]
  \label{pa-thm:trdeg_height_standard_smooth}
  \lean{Algebra.IsStandardSmoothOfRelativeDimension.exists_le_trdeg_and_natCast_le_height_add}
  \leanok
  \dcref{custom}

  \uses{pa-thm:trdeg_height_polynomial}
  Let \(k\) be a field and \(S\) a standard-smooth \(k\)-algebra of relative dimension
  \(n\). For every prime ideal \(\mathfrak q \subseteq S\) there is a natural number
  \(d\) with
  \[
    d \leq \mathrm{trdeg}_k(S/\mathfrak q)
    \quad\text{and}\quad
    n \leq \mathrm{ht}(\mathfrak q) + d.
  \]
\end{theorem}
\begin{proof}
  \leanok
  As in \ref{pa-thm:standard_smooth_height_lower}: choose a submersive presentation
  \(P = k[x_i : i \in \iota] \twoheadrightarrow S\) with kernel generated by
  \(\#\sigma\) relations and \(n = \#\iota - \#\sigma\). Let \(M \subseteq P\) be the
  pullback of \(\mathfrak q\), a prime containing the kernel. By
  \ref{pa-thm:trdeg_height_polynomial} pick \(d\) with \(d \leq \mathrm{trdeg}_k(P/M)\) and
  \(\#\iota \leq \mathrm{ht}(M) + d\). The induced map \(P/M \to S/\mathfrak q\) is an
  injective \(k\)-algebra map (indeed an isomorphism by the first isomorphism theorem),
  so \(d \leq \mathrm{trdeg}_k(S/\mathfrak q)\). Krull's height theorem bounds the
  height defect by the number of relations,
  \(\mathrm{ht}(M) \leq \mathrm{ht}(\mathfrak q) + \#\sigma\), so
  \(\#\iota \leq \mathrm{ht}(\mathfrak q) + \#\sigma + d\) and
  \(n = \#\iota - \#\sigma \leq \mathrm{ht}(\mathfrak q) + d\).
\end{proof}

\begin{theorem}[K\"ahler rank equals transcendence degree over a perfect field]
  \label{pa-thm:kaehler_rank_trdeg}
  \lean{Algebra.rank_kaehlerDifferential_eq_trdeg}
  \leanok
  \dcref{custom}

  Let \(k\) be a perfect field and \(K\) a field extension of \(k\) that is essentially
  of finite type. Then
  \[
    \dim_K \Omega_{K/k} = \mathrm{trdeg}_k(K).
  \]
\end{theorem}
\begin{proof}
  \leanok
  Since \(k\) is perfect and \(K\) essentially of finite type over \(k\), \(K\) admits a
  \emph{separating} transcendence basis \(s \subseteq K\): a transcendence basis such
  that \(K\) is separable algebraic over the intermediate field \(F = k(s)\). A
  separable algebraic extension is formally \'etale, so base change along \(F \to K\)
  identifies \(\Omega_{K/k} \cong K \otimes_F \Omega_{F/k}\); in particular
  \(\dim_K \Omega_{K/k} = \dim_F \Omega_{F/k}\). The field \(F\) is \(k\)-isomorphic to
  the fraction field of the polynomial ring \(k[x_i : i \in s]\), and K\"ahler
  differentials localise: \(\Omega_{F/k}\) is the localisation of the free module
  \(\Omega_{k[x_i]/k}\) with basis \(\{dx_i : i \in s\}\), hence is \(F\)-free on
  \(\{dx_i\}\). Therefore
  \(\dim_F \Omega_{F/k} = \#s = \mathrm{trdeg}_k(K)\), the last equality because \(s\)
  is a transcendence basis.
\end{proof}

\begin{theorem}[The conormal dimension identity]
  \label{pa-thm:conormal_dim_identity}
  \lean{finrank_cotangentSpace_add_finrank_kaehler_residueField}
  \leanok
  \source{stacks-project}

  \dcref{custom}
  Let \(R\) be a commutative ring and \((S, \mathfrak m)\) a local \(R\)-algebra that is
  formally smooth over \(R\), whose residue field \(\kappa = S/\mathfrak m\) is also
  formally smooth over \(R\), and whose K\"ahler differentials \(\Omega_{S/R}\) are
  \(S\)-free of rank \(n \in \mathbb{N}\). Then
  \[
    \dim_\kappa \mathfrak m/\mathfrak m^2 + \dim_\kappa \Omega_{\kappa/R} = n.
  \]
\end{theorem}
\begin{proof}
  \leanok
  Consider the conormal (second fundamental) exact sequence of the surjection
  \(S \twoheadrightarrow \kappa\) with kernel \(\mathfrak m\) (Stacks 00RU):
  \[
    \mathfrak m/\mathfrak m^2 \xrightarrow{\ c\ }
      \kappa \otimes_S \Omega_{S/R} \longrightarrow \Omega_{\kappa/R} \longrightarrow 0,
    \qquad c(\bar f) = df \otimes 1.
  \]
  Because \(\kappa\) is formally smooth over \(R\), the surjection
  \(S \twoheadrightarrow \kappa\) admits enough infinitesimal liftings to \emph{split}
  the conormal map: \(c\) is (split) injective. All three terms are
  \(\kappa\)-vector spaces and the middle has dimension
  \(\dim_\kappa (\kappa \otimes_S \Omega_{S/R}) = n\) (base change of a free rank-\(n\)
  module). The map onto \(\Omega_{\kappa/R}\) is surjective with kernel exactly the
  image of \(c\), which has dimension \(\dim_\kappa \mathfrak m/\mathfrak m^2\) by
  injectivity; rank--nullity gives the identity. (Finite-dimensionality of all terms
  follows from that of the middle term.)
\end{proof}

\begin{theorem}[Smooth-prime regularity over a perfect field]
  \label{pa-thm:smooth_prime_regularity}
  \lean{isRegularLocalRing_of_isLocalization_atPrime_of_isStandardSmooth_of_perfectField}
  \leanok
  \source{stacks-project}

  \uses{pa-thm:trdeg_height_standard_smooth, pa-thm:kaehler_rank_trdeg,
    pa-thm:conormal_dim_identity, pa-lem:regular_of_cotangent_le_dim}
  \dcref{custom}
  Let \(k\) be a perfect field (for instance an algebraically closed field) and \(S\) a
  nonzero standard-smooth \(k\)-algebra. For every prime ideal
  \(\mathfrak q \subseteq S\), every localisation \(S_{\mathfrak q}\) of \(S\) at
  \(\mathfrak q\) is a regular local ring.
\end{theorem}
\begin{proof}
  \leanok
  \(S\) is of finite type over \(k\), hence Noetherian, and so is the localisation
  \(A := S_{\mathfrak q}\), a local ring with maximal ideal \(\mathfrak m\) and residue
  field \(\kappa\). Fix \(n\) with \(S\) standard smooth of relative dimension \(n\)
  over \(k\).

  \emph{The K\"ahler package at \(A\).} \(\Omega_{S/k}\) is \(S\)-free of rank \(n\)
  (standard smoothness), and K\"ahler differentials localise, so \(\Omega_{A/k}\) is
  \(A\)-free of rank \(n\). Localisations of (formally) smooth algebras are formally
  smooth, so \(A\) is formally smooth over \(k\); and \(\kappa\), a field extension of
  the perfect field \(k\) essentially of finite type, is formally smooth over \(k\) as
  well. The conormal identity \ref{pa-thm:conormal_dim_identity} thus gives
  \[
    \dim_\kappa \mathfrak m/\mathfrak m^2 \;=\; n - \dim_\kappa \Omega_{\kappa/k}.
  \]

  \emph{The dimension count.} By \ref{pa-thm:kaehler_rank_trdeg},
  \(\dim_\kappa \Omega_{\kappa/k} = \mathrm{trdeg}_k(\kappa)\). By
  \ref{pa-thm:trdeg_height_standard_smooth} there is \(d\) with
  \(d \leq \mathrm{trdeg}_k(S/\mathfrak q)\) and
  \(n \leq \mathrm{ht}(\mathfrak q) + d\); and \(S/\mathfrak q\) embeds into \(\kappa\)
  as \(k\)-algebras (localisation at \(\mathfrak q\) followed by reduction is injective
  on \(S/\mathfrak q\)), so
  \(d \leq \mathrm{trdeg}_k(\kappa) = \dim_\kappa \Omega_{\kappa/k}\). Combining,
  \[
    \dim_\kappa \mathfrak m/\mathfrak m^2
      = n - \dim_\kappa \Omega_{\kappa/k}
      \leq \bigl(\mathrm{ht}(\mathfrak q) + d\bigr) - \dim_\kappa \Omega_{\kappa/k}
      \leq \mathrm{ht}(\mathfrak q) = \dim A.
  \]
  By \ref{pa-lem:regular_of_cotangent_le_dim}, \(A\) is a regular local ring.
\end{proof}
